\documentclass[11pt]{article}
\usepackage[T1]{fontenc}

\usepackage[margin=1.02in]{geometry}
\usepackage{amsmath,amssymb,amsthm,mathtools,mathrsfs}
\usepackage{bm}
\usepackage{aliascnt}
\usepackage{enumitem}
\usepackage{booktabs,array,tabularx}
\usepackage{microtype}
\usepackage[colorlinks=true,linkcolor=blue,citecolor=blue,urlcolor=blue]{hyperref}
\usepackage[nameinlink,capitalise]{cleveref}

% Separate navigation for the exposition and both parts of the supplement.
\newif\ifmathsupplement
\mathsupplementfalse
\makeatletter
\let\RG@addcontentsline\addcontentsline
\renewcommand{\addcontentsline}[3]{%
  \def\RG@tocname{#1}\def\RG@toc{toc}%
  \ifx\RG@tocname\RG@toc
    \RG@addcontentsline{toc}{#2}{#3}% global outline; local lists below
    \ifmathsupplement
      \RG@addcontentsline{supptoc}{#2}{#3}%
    \else
      \RG@addcontentsline{maintoc}{#2}{#3}%
    \fi
  \else
    \RG@addcontentsline{#1}{#2}{#3}%
  \fi}
\newcommand{\maintableofcontents}{\section*{Contents}\@starttoc{maintoc}}
\newcommand{\supplementtableofcontents}{%
  \section*{Contents of the Mathematical Supplement}%
  \begingroup
    \renewcommand*\l@section[2]{%
      \addvspace{0.65em}%
      \begingroup\bfseries
      \@dottedtocline{1}{0em}{2.8em}{##1}{##2}%
      \endgroup}%
    \renewcommand*\l@subsection{\@dottedtocline{2}{2.8em}{3.6em}}%
    \@starttoc{supptoc}%
  \endgroup}
\makeatother

\usepackage{tikz-cd}
\setlength{\emergencystretch}{4em}
\raggedbottom
\hbadness=10000
\hypersetup{pdftitle={From Predictive Dynamics to Spectral Geometry},pdfauthor={Aurelien Pelissier}}

% -----------------------------------------------------------------------------
% theorem environments
% -----------------------------------------------------------------------------
% theorem environments
\newtheorem{theorem}{Theorem}[section]
\newaliascnt{proposition}{theorem}
\newtheorem{proposition}[proposition]{Proposition}
\aliascntresetthe{proposition}
\crefname{proposition}{proposition}{propositions}
\Crefname{proposition}{Proposition}{Propositions}
\newaliascnt{lemma}{theorem}
\newtheorem{lemma}[lemma]{Lemma}
\aliascntresetthe{lemma}
\crefname{lemma}{lemma}{lemmas}
\Crefname{lemma}{Lemma}{Lemmas}
\newaliascnt{corollary}{theorem}
\newtheorem{corollary}[corollary]{Corollary}
\aliascntresetthe{corollary}
\crefname{corollary}{corollary}{corollaries}
\Crefname{corollary}{Corollary}{Corollarys}
\theoremstyle{definition}
\newaliascnt{definition}{theorem}
\newtheorem{definition}[definition]{Definition}
\aliascntresetthe{definition}
\crefname{definition}{definition}{definitions}
\Crefname{definition}{Definition}{Definitions}
\newaliascnt{construction}{theorem}
\newtheorem{construction}[construction]{Construction}
\aliascntresetthe{construction}
\crefname{construction}{construction}{constructions}
\Crefname{construction}{Construction}{Constructions}
\theoremstyle{remark}
\newaliascnt{remark}{theorem}
\newtheorem{remark}[remark]{Remark}
\aliascntresetthe{remark}
\crefname{remark}{remark}{remarks}
\Crefname{remark}{Remark}{Remarks}
\newaliascnt{example}{theorem}
\newtheorem{example}[example]{Example}
\aliascntresetthe{example}
\crefname{example}{example}{examples}
\Crefname{example}{Example}{Examples}


\theoremstyle{definition}
\newtheorem{assumption}[theorem]{Assumption}
\theoremstyle{remark}
\newtheorem{scope}[theorem]{Scope statement}
\newtheorem{interpretation}[theorem]{Interpretation}
\newtheorem{countertheorem}[theorem]{Countertheorem}

% -----------------------------------------------------------------------------
% notation
% -----------------------------------------------------------------------------
\newcommand{\N}{\mathbb N}
\newcommand{\Z}{\mathbb Z}
\newcommand{\R}{\mathbb R}
\newcommand{\C}{\mathbb C}
\newcommand{\Torus}{\mathbb T}
\newcommand{\one}{\mathbf 1}
\newcommand{\id}{\mathrm{id}}
\newcommand{\dd}{\mathrm d}
\newcommand{\Tr}{\operatorname{Tr}}
\newcommand{\tr}{\operatorname{tr}}
\newcommand{\rank}{\operatorname{rank}}
\newcommand{\Ker}{\operatorname{Ker}}
\newcommand{\Ran}{\operatorname{Ran}}
\newcommand{\Span}{\operatorname{span}}
\newcommand{\Hom}{\operatorname{Hom}}
\newcommand{\End}{\operatorname{End}}
\newcommand{\supp}{\operatorname{supp}}
\newcommand{\spec}{\operatorname{spec}}
\newcommand{\diag}{\operatorname{diag}}
\newcommand{\norm}[1]{\lVert#1\rVert}
\newcommand{\abs}[1]{\lvert#1\rvert}
\newcommand{\ip}[2]{\left\langle #1,#2\right\rangle}
\newcommand{\A}{\mathcal A}
\newcommand{\Hh}{\mathcal H}
\newcommand{\Om}{\Omega}
\newcommand{\GT}{\boldsymbol{\mathbb T}_{\mathrm{Grand}}}
\newcommand{\GTX}{\boldsymbol{\mathbb T}_{\mathrm{Grand},X}}
\newcommand{\Lip}{\operatorname{Lip}}
\newcommand{\Cl}{\mathrm{Cl}}
\newcommand{\HS}{\mathrm{HS}}
\newcommand{\vol}{\operatorname{vol}}
\newcommand{\supplement}{Mathematical Supplement}
\newcommand{\EssIm}{\operatorname{EssIm}}
\newcommand{\Curr}{\operatorname{Curr}}
\newcommand{\RReal}{\operatorname{RReal}}
\newcommand{\Corr}{\operatorname{Corr}}
\newcommand{\Stab}{\operatorname{Stab}}
\newcommand{\McM}{\operatorname{McM}}
\newcommand{\RG}{\operatorname{RG}}
\newcolumntype{L}[1]{>{\raggedright\arraybackslash}p{#1}}


\title{\textbf{From Predictive Dynamics to Spectral Geometry}}

\author{%
  Aurelien Pelissier$^{\dagger}$\\
  {\small $^{\dagger}$Correspondence: \href{mailto:aurelien.pelissier.38@gmail.com}{aurelien.pelissier.38@gmail.com}}%
  \vspace{-0.5cm}
}
\date{}



\begin{document}
\maketitle

\begin{abstract}
In Alain Connes' noncommutative geometry, geometry is encoded by an algebra, a Hilbert space, and a Dirac operator taken as fundamental spectral data. We show that these structures can instead emerge from a more primitive predictive dynamics. A compatible first-order variation on the induced history algebra canonically generates a real even Hodge--Dirac spectral triple, while the algebra alone determines neither its metric scale nor the full Dirac operator.
%
The resulting spectral geometry is an information-reduced description: dynamically inequivalent predictive systems can share the same effective spectral data. Nevertheless, a finite commutator derivation determines a canonical centered Dirac implementer, with the remaining ambiguity recovered by the sharp finite number of independent real linear commutant coordinates, and on periodic three-dimensional lattice approximants the emergent spectral distance converges to a Riemannian metric reconstructed independently from the dynamics. At fixed vertex masses and symmetric conductances, the strictly positive stationary-current fibre is globally parametrized by the cohomology class of its edge affinities.
%
Thus Connes-type spectral geometry can arise as an effective operator-geometric description of richer predictive dynamics, providing a concrete setting in which spectral structure is reconstructed rather than assumed at the microscopic entrance.
\end{abstract}

\medskip
\noindent\textbf{Keywords.}
Noncommutative geometry; spectral triples; Renewal Geometry; predictive dynamics;
Hodge--Dirac operators; Connes distance; stationary currents; spectral reconstruction.

\setcounter{tocdepth}{2}
\clearpage
\maintableofcontents
\clearpage

\section{Introduction}
\label{sec:introduction}
% =============================================================================

Many spaces that arise naturally in mathematics and physics do not admit a
satisfactory description as ordinary geometric spaces. Quotients may be singular,
points may fail to capture the relevant observables, and the structures of interest
may be encoded more faithfully by operators than by a classical manifold. In such
situations the usual geometric language---built from points, coordinates, and
functions on a well-behaved space---can cease to be adequate, even though an
algebra of observables remains meaningful.

Noncommutative geometry was developed to address this problem. Connes' programme
replaces the requirement of an underlying point-set space by an operator-algebraic
description in which an algebra plays the role of generalized coordinates. In the
commutative case this remains compatible with ordinary geometry, since a suitable
commutative algebra determines the underlying space through Gelfand duality.
Allowing the algebra to become noncommutative extends the same principle to
settings for which no classical space of points need exist.

The spectral formulation equips this generalized coordinate algebra with geometry.
A compact spin manifold can be encoded by an algebra $\mathcal A$ represented on
a Hilbert space $\mathcal H$, together with a Dirac operator $D$
\cite{Connes1994}. The representation supplies the operator-theoretic carrier,
while $D$ encodes differential and metric information.  Under the reconstruction
axioms for commutative spectral triples, this operator data recovers compact spin
geometry \cite{RennieVarilly2008,connesreconstruction}; in the genuinely
noncommutative case the same operator-geometric language continues to apply even
when no ordinary underlying space exists.

Spectral triples consequently provide a common language for noncommutative spaces,
differential calculi, index theory, and metric geometry. The Connes distance
formula turns the commutator with $D$ into a metric notion, making the Dirac
operator simultaneously a differential and geometric object. Related work on
compact quantum metric spaces formulates metric information at the level of an
algebra equipped with a Lipschitz seminorm \cite{Rieffel1999}.  Rieffel's quantum
Gromov--Hausdorff distance and the later propinquity and spectral propinquity
provide frameworks for convergence of quantum metric spaces and metric spectral
triples \cite{Rieffel2004,Latremoliere2022,FarsiLatremolierePacker2024}, while real
structures and gradings encode analogues of spin and orientation data.

A different structural question is whether the objects entering this spectral
description must themselves be regarded as primitive. This is logically prior to
spectral reconstruction: rather than asking which geometry is encoded by a given
spectral triple, one asks whether the algebra, representation, differential
structure, and Dirac data can arise from an underlying dynamics in which the persistent algebraic and Hilbert-space carrier is not initially prescribed,
while the compatible differential is specified separately.

\paragraph{From predictive dynamics to spectral geometry.}

Renewal Geometry provides one setting in which this question can be posed
operationally \cite{PelissierRenewalSpacetime2026}. Its starting point is not a
manifold, graph, causal order, algebra, Hilbert space, or spectral triple, but a
finite operational process specified by composable intervention histories and
their conditional future statistics. Histories that cannot be distinguished by
any admissible future continuation are identified. The resulting
future-equivalence quotient reconstructs the minimum persistent predictive
carrier of the process rather than requiring such a carrier to be supplied
independently.

Complete operational statistics then reconstruct a canonical positive history
object, the \emph{Renewal Grand Tensor}.  We denote the Grand Tensor associated
with a finite predictive system $X$ by $\GTX$.  Its role is pregeometric: it
records which distinctions survive prediction, how represented histories compose,
and how they overlap, before positions, distances, metrics, or continuum geometry
have been assigned.  The corresponding history algebra and positive moment
structure therefore arise from the operational process rather than being inserted
as geometric input.

For the present paper, the relevant interface is the algebraic content carried by
$\GTX$.  The Grand Tensor canonically supplies the future-minimal predictive
carrier, the represented history algebra $\A_X^{\mathrm{hist}}$, and its positive
moment structure.  On the finite algebra there is in particular an intrinsic
normalized regular trace, which will provide one canonical choice of $\tau_X$.
To pass from this reconstructed algebraic object to spectral geometry, one then
specifies a compatible first-order variation $\delta_X$ of the history algebra.
Thus the interface used here is
\begin{equation}
\GTX
\longmapsto
(\A_X^{\mathrm{hist}},\tau_X)
\overset{\delta_X}{\longrightarrow}
(\A_X,\Hh_X,D_X,J_X,\gamma_X).
\label{eq:RGT-spectral-interface}
\end{equation}
The first arrow is reconstructed directly from the operational statistics and is
summarized in Section~\ref{sec:finite-main} and developed self-containedly
in Section~\ref{app:operational-recognition}.  The second is the
Hodge--Dirac spectralization studied in this paper.  The Grand Tensor is not
identified with the displayed algebraic or spectral data: it retains the
predictive carrier and complete positive history structure, including information
that need not survive in the spectral image.  Nor do we assume that the Grand
Tensor or its algebra alone fixes $\delta_X$; the first-order variation is exactly
the additional datum whose metric scale and Dirac ambiguity are analyzed below.

The same $\GTX$ also supports the separate Lorentzian spacetime construction
developed within the Renewal Geometry framework~\cite{PelissierRenewalSpacetime2026}.
The present work asks instead whether
Connes-type spectral geometry can arise from a compatible first-order variation
of its reconstructed history algebra.  The additional spacetime-specific assumptions of the Lorentzian construction
are not used below; once a finite-dimensional unital $C^*$-algebra $\A$, a
faithful tracial state $\tau$, and a compatible real Hilbert-bimodule derivation
$\delta$ are supplied, the finite spectralization results are ordinary
operator-theoretic statements independent of the operational route by which
$\A$ was obtained.

The resulting problem has both a forward and an inverse aspect. On the forward
side one asks whether a first-order predictive differential canonically produces
the structures of spectral geometry, and whether the resulting metric has an
independent dynamical meaning. On the inverse side one asks which spectral
geometries can be reached, what additional data are required to reconstruct a
full Dirac operator, and whether different microscopic dynamics can have the same
spectral image.

\paragraph{What lies before a spectral triple?}

Finite spectral geometry already has a substantial reconstruction and
classification theory.  Krajewski classified finite spectral triples in terms of
matrix data and diagrams \cite{Krajewski1998}, while approximately finite
$C^*$-algebras, weighted trees, ultrametric Cantor sets, fractals, stationary
Bratteli diagrams, and aperiodic systems admit natural spectral-triple
constructions from prescribed filtrations or hierarchical data
\cite{ChristensenIvan2006,PearsonBellissard2009,
ChristensenIvanSchrohe2012,KellendonkSavinien2012}.  Related constructions attach
spectral data to given $C^*$-dynamical systems, including crossed products,
symbolic dynamics, and irreversible dynamics
\cite{Paterson2014,JulienPutnam2016,AielloGuidoIsola2022}.  In these settings the
algebraic or geometric carrier is already present before spectralization.

A closer operator-theoretic precedent comes from quantum Riemannian geometry and
noncommutative potential theory.  Majid constructs geometrically realized Dirac
operators from a prescribed noncommutative coordinate algebra, differential
calculus, quantum metric, and compatible bimodule connection \cite{Majid2025}.
Cipriani and Sauvageot realize tracially symmetric Dirichlet forms through
closable derivations into Hilbert bimodules \cite{CiprianiSauvageot2003}; this
calculus underlies Dirac constructions on fractals and spectral triples associated
with Dirichlet, resistance, and sub-Markovian structures
\cite{HinzTeplyaev2013,HinzKelleherTeplyaev2015,Cipriani2016,Arhancet2024}.
These works are particularly close to the Hodge--Dirac layer used below, but they
begin with the algebra and differential or Markov structure already specified.

The operational entrance has complementary precedents.  Quantum combs and
process-tensor formalisms encode complete multitime intervention statistics by
positive causal process objects and describe memory-bearing dynamics without a
Markov approximation
\cite{ChiribellaDArianoPerinotti2009,PollockProcessTensor2018,PollockEtAl2018}.
Classically, observable operator models, predictive-state representations, and
computational mechanics define state through future predictions rather than hidden
microscopic labels
\cite{Jaeger2000,LittmanSuttonSingh2001,SinghJamesRudary2004,
CrutchfieldYoung1989,shalizicrutchfield,TraversCrutchfield2025,
MarzenCrutchfield2017}.  The predictive construction used here combines this
future-equivalence principle with multitime operational composition and positive
process structure: the persistent carrier and represented history algebra are
reconstructed only after future-invisible distinctions have been removed, and a
compatible first-order differential is then spectralized on that reconstructed
algebra.

The resulting question is whether the algebra, Hilbert-space carrier,
differential structure, and Dirac data of noncommutative geometry can arise from
an underlying predictive dynamics in which the persistent carrier is not initially prescribed and the differential
input is made explicit.
A nontrivial answer requires more than reproducing a desired spectral triple by
definition: one must identify what the first-order differential actually fixes,
which spectral geometries are obtainable under different reconstruction
categories, whether the resulting metric has an independent dynamical meaning,
and what information is lost under spectralization.

The last point is essential.  If distinct predictive systems have the same
spectral image, then a spectral triple is not, in general, a complete invariant of
the dynamics from which it arose.  It may instead describe an equivalence class
of microscopically distinct systems that agree on their effective metric and
differential structure.  In this sense spectral geometry can be
\emph{information-reduced}: complete for the operator-geometric information it
encodes while insufficient, in general, to reconstruct the richer underlying
dynamics.

\paragraph{Our contribution.}

We first construct a canonical spectralization of finite differential data. A
star-compatible derivation $\delta$ into a real Hilbert $\A$-bimodule determines
a minimal one-form carrier and a real even Hodge--Dirac spectral triple with exact
order-zero and first-order relations. Two rigidity results identify what this
construction does not determine. The algebra and faithful trace do not fix an
absolute Dirac scale, and a fixed commutator derivation determines the full
compatible Dirac operator only modulo a finite-dimensional self-adjoint commutant.
The orthogonally centered representative is the unique
minimum-Hilbert--Schmidt-norm implementer, and the dimension of this commutant is
the sharp number of independent real linear coordinates required to distinguish the
remaining Dirac ambiguity.

This distinction separates an \emph{austere Hodge construction} from an enlarged
\emph{calibrated marked reconstruction}. The bare Hodge operator always has the
unit zero-mode and therefore cannot realize arbitrary finite Dirac operators.
When the exact $q_{\mathsf S}=\dim_{\R}\mathfrak c_{\mathsf S}$ scalar
commutant anchors are retained, every finite spectral packet is reconstructible,
and fewer independent real linear coordinates cannot distinguish the full affine Dirac fibre.  Every
separable compact-resolvent packet arises as a strong cofinal limit of finite
marked screens. A stricter requirement,
in which multiplication must be transported asymptotically in norm, is genuinely
selective and is characterized exactly by spectral quasidiagonality.

We next test whether the construction has independent geometric content. For
stationary graph dynamics, the symmetric jump data determine a Hodge--Dirac
operator whose square on zero-forms is $-2$ times the symmetric Markov generator and whose commutator
seminorm is its local quadratic form. On periodic $A_3$ approximants, the
corresponding Connes distance converges uniformly to a Riemannian metric
reconstructed independently from the jump second moments. In the spin extension, the local positive-stencil and covariant Wilson construction gives
the required chartwise estimates; on flat periodic spin tori these estimates close
globally and give unconditional norm-resolvent convergence, while on general closed
Riemannian spin manifolds the same conclusion follows under the stated compatible
stable-spin discretization hypotheses. The independent scalar metric reconstruction
is unconditional.

Finally, we characterize several mechanisms by which spectralization loses
information. Directed stationary current can vary at fixed metric spectral
geometry, hidden dynamical memory can vary at fixed static spectral data, and
locally identical spectral packets can differ by global marked monodromy. In their stated classes, these
inverse fibres have intrinsic coordinates and compatible coarse-graining maps.
Thus spectralization is not merely noninjective: the information it forgets has
its own mathematical structure and transformation laws.

The two principal graph results are the uniform $A_3$ metric limit and the
identification of the strictly positive stationary-current fibre with
$H^1(G;\R)$ through its affinity class.  They form the mathematical core,
together with the finite differential and operational reconstruction theory.
The spin, rational-memory, and completed-fibre constructions are further
extensions, each with its own additional hypotheses; they are not needed for
those two graph theorems.

The resulting interpretation is that Connes-type spectral geometry can arise as
an effective geometric description of a richer predictive dynamics rather than
necessarily as its most primitive level.  The independent Lorentzian construction
\cite{PelissierRenewalSpacetime2026} shows that Lorentzian spacetime and spectral
geometry can both arise downstream of the same predictive structure, without
either being required as the primitive entrance to the other.  The two
constructions therefore lead to the schematic picture
\begin{equation}
\boxed{
\begin{array}{c}
\text{operational histories}\\[1mm]
\downarrow\\[1mm]
\text{predictive history reconstruction}\\[1mm]
\downarrow\\[1mm]
\text{Renewal Grand Tensor}\\[1mm]
\swarrow\hspace{2.5cm}\searrow\\[-1mm]
\text{Lorentzian spacetime}
\hspace{1.4cm}
\text{Connes-type spectral geometry}.
\end{array}}
\label{eq:intro-renewal-geometric-branches}
\end{equation}

The Lorentzian branch is developed in
\cite{PelissierRenewalSpacetime2026}.  The present paper develops the spectral
branch: its canonical finite construction, its category-dependent image, its
independent metric test, and its inverse fibres.  The central question is therefore
not only how spectral geometry can arise, but what a richer underlying dynamics
retains that must be forgotten when it is compressed to a spectral description.

















% =============================================================================

\paragraph{Reading the paper.}
The main text gives the predictive construction, the principal theorem statements,
and their interpretation.  Part~I of the Mathematical Supplement contains the
self-contained core proofs.  Part~II retains the additional spin, dilation,
indefinite-memory, and joint-realization results.  The shared notation and
cross-references make the extensions available without making them prerequisites
for the graph theorems.

\section{From predictive dynamics to spectral geometry}
\label{sec:finite-main}

The spectralization theorem is algebraically self-contained, but its motivating
entrance is a finite predictive process.  We retain the operational description
and the Renewal Grand Tensor interface here; the full reconstruction,
realization, and recognition arguments are given in
Section~\ref{app:operational-recognition}.  This separates the meaning of the
construction from its technical proofs without changing the input data.

\subsection{The Renewal Grand Tensor and the first-order differential interface}
\label{subsec:grand-tensor-main}

Fix a finite cutoff $X$, a finite typed event grammar, and the reachable
histories.  The conditional event law determines the future cylinder probabilities
\begin{equation}
 p_X(w\mid h)=\prod_{j=1}^k p_X(a_j\mid ha_1\cdots a_{j-1}),
 \qquad w=a_1\cdots a_k.
 \label{eq:overview-chain-law}
\end{equation}
Histories ending at the same cut type are identified when every declared future
continuation has the same conditional probability.  Thus
\begin{equation}
 h\sim_Xh'\quad\Longleftrightarrow\quad
 p_X(w\mid h)=p_X(w\mid h')\ \text{for every admissible future }w,
 \qquad Z_{X,x}^{\min}=\mathsf H_{X,x}/\!\sim_X.
 \label{eq:overview-predictive-state}
\end{equation}
The branch probabilities are $\kappa_{X,a}([h])=p_X(a\mid h)$, and the update
$T_{X,a}([h])=[ha]$ is defined only on branches with positive probability.
This partial-transition convention matters: no conditional successor is inferred
from a zero-probability extension.  The quotient is the coarsest reachable
deterministic predictive realization for the specified futures.

To obtain the operator-algebraic history object, the input is strengthened to
complete multitime intervention tables on specified finite operator frames.
The reconstructed Choi tensors must be positive and satisfy the causal prefix
identities; nonnegative entries in a probability table alone do not replace
these conditions.  The external intervention ports are part of this operational
input.  The persistent memory representation is reconstructed using the coherent
isometric realization that is minimal at every prefix.

Its reduced branches and their formal Hilbert--Schmidt adjoints generate a
represented process algebra.  The formal adjoints are algebraic operations,
not necessarily executable reverse channels.  Quotienting the two-sided ideal
invisible to all retained readout contexts gives the future-separated history
algebra $\A_X^{\mathrm{hist}}$.  Its intrinsic normalized regular trace is
\begin{equation}
 \widehat\tau_{X,\mathrm{reg}}(a)
 =\frac{\operatorname{Tr}_{\C}(L_a)}{\dim_{\C}\A_X^{\mathrm{hist}}},
 \qquad
 \mathbb K_{X,r}(v,w)
 =\widehat\tau_{X,\mathrm{reg}}([v]^*[w]),\quad |v|,|w|\le r.
 \label{eq:overview-history-gram}
\end{equation}
The Gram kernels record the word relations.  When two successive word spans
have equal rank, the finite panel determines the multiplication, involution,
and all source words; see Theorem~\ref{thm:ncg-flat-history-reconstruction}.
These are exact finite-data statements, rather than a statistical recovery
claim from an incomplete or noisy table.

\begin{definition}[Finite Renewal Grand Tensor]
\label{def:ncg-grand-tensor}
The finite Renewal Grand Tensor associated with the operational process $X$ is
\begin{equation}
 \boxed{
 \GTX
 =\left(
 \A_X^{\mathrm{hist}},
 \{\mathbb K_{X,r}\}_{r\ge0},
 Z_X^{\min}
 \right).}
\label{eq:ncg-grand-tensor}
\end{equation}
It is a complete positive finite history-moment object: $Z_X^{\min}$ records
which past distinctions remain predictively effective, $\A_X^{\rm hist}$ records
how future-visible represented histories compose, and the Gram hierarchy records
their intrinsic positive overlaps.
\end{definition}

The reconstruction in Proposition~\ref{prop:ncg-grand-tensor-reconstruction}
therefore supplies this finite history object.  Conversely, the boundary-relative
realization and recognition results retain representation multiplicities:
preparations and terminal readouts remain external contexts, and are not silently
adjoined to the internal history algebra.  The separate dimension-minimal
recognition argument is not a uniqueness theorem for arbitrary physical
implementations of the coherent process representation.

\paragraph{First-order differential interface.}

The Grand Tensor is logically upstream of spectral geometry, but it is important
not to identify algebraic reconstruction with metric differentiation.  From
\eqref{eq:ncg-grand-tensor} we canonically obtain the history algebra and its
positive finite-dimensional structure.  We take
\begin{equation}
 \A=\A_X^{\mathrm{hist}},
 \qquad
 \tau=\widehat\tau_{X,\mathrm{reg}},
\label{eq:ncg-grand-algebra-trace}
\end{equation}
or, more generally, any declared faithful tracial state on $\A$.  To spectralize
this algebra one additionally specifies a compatible first-order variation.
Abstractly, this is a Hilbert $\A$-bimodule $\mathcal K$, an antiunitary
involution $C_{\mathcal K}$ reversing the bimodule actions, and a complex-linear
map
\begin{equation}
 \delta:\A\longrightarrow\mathcal K
\label{eq:ncg-grand-delta}
\end{equation}
satisfying
\begin{equation}
 \delta(\one)=0,
 \qquad
 \delta(ab)=a\delta(b)+\delta(a)b,
 \qquad
 C_{\mathcal K}\delta(a)=\delta(a^*).
\label{eq:ncg-grand-delta-rules}
\end{equation}
When this datum comes from a Renewal process, $\delta$ represents the declared
first-order variation of the reconstructed predictive history algebra.  The
results below do not assume that $\GTX$ or $\A$ alone uniquely determines this
variation.  This distinction is essential: the no-scale result below shows that
rescaling $\delta$ changes the spectral metric while leaving the algebra and trace
fixed.

The logical interface is therefore
\begin{equation}
\boxed{
\text{operational histories}
\longrightarrow
\GTX
\longmapsto
(\A_X^{\mathrm{hist}},\tau_X)
\overset{\delta_X}{\longrightarrow}
(\A_X,\Hh_X,D_X,J_X,\gamma_X).}
\label{eq:ncg-grand-spectral-interface}
\end{equation}
The first arrows reconstruct predictive and algebraic structure; the final arrow
is the Hodge--Dirac spectralization studied in this paper.

The zero-form carrier is
\begin{equation}
\Hh_0=L^2(\A,\tau),
\label{eq:summary-H0}
\end{equation}
and the one-form carrier is the minimal quotient of the universal calculus
generated by $\delta$,
\begin{equation}
\Hh_1(\delta)=\Om_u^1(\A)/\Ker T_\delta,
\qquad
\Om_u^1(\A)=\Ker(m:\A\otimes\A\to\A).
\label{eq:summary-H1}
\end{equation}
On $\Hh=\Hh_0\oplus\Hh_1(\delta)$, the induced derivation $\partial$ gives
\begin{equation}
D=\begin{pmatrix}0&\partial^*\\ \partial&0\end{pmatrix},
\qquad
\gamma=\begin{pmatrix}I&0\\0&-I\end{pmatrix},
\label{eq:summary-D}
\end{equation}
while reversal of universal one-forms gives the real structure $J$.  Thus the
spectral triple is constructed from a first-order differential on an algebra that
has itself already been reconstructed from predictive histories.

\subsection{What the differential determines}

The finite Hodge construction is canonical once $\delta$ is fixed.  It produces a
real even spectral triple with exact order-zero and first-order relations.  The
spectral scale is not already hidden in $(\A,\tau)$: replacing $\delta$ by
$s\delta$ leaves the algebra and trace unchanged but sends
\begin{equation}
D\longmapsto sD,
\qquad
L_D\longmapsto sL_D,
\qquad
d_D\longmapsto s^{-1}d_D.
\label{eq:summary-scale}
\end{equation}
Thus an absolute metric scale enters through the differential data.

There is a second, finite ambiguity.  For a marked spectral packet
$\mathsf S=(\A,\Hh,\pi,D,J,\gamma)$ define
\begin{equation}
\mathfrak c_{\mathsf S}
=\{C\in\pi(\A)'_{\rm sa}: C\gamma+\gamma C=0,\ JC=\varepsilon'CJ\}.
\label{eq:summary-commutant}
\end{equation}
All compatible self-adjoint Dirac operators implementing the same commutator
derivation form exactly the affine space
\begin{equation}
D+\mathfrak c_{\mathsf S}.
\label{eq:summary-affine}
\end{equation}
The orthogonally centered representative
$D_{\rm der}=D-P_{\mathfrak c_{\mathsf S}}D$ is the unique
minimum-Hilbert--Schmidt-norm implementer, and
$q_{\mathsf S}=\dim_{\R}\mathfrak c_{\mathsf S}$ real linear scalar measurements are
necessary and sufficient to distinguish every point of the fibre.

The austere Hodge construction has an additional visible restriction:
$D(\one,0)=0$.  It therefore cannot be surjective onto arbitrary finite Dirac
operators.  This is why the inverse statement below uses an enlarged calibrated
marked reconstruction rather than the bare Hodge functor.

\subsection{Which spectral geometries can be reconstructed?}

At the finite marked level, the represented algebra and symmetry marks are retained
together with the commutator derivation and exactly the independent real linear coordinates required
to resolve its compatible commutant fibre.  If
$q_{\mathsf S}=\dim_{\R}\mathfrak c_{\mathsf S}$, then $q_{\mathsf S}$
independent real linear coordinates are sufficient and, in general, necessary.  At the strong
compact level, finite spectral screens converge strongly on fixed vectors,
resolvents, commutators, and matrix coefficients.  At the norm-multiplicative spectral-compression
level, multiplication must additionally be transported asymptotically in norm.

\begin{theorem}[Finite Hodge--Dirac reconstruction and commutant coordinates]
\label{thm:summary-finite}
The following statements hold.
\begin{enumerate}[label=\textup{(\roman*)},leftmargin=2.8em]
\item Every finite real differential datum $(\A,\tau,\delta)$ determines the real even
Hodge--Dirac spectral triple above, functorially under trace- and
 differential-preserving $*$-isomorphisms.
\item The algebra and trace do not determine the Dirac scale, and a fixed commutator
 derivation determines the full compatible Dirac operator exactly modulo
 $\mathfrak c_{\mathsf S}$.
\item For the enlarged calibrated marked reconstruction,
\begin{equation}
\begin{aligned}
\EssIm(\mathfrak S_{\rm fin}^{\rm cal})
  &=\mathbf{Spec}_{\rm fin},\\
\EssIm(\mathfrak S_{\rm cpt}^{\rm strong,cal})
  &=\mathbf{Spec}_{\rm cpt}^{\rm sep},\\
\EssIm(\mathfrak S_{\rm cpt}^{\rm norm,cal})
  &=\mathbf{Spec}_{\rm cpt}^{\rm sqd}
  \subsetneq\mathbf{Spec}_{\rm cpt}^{\rm sep}.
\end{aligned}
\label{eq:summary-essential-image}
\end{equation}
Thus every finite spectral packet is reconstructible from its represented
commutator data plus exactly $q_{\mathsf S}$ independent real linear commutant coordinates; every
separable compact-resolvent packet is a strong cofinal limit of finite anchored
screens; and norm-multiplicative spectral-compression realizability is equivalent to spectral
quasidiagonality.
\end{enumerate}
\end{theorem}

The distinction between the austere Hodge image and the anchored reconstruction is
essential.  The commutator derivation fixes a canonical centered implementer, while
the $q_{\mathsf S}$ anchors recover precisely the finite component invisible to
commutators; no full generator germ is retained.  The complete definitions, sharp
minimality statement, and essential-image proof are given in
\Cref{sec:finite} of the mathematical supplement.


% =============================================================================

At the compact-resolvent level, ``compression'' refers to the represented
matrices and the convergence conditions of
Definition~\ref{def:reconstruction-categories}.  The compressed observables are
generally an operator system rather than a representation of the original
algebra, and are not asserted to inherit every real first-order axiom.
The strong and norm-multiplicative assertions are separated in
Theorem~\ref{thm:spectral-compressions}.

\section{Independent metric reconstruction}
\label{sec:a3-main}
% =============================================================================

Producing an operator called $D$ is not by itself evidence that geometry has emerged.
A stronger test is available when the same finite dynamics determines a metric by a
route that does not use the Dirac operator.

\subsection{An independent metric test}

For a stationary finite graph, let $m_x>0$ be the stationary masses,
$q_{xy}=m_xk_{xy}$ the stationary fluxes, and
\begin{equation}
c_{xy}=\frac{q_{xy}+q_{yx}}2,
\qquad
j_{xy}=\frac{q_{xy}-q_{yx}}2.
\label{eq:summary-cj}
\end{equation}
The symmetric conductances determine the graph Hodge--Dirac operator $D_{m,c}$.
Exactly,
\begin{equation}
D_{m,c}^2|_{\ell^2(V,m)}=-2\mathcal L_s,
\qquad
\|[D_{m,c},f]\|^2
=\max_x\frac1{m_x}\sum_yc_{xy}|f(y)-f(x)|^2.
\label{eq:summary-graph-identities}
\end{equation}
Thus the Dirac square on zero-forms is $-2$ times the symmetric Markov generator, and its commutator seminorm
is the local quadratic form of that generator.

For the periodic root graphs, let
$W_0=\{x\in\R^4:\sum_i x_i=0\}$,
$\Lambda_{A_3}=W_0\cap\Z^4$, and
$V_h=h\Lambda_{A_3}/\Lambda_{A_3}$, where $h=d^{-1}$ and $d\ge3$.
The twelve oriented roots are $\Phi_{A_3}=\{e_i-e_j:i\ne j\}$.
With masses $h^3$ and root-edge rates $1/(8h^2)$, their jump covariance is
\begin{equation}
 Q_h=\sum_{\alpha\in\Phi_{A_3}}
       \frac{(h\alpha)\otimes(h\alpha)}{8h^2}
     =I_{W_0},\qquad g=Q_h^{-1}.
 \label{eq:overview-A3-covariance}
\end{equation}
This calculation specifies the flat metric without using $D_h$.
The commutator normalization then gives
\begin{equation}
L_h(f)^2
=\max_{x\in V_h}\frac1{8h^2}
\sum_{\alpha\in\Phi_{A_3}}|f(x+h\alpha)-f(x)|^2,
\label{eq:summary-A3-Lip}
\end{equation}
with the following uniform metric conclusion.

\begin{theorem}[Uniform $A_3$ metric recovery]
\label{thm:summary-metric}
The pure-state Connes distance satisfies
\begin{equation}
\sup_{x,y\in V_h}
|d_h^{\rm Connes}(x,y)-d_g(x,y)|\longrightarrow0,
\label{eq:summary-A3-distance}
\end{equation}
where $d_g$ is the flat geodesic distance determined by the jump covariance.
\end{theorem}

This is the principal independent metric test of the paper.  The proof in
Section~\ref{subsec:A3-distance} controls the entire discrete Lipschitz ball,
not merely samples of smooth functions: interpolation gives compactness, the
quadratic root constraint identifies the limiting unit gradient bound, and
uniformly smoothed distance functions give the reverse estimate.

\paragraph{What the limit does and does not assert.}
The theorem compares vertices, that is, pure states.  It is not by itself a
convergence theorem for the full quantum state-space Lip-norms or for the
Hodge operator in norm resolvent.  Indeed, retaining both oriented edges gives
\begin{equation}
 \dim\Ker D_{m,c}=2|E|-|V|+2,
 \label{eq:overview-hodge-kernel}
\end{equation}
which is $11|V|+2$ on the periodic $A_3$ root graph.  A compatible limit retaining
all these zero modes cannot have compact resolvent.  The spin construction in
Section~\ref{sec:extensions-main} uses a different, additionally marked operator
and does not remove this obstruction by assertion.

There is also an independent endpoint-block construction on every closed connected
Riemannian manifold.  Its seminorm is the ordinary graph Lipschitz seminorm and its
Connes distance is exactly the graph shortest-path metric; refining sufficiently
dense meshes gives uniform geodesic convergence
(Theorem~\ref{thm:supp-global-metric}).  This is a complementary construction,
not an extension of the quadratic $A_3$ Hodge theorem by the same proof.

\section{Spectral geometry as an information-reduced description}
\label{sec:universality-main}

The spectral map is not injective.  The stationary-current theorem makes the
loss of directional information exact, while positive dynamical memory and
marked monodromy give two complementary inverse problems.  Keeping these
reductions distinct is essential: they forget different data and obey different
coarse-graining laws.

\subsection{The stationary-current fibre}

For a connected graph with fixed masses $m$ and conductances $c$, write
\begin{equation}
 \overline{\mathcal J}(G,c)=\{j\in\Ker B:|j_e|\le c_e\},\qquad
 \mathcal J(G,c)=\{j\in\Ker B:|j_e|<c_e\}.
 \label{eq:summary-current-fibre}
\end{equation}
The closed polytope is the full nonnegative-rate stationary realization fibre.
Its relative interior $\mathcal J(G,c)$ is the fibre with strictly positive
flux in both directions, and has dimension $b_1(G)=|E|-|V|+1$.
Every point has the same metric Hodge--Dirac triple, although directed path
probabilities, time-reversal asymmetry, and entropy production may differ.

For a strictly positive realization, its edge affinity is
\begin{equation}
 a_e=\operatorname{arctanh}\frac{j_e}{c_e}
     =\frac12\log\frac{q_e}{q_{\bar e}}.
 \label{eq:summary-affinity}
\end{equation}
Adding a vertex potential $B^{\mathsf T}\phi$ changes the representative but
not the cohomology class, or equivalently its sums around cycles.  The
nonlinear stationarity equation selects precisely one representative from
each such class.

\begin{theorem}[Cohomological reconstruction of directed stationary dynamics]
\label{thm:summary-affinity}
For fixed connected $(G,m,c)$ with $c_e>0$, the map
\begin{equation}
 \mathfrak A:\mathcal J(G,c)\longrightarrow H^1(G;\R),\qquad
 j\longmapsto[\operatorname{arctanh}(j/c)]
 \label{eq:overview-affinity-diffeomorphism}
\end{equation}
is a diffeomorphism.  For a representative $a_0$ of a class $\alpha$, minimize
\begin{equation}
 \sum_e c_e\log\cosh\bigl(a_{0,e}+(B^{\mathsf T}\phi)_e\bigr)
 \quad\text{over }\sum_x\phi_x=0.
 \label{eq:overview-affinity-minimization}
\end{equation}
The unique minimizer gives $a_\alpha=a_0+B^{\mathsf T}\phi_\alpha$,
$j_\alpha=c\tanh a_\alpha$, and the complete directed generator
$k_{xy}=(c_{xy}+j_{\alpha,xy})/m_x$.
\end{theorem}

The full proof, including coercivity, strict convexity, the positivity domain,
and smooth dependence, is Theorem~\ref{thm:supp-nonlinear-Hodge}.
The variational selection specializes nonlinear Hodge theory
\cite{PardoGuerraThapaWashburn2026}; the reconstruction statement identifies
that selector with the entire fixed-traffic stationary Markov fibre.
The affinity is half the logarithmic forward/reverse stationary-flux ratio;
this convention fixes the factors in the entropy formula.

In particular, geometry need not determine irreversibility.  On a three-cycle
with equal conductance $c$, every constant current $j_e=s$, $|s|<c$, gives the
same metric triple, whereas
\begin{equation}
 \sigma(s)=6s\log\frac{c+s}{c-s}
 \label{eq:overview-cycle-entropy}
\end{equation}
changes with $s$.  More generally $\sigma(j)=4\sum_e j_ea_e$, so the
current--affinity pairing is invariant under adding an exact cochain.
The additional cohomology class restores the directed generator on the
strictly positive fibre without changing its metric triple.

\subsection{Positive memory and global monodromy}

Hidden positive memory gives a second inverse problem.  Eliminating auxiliary
degrees of freedom produces a Stieltjes self-energy
\begin{equation}
 \Sigma_\mu(z)=\int_{(0,\infty)}\frac{1}{\lambda-z}\,\dd\mu(\lambda)
             =B_\mu(C_\mu-z)^{-1}B_\mu^*.
 \label{eq:summary-Stieltjes}
\end{equation}
The full operator-valued measure determines the source-minimal positive dilation,
whereas a static effective operator may retain only the inverse moment
$\Sigma_\mu(0)=\int\lambda^{-1}\,\dd\mu(\lambda)$.
Section~\ref{app:memory-details} distinguishes static data from the full
self-energy, including examples where all integer moments agree but the
self-energies differ.  Section~\ref{app:positive-dilation} retains the explicit
minimal dilation construction and proof.

The third inverse problem is global.  Locally marked-isomorphic spectral triples
with specified edge identifications on a connected graph are classified by
\begin{equation}
 [\rho]\in\Hom(\pi_1(\Gamma),G_{\mathsf S})/G_{\mathsf S}.
 \label{eq:summary-monodromy}
\end{equation}
Trivial monodromy gives a marked trivialization compatible with those
identifications; nontrivial monodromy obstructs it.  This concerns marked flat
transport, not a claim that every underlying unmarked vector bundle is
nontrivial.  The complete classification is in Section~\ref{app:marked-monodromy}.

\subsection{Componentwise coarse graining}

These inverse data have different natural effective maps.  Current is pushed
through the graph incidence chain map; a positive-memory cutoff removes the
high-energy measure while retaining its static Schur correction; and monodromy
descends only when the contracted cycles have trivial holonomy.

\begin{theorem}[Coarse graining of the core inverse data]
\label{thm:summary-universality}
Under the separate quotient, integrability, and descent hypotheses of
Theorems~\ref{thm:supp-current-chain}, \ref{thm:supp-memory-RG},
and~\ref{thm:supp-monodromy-descent}, current aggregation contracts entropy
production, positive-memory cutoff preserves the zero-energy Schur complement,
and marked monodromy descends exactly when
$\Ker q_*\subseteq\Ker\rho$.  These operations compose componentwise under
compatible graph quotients, nested thresholds, and nested effective-space
isometries.
\end{theorem}

Surjectivity on linear cycle spaces does not assert surjectivity of the bounded
current polytopes.  Nor do these three examples alone exhaust the inverse
fibre of an arbitrary predictive system.  The additional joint-realization
and completed-fibre framework is stated separately in
Section~\ref{app:completed-fibres}, with its retained-coordinate and completion
hypotheses explicit.

\section{Further mathematical extensions}
\label{sec:extensions-main}

The preceding graph theorems and their finite framework form the common
mathematical core.  The extended version also retains results on continuum spin
operators, indefinite memory, and joint realization spaces.  Their proofs are
collected in Part~II of the Mathematical Supplement.  None is a prerequisite
for the core $A_3$ metric or stationary-current theorem.

\subsection{From local spin discretization to a global convergence criterion}

The spin problem is subtler.  Local positive stencils reconstruct the inverse metric,
and a density-symmetric covariant Wilson discretization removes lattice doublers while
remaining $O(h)$ on sampled smooth spinors.  These chartwise estimates do not by
themselves guarantee that independently discretized charts assemble into one stable
global operator.  The mathematical supplement therefore isolates the additional
global approximation properties in the notion of a \emph{compatible stable spin
discretization}.

\begin{theorem}[Additional spin continuum results]
\label{thm:summary-spin}
The following statements hold.
\begin{enumerate}[label=\textup{(\roman*)},leftmargin=2.8em]
\item The covariant Wilson construction has one low-energy species and satisfies the
local uniform graph estimates required for a spin continuum limit.
\item On a flat periodic spin torus with constant metric and the periodic spin
structure, the uniform covariant Wilson discretization satisfies the global
stability and core-consistency conditions and converges in generalized
norm-resolvent sense to the doubled continuum Dirac operator.
\item Let $(M^d,g,\mathfrak s)$ be closed.  If the finite covariant Wilson operators
form a compatible stable spin discretization in the precise sense of
Definition~\ref{def:stable-spin-atlas}, then for every $z\in\C\setminus\R$,
\begin{equation}
\left\|W_h(D_h-z)^{-1}W_h^*
-(\widehat D_{M,\mathfrak s}-z)^{-1}\right\|\longrightarrow0,
\qquad
\widehat D_{M,\mathfrak s}=D_{M,\mathfrak s}\otimes\sigma_1.
\label{eq:summary-spin-resolvent}
\end{equation}
\end{enumerate}
\end{theorem}

The theorem separates three issues that should not be conflated: the extensive
kernel of the unreduced universal graph calculus, the Wilson control of
the spinor sector, and the additional atlas compatibility needed on a general
closed manifold.  Flat periodic spin tori provide an unconditional global spin
class; the general norm-resolvent statement is a convergence criterion for
compatible stable assemblies.  The scalar metric approximation is unconditional.
The precise local estimates, atlas definition, periodic corollary, and compactness
argument are given in \Cref{app:compact-spin}.


% =============================================================================

The spin lift, global spin structure, and stable reconstruction maps are
additional inputs.  The Wilson family is a separate discretization of the doubled
spin Dirac operator; it is not identified with the unreduced one-step graph Hodge
operator.  Its unique low-energy lattice species is understood within that fixed
doubled Clifford convention.

\subsection{Indefinite memory beyond the positive branch}

For normalized rational Hermitian self-energies, the coefficient-Hankel
construction gives a finite source-minimal Pontryagin realization.  Its
ordinary dimension is the stabilized Hankel rank and its negative index is
the number of negative squares of the corresponding Nevanlinna kernel
(Theorem~\ref{thm:supp-complete-rational-Krein}).  The pole--Hankel invariants and
the finite descriptor extension for polynomial parts are retained in
Section~\ref{app:rational-krein} rather than incorporated into the positive
Stieltjes theorem.  They use their own indefinite inner-product and realization
conventions; no arbitrary positive-memory model is assigned a negative index.

\subsection{Joint realizations and completed inverse fibres}

Prescribed marginal realizations need not have a common realization merely
because their pairwise overlaps are admissible.  Positive block correlation
operators classify minimal common Gram realizations, with support rank giving
the minimum common dimension and a stabilizer recording the remaining gauge
freedom (Theorem~\ref{thm:supp-common-source}).  The three-block example in
Proposition~\ref{cth:supp-pairwise-not-global} makes the obstruction explicit.

The broader finite-fibre and completion statements retain this joint Gram
coordinate together with the declared current, memory, marking, calibration,
and monodromy data.  They concern the specified fully reduced
retained-coordinate class, not an unconditional exhaustive classification of
all microscopic predictive dynamics.  Compactification, compatible restriction
maps, resolvent and weight data, and cofinal separation enter the completion
hypotheses.  Within that setting, the stable-image pruning lemma distinguishes
finite-stage points that extend indefinitely from those that do not.  The
resulting inverse-limit description is kept in
Section~\ref{app:completed-fibres} with these assumptions rather than promoted
to a consequence of the two graph theorems.

\section{Interpretation and discussion}
\label{sec:discussion}
% =============================================================================

The results separate predictive reconstruction from operator-theoretic
spectralization in a useful way.  Renewal Geometry first reconstructs the Grand
Tensor $\GTX$, including the future-minimal predictive carrier, represented
history algebra, and positive moment hierarchy.  Spectralization then equips that
reconstructed algebra with a compatible first-order differential datum
$\delta_X$.  Once $(\A,\tau,\delta)$ is given, the finite Hodge--Dirac theorem
is an ordinary statement about $C^*$-algebras, Hilbert bimodules, and derivations.
The predictive reconstruction is therefore logically prior to the spectral
triple without being an extra hypothesis of the abstract operator-theoretic
result.

\subsection{Spectral geometry as an effective description of predictive dynamics}

The first consequence is a change in the logical status of the spectral triple.
In Connes' formulation, the algebra, its Hilbert-space representation, and the
Dirac operator constitute the spectral data from which geometry is extracted.  In
the present construction these objects can instead occur downstream of the
Renewal Grand Tensor:
\begin{equation}
\text{predictive dynamics}
\longrightarrow
\GTX
\longmapsto
(\A_X^{\mathrm{hist}},\tau_X)
\overset{\delta_X}{\longrightarrow}
(\A_X,\Hh_X,D_X,J_X,\gamma_X).
\label{eq:discussion-layered}
\end{equation}
A spectral triple therefore need not be interpreted as the most microscopic
object in the description.  Within the framework considered here, it is the
spectral image of a specific algebraic and differential sector of a richer
predictive object.

This does not mean that an arbitrary Dirac operator is determined by the algebra.
The finite Hodge construction shows precisely where additional information enters.
The algebra and faithful trace do not fix an absolute metric scale, and the
commutator derivation determines the full compatible Dirac operator only modulo
the appropriate self-adjoint commutant.  Likewise, finite marked surjectivity uses exactly
$q_{\mathsf S}=\dim_{\R}\mathfrak c_{\mathsf S}$ independent real linear commutant coordinates:
the derivation fixes the centered implementer and the anchors recover precisely the
remaining compatible affine component.  It is not a claim that an arbitrary target
Dirac operator has been derived from the austere Hodge datum alone.  At the
norm-multiplicative spectral-compression level the reconstruction becomes genuinely selective, with
spectral quasidiagonality providing the exact criterion.

The independent metric construction is important for the same reason.  On the
periodic $A_3$ family, the Riemannian metric reconstructed from jump second moments
and the Connes distance reconstructed from the Hodge--Dirac operator are obtained
by different routes and converge to the same continuum geometry.  The agreement
therefore tests the spectralization rather than defining the target metric by
construction.  For flat periodic spin tori the Wilson construction also gives unconditional
norm-resolvent convergence.  For general closed Riemannian manifolds the scalar
metric approximation remains unconditional, while the spin norm-resolvent statement
requires the compatible stable-discretization hypotheses stated above.

\paragraph{A common predictive origin for two geometric constructions.}

The independent Lorentzian construction
\cite{PelissierRenewalSpacetime2026} uses the same Renewal Grand Tensor $\GTX$
to reconstruct Lorentzian spacetime without requiring a spectral triple as an
intermediate primitive.  Together, the two constructions therefore place the
Grand Tensor as a common predictive object upstream of two distinct downstream
descriptions, as summarized in
\eqref{eq:intro-renewal-geometric-branches}.

This matters conceptually.  If spectral geometry were the only downstream output,
one could interpret the predictive construction merely as an alternative
procedure for generating spectral triples.  The independent spacetime branch
instead shows that Lorentzian geometry and Connes-type spectral geometry can be
different projections of a common predictive structure rather than successive
fundamental layers.

The claim is structural rather than ontological.  The results do not establish
that Renewal Geometry is the unique microscopic theory underlying all
noncommutative geometries or all physical spacetimes.  They show that, within the
framework studied here, neither spectral geometry nor spacetime geometry needs to
be placed at the primitive entrance.  Both may arise from a predictive structure
defined before either metric geometry or a spectral triple is introduced.

\paragraph{Information loss and nonfaithfulness.}

The inverse results sharpen this effective interpretation.  Spectralization is
generically many-to-one:
\begin{equation}
\mathfrak S(\mathcal R_1)=\mathfrak S(\mathcal R_2)
\qquad\text{need not imply}\qquad
\mathcal R_1\simeq\mathcal R_2,
\label{eq:discussion-nonfaithful}
\end{equation}
where $\mathcal R_i$ denote predictive or dynamical realizations and
$\mathfrak S$ their spectral image.  A spectral triple is therefore not, in
general, a complete invariant of the underlying dynamics.

The stationary-current fibre provides the simplest example.  Fixing the masses
and symmetric conductances fixes the metric Hodge--Dirac triple, while a
divergence-free stationary current remains undetermined.  Moving within this
fibre changes path probabilities, time-reversal asymmetry, and entropy production
without changing the spectral metric.  The spectral geometry retains the
symmetric metric sector while forgetting part of the directed dynamics.

The memory fibre exhibits a different form of information loss.  Eliminating
auxiliary degrees of freedom produces a frequency-dependent self-energy whose
minimal realization is encoded by an operator-valued measure.  A static spectral
description may retain only selected boundary values, moments, or an effective
operator, so distinct dynamical tails can have the same static spectral image.
What is lost here is temporal organization rather than merely a finite parameter.

Global marked monodromy provides a third mechanism.  Locally equivalent spectral
packets can differ in how their marked realizations are assembled globally.
Consequently spectralization can forget directional, dynamical, and global
information for mathematically different reasons.

These examples motivate the description of spectral geometry as an
\emph{information-reduced effective geometry}.  This should not be interpreted as
a claim that noncommutative geometry is mathematically incomplete.  A spectral
triple may completely encode the geometric observables for which it is designed.
The point is that it need not completely encode the predictive dynamics from which
it arose.  A single spectral geometry can therefore represent an entire family of
dynamically inequivalent microscopic realizations.

The analogy with effective descriptions in statistical physics is useful.  A set
of macroscopic variables may completely characterize an effective state while
failing to determine a unique microscopic realization.  Likewise, a spectral
triple may fully characterize an effective metric and differential sector while
identifying systems that remain distinct at the predictive level.

\paragraph{Why the quotient has a universality interpretation.}

The inverse fibre
\begin{equation}
\mathfrak S^{-1}(\mathsf S)
\label{eq:discussion-fibre}
\end{equation}
is therefore more than a technical obstruction to inversion.  It describes the
microscopic information quotiented out in the passage to a fixed spectral packet.

Moreover, the discarded variables are themselves organized under effective
reduction.  Directed currents are pushed through the induced graph chain map;
high-energy memory can be absorbed into a zero-energy contact correction while
the retained low-energy measure survives; and marked monodromy descends exactly
when the cycles removed by the quotient carry trivial holonomy.  These maps
compose under successive coarse grainings.

This is the sense in which the spectral packet acts as a universality-level
observable.  Different microscopic systems can share the same spectral geometry
because distinctions that remain meaningful at the predictive level no longer
affect the operator-geometric structure retained at that level of description.

The coarse-graining language should nevertheless be interpreted precisely.  The
results do not prove that spectralization is universally generated by a
renormalization-group flow, nor that arbitrary predictive dynamics converges
toward a spectral fixed point.  What is established is a controlled many-to-one
map, explicit inverse fibres, and composable effective maps on their principal
coordinates.  The universality interpretation is therefore structural rather than
a claim of a universal dynamical attractor.

\subsection{Inverse fibres, scope, and open problems}

Once the lost information has been identified, a natural question is whether it
should be retained when it is physically relevant.  The present examples suggest
several possible enrichments.  A metric spectral triple can be supplemented by an
affinity or current class to retain irreversible directed information; an
operator-valued self-energy can retain frequency-dependent memory; and a marked
holonomy class can retain global assembly data.  Schematically,
\begin{equation}
(\A,\Hh,D,J,\gamma)
\quad\longrightarrow\quad
(\A,\Hh,D,J,\gamma;
\text{current},\text{memory},\text{monodromy},\ldots).
\label{eq:discussion-enrichment}
\end{equation}
The ordinary spectral packet is recovered by forgetting the additional
coordinates.

This perspective changes the natural inverse problem.  Rather than expecting one
microscopic dynamics to correspond to one spectral geometry, one should generally
expect a moduli space or stratified fibre of realizations.  Under their declared retained-coordinate and completion hypotheses, the
finite stabilizer quotient and its completed inverse-limit version provide
examples of such a description.  The
natural question is therefore not to identify a unique dynamics represented by a
spectral triple, but to classify the predictive dynamics that share it as a common
effective geometry.

The independent Lorentzian construction
\cite{PelissierRenewalSpacetime2026} makes this inverse problem richer.  Two
predictive systems with the same spectral image need not have the same image under
a distinct downstream geometric reconstruction; conversely, systems may agree on
a spacetime sector while differing in spectral or dynamical refinements.  Joint
fibres of the spectral and Lorentzian projections therefore provide a natural
object for future study.

\paragraph{Mathematical scope.}

Several limitations of the present results should remain explicit.  The strict
spectral image would benefit from an intrinsic characterization independent of a
chosen spectral projection sequence.  For spin geometry, the local Wilson
construction provides the required analytic estimates and the flat periodic class
closes unconditionally, but an unconditional global realization theorem for every
closed spin manifold would require verification of the compatibility and stability
hypotheses in general.  Likewise, the extensive kernel theorem shows that any
cellular or relation-adapted reduction of the universal one-step graph calculus
requires its own continuum certificate; no such reduction theorem is assumed here.  Noncompact and genuinely
semifinite geometries require further control of domains, localization, and
weights.  The positive Stieltjes memory classification has a finite
rational indefinite extension in Part~II, but the unrestricted nonrational
infinite-dimensional indefinite problem is not treated here.

Nor does the finite compact-quantum-metric-space result by itself imply
quantum Gromov--Hausdorff, propinquity, or spectral-propinquity convergence in the
continuum limit \cite{Rieffel2004,Latremoliere2022,FarsiLatremolierePacker2024}.
Such statements require uniform control of the full state-space Lip-norms and, for
spectral propinquity, compatible control of the spectral data beyond the
pure-state distance convergence proved for the $A_3$ approximants.

\paragraph{Open questions and broader interpretation.}

The effective-geometric viewpoint raises several further questions.  Which parts
of the inverse fibre can be reconstructed from natural enrichments of a spectral
triple?  Which fibre coordinates remain stable under broader classes of
coarse graining?  Can additional dynamical principles select canonical
representatives inside a spectral fibre?  And how much of the many-to-one
structure persists in continuum classes substantially larger than the finite and
controlled compact-resolvent settings considered here?

The coexistence of the spectral and spacetime branches suggests a second family
of questions.  One may ask when the two downstream reconstructions determine one
another, when one retains information absent from the other, and which microscopic
coordinates are invisible to both.  More generally, one can ask whether other
geometric or physical structures arise as additional quotients of the same
predictive object.

Within the scope established in this paper, the hierarchy is nevertheless clear.
Predictive dynamics carries temporal, directional, and global information before
a metric or spectral triple is introduced.  The Renewal Grand Tensor $\GTX$
organizes the future-minimal carrier, represented history algebra, and positive
moment structure.  Lorentzian spacetime can arise from one downstream sector,
while on the spectral branch the history algebra becomes Connes-type spectral
geometry after a compatible first-order variation is specified.  The resulting
spectral packet retains the metric and differential structures relevant to its
level while identifying classes of underlying dynamics that remain distinct
microscopically.

Thus, within this framework, Connes-type spectral geometry is naturally
interpreted as a stable effective geometry of a richer predictive dynamics:
powerful and potentially complete for its geometric sector, but not in general a
complete description of the dynamics from which that geometry emerges.

\section*{Formal Verification and Code Availability}
All definitions, theorems, and proofs in this work have been formally verified in \texttt{Lean~4} using \texttt{Mathlib} without unproven assumptions (\texttt{sorry}-free). The complete source code, build configuration, and documentation are available at \url{https://github.com/AI-MathPhys/renewal-geometry}.

\clearpage
\phantomsection
\addcontentsline{toc}{section}{Mathematical Supplement}
\mathsupplementtrue
\appendix
\renewcommand{\thesection}{S\arabic{section}}
\renewcommand{\theHsection}{supp.\arabic{section}}
\setcounter{section}{0}

\begin{center}
 {\LARGE\bfseries Mathematical Supplement}
\end{center}
\medskip

The supplement is divided into two parts.  Part~I contains the self-contained
mathematical core of the paper.  It first develops the finite Hodge--Dirac
construction from a real bimodule derivation, including the scale ambiguity of
the resulting metric, the affine commutant fibre of compatible Dirac
implementers, and the distinction between strong and norm-multiplicative
spectral compression.  It then treats stationary Markov graphs.  For the
covariance-normalized periodic $A_3$ family, the graph commutator seminorm is
identified explicitly and the associated pure-state Connes distances are proved
to converge uniformly to the flat Riemannian metric reconstructed independently
from the jump covariance.  At fixed vertex masses and symmetric conductances,
the directed stationary realizations form a cycle-current fibre whose strictly
positive interior is globally parametrized by
$H^1(G;\mathbb R)$ through the affinity class
$[\operatorname{arctanh}(j/c)]$; the inverse is obtained by a weighted
$\log\cosh$ minimization, and graph quotients induce the corresponding
current and entropy-contraction maps.

The remaining sections of Part~I supply the mathematical structures supporting
these results.  They include an independent endpoint construction recovering
Riemannian distance from finite metric triples, edge-ratio and cellular Hodge
descriptions of the affinity class, operational reconstruction and recognition
of the finite history algebra from complete intervention data, positive
operator-valued Stieltjes memory and its static-preserving cutoff maps, and the
classification and descent of marked monodromy.  These results are stated with
their own hypotheses and are logically separate from the proofs of the two
principal graph theorems.

Part~II collects further mathematical extensions that are not required for the
core results.  It develops positive local stencils and covariant Wilson
discretizations for spin geometry, including unconditional norm-resolvent
convergence on flat periodic spin tori and a general convergence theorem under
explicit compatible-stability hypotheses.  It also extends the positive-memory
theory to finite rational Hermitian dynamics through minimal
Pontryagin/descriptor realizations, and studies joint and completed inverse
fibres using correlation-Gram data, stabilizer quotients, and compact
finite-stage inverse systems.  Thus the spin-discretization,
indefinite-memory, and completed-fibre results enlarge the scope of the arXiv
version without entering the assumptions or proofs of the finite
Hodge--Dirac, $A_3$ metric, or stationary-current reconstruction theorems.

\supplementtableofcontents
\clearpage
\phantomsection
\addcontentsline{toc}{part}{Part I. Core mathematical results}
\section*{Part I. Core mathematical results}
\label{part:core}

\section{Finite Hodge--Dirac reconstruction}
\label{sec:finite}
% =============================================================================

We first construct a finite Hodge--Dirac triple from a real bimodule
derivation and identify the ambiguity of a Dirac implementer with the same
commutators.  These operator-algebraic facts provide the framework for the graph
metric and stationary-current theorems.  The final subsection distinguishes
strong spectral compression from recovery of multiplication in operator norm.

\subsection{Finite real differential data}
\label{subsec:algebraic-input}

The finite construction requires only a finite-dimensional unital $C^*$-algebra
$\A$, a faithful tracial state $\tau$, and the differential datum specified
below.  The Hilbert space of zero-forms is the vector space $\A$ with
inner product
\begin{equation}
 \Hh_0=L^2(\A,\tau),\qquad
 \langle a,b\rangle_0=\tau(a^*b).
 \label{eq:history-H0}
\end{equation}
Faithfulness makes this inner product nondegenerate.  Left and right
multiplication are bounded commuting actions, with adjoints $L_a^*=L_{a^*}$
and $R_a^*=R_{a^*}$.  The right action is a representation of the opposite
algebra: $R_aR_b=R_{ba}$.

There is a distinguished possible choice of trace, but the results do not
require it.  If $\A\cong\bigoplus_\lambda M_{n_\lambda}(\C)$, its normalized
regular trace is
\begin{equation}
 \widehat\tau_{\rm reg}(a)
 =\frac{\Tr_{\C}(L_a)}{\dim_{\C}\A}
 =\frac{\sum_\lambda n_\lambda\Tr(a_\lambda)}
        {\sum_\lambda n_\lambda^2}.
 \label{eq:regular-trace-selfcontained}
\end{equation}
This is a faithful tracial state invariant under $*$-isomorphisms.  Thus both
$\A$ and $\widehat\tau_{\rm reg}$ have ordinary finite-dimensional algebraic
meanings independent of the motivating history construction.

A Hilbert-space $\A$-bimodule means a Hilbert space $\mathcal K$ with commuting
unital $*$-representations of $\A$ and $\A^{\rm op}$.  We denote these by
$\lambda_{\mathcal K}(a)$ and $\rho_{\mathcal K}(b)$, so that
$a\xi b=\lambda_{\mathcal K}(a)\rho_{\mathcal K}(b)\xi$ and
$\rho_{\mathcal K}(a)\rho_{\mathcal K}(b)=\rho_{\mathcal K}(ba)$.

\begin{definition}[Finite real differential datum]
\label{def:finite-real-differential}
A finite real differential datum on $(\A,\tau)$ consists of a Hilbert-space
$\A$-bimodule $\mathcal K$, an antiunitary involution $C_{\mathcal K}$ with
\[
 C_{\mathcal K}(a\xi b)=b^*C_{\mathcal K}(\xi)a^*,
\]
and a complex-linear map $\delta:\A\to\mathcal K$ satisfying
\begin{equation}
 \delta(\one)=0,\qquad
 \delta(ab)=a\delta(b)+\delta(a)b,\qquad
 C_{\mathcal K}\delta(a)=\delta(a^*).
 \label{eq:predictive-derivation}
\end{equation}
Only the bimodule generated by $\delta(\A)$ is retained.  No assumption is
made about the origin of these algebraic data.
\end{definition}

For completeness, the one-form construction can be performed entirely within
this datum.  The universal first-order calculus is
\begin{equation}
 \Om_u^1(\A)=\Ker(m:\A\otimes\A\to\A),\qquad
 \dd_u a=\one\otimes a-a\otimes\one,
 \label{eq:universal-oneforms}
\end{equation}
with bimodule action $c(a\otimes b)d=ca\otimes bd$ and conjugate-linear reversal
\begin{equation}
 \left(\sum_i a_i\otimes b_i\right)^\natural
 =-\sum_i b_i^*\otimes a_i^*.
 \label{eq:oneform-reversal}
\end{equation}
Direct expansion gives $(\dd_u a)^\natural=\dd_u(a^*)$ and
$\dd_u(ab)=a\dd_u b+(\dd_u a)b$.
The map
\[
 T_\delta\left(\sum_i a_i\otimes b_i\right)
 =\sum_i a_i\delta(b_i),\qquad \sum_i a_ib_i=0,
\]
is the unique bimodule map with $T_\delta\dd_u=\delta$.  Indeed, the
Leibniz rule makes it right linear because the additional term is
$(\sum_i a_ib_i)\delta(c)=0$; left linearity is immediate.
Star compatibility gives
$T_\delta(\omega^\natural)=C_{\mathcal K}T_\delta(\omega)$.
Consequently
\begin{equation}
 Q_\delta(\omega,\eta)
 =\langle T_\delta\omega,T_\delta\eta\rangle_{\mathcal K}
 \label{eq:differential-Gram}
\end{equation}
is positive, and its null space is exactly $\Ker T_\delta$.  The resulting
one-form Hilbert space is
\begin{equation}
 \Hh_1(\delta)=\Om_u^1(\A)/\Ker T_\delta
 \cong\Ran T_\delta
 =\Span\{a\delta(b)c:a,b,c\in\A\}.
 \label{eq:H1-delta}
\end{equation}
It is finite dimensional even if the initially supplied $\mathcal K$ is not.
The displayed identification is isometric, so this quotient removes exactly
the unused bimodule directions.  We call the resulting one-form bimodule
\emph{minimal} in this generated-bimodule sense.  Write
$\partial a=[\dd_u a]$ for the induced derivation.

\subsection{The associated spectral triple}

Set $\Hh_0=L^2(\A,\tau)$ and let $L_a$ and $R_a$ denote left and right
multiplication on $\Hh_0$.  On
\begin{equation}
 \Hh=\Hh_0\oplus\Hh_1(\delta)
 \label{eq:H-total}
\end{equation}
define
\begin{equation}
 \pi(a)=L_a\oplus\lambda(a),
 \qquad
 D=\begin{pmatrix}0&\partial^*\\ \partial&0\end{pmatrix},
 \qquad
 \gamma=\begin{pmatrix}I&0\\0&-I\end{pmatrix}.
 \label{eq:finite-D}
\end{equation}
On $\Hh_0$ put $J_0a=a^*$, and on $\Hh_1(\delta)$ let
$J_1[\omega]=[\omega^\natural]$; set $J=J_0\oplus J_1$.

\begin{theorem}[Finite Hodge--Dirac spectral triple]
\label{thm:finite-Hodge-Dirac}
The quintuple $(\A,\Hh,D,J,\gamma)$ is a real even finite spectral triple.  With the
standard sign convention its real-structure signs are those of KO-dimension $0$:
\begin{equation}
 J^2=I,\qquad JD=DJ,\qquad J\gamma=\gamma J.
 \label{eq:KO-zero-signs}
\end{equation}
More precisely, $\pi$ is faithful, $D=D^*$ has compact resolvent, every
$[D,\pi(a)]$ is bounded,
\begin{equation}
 \gamma^*=\gamma,\qquad \gamma^2=I,\qquad
 [\gamma,\pi(a)]=0,\qquad D\gamma=-\gamma D,
 \label{eq:even-grading-axioms}
\end{equation}
and $J$ is antiunitary.  The opposite representation is
\[
 \pi^o(b)=J\pi(b^*)J^{-1}=R_b\oplus\rho(b).
\]
For all $a,b\in\A$ the order-zero and first-order conditions hold:
\begin{equation}
 [\pi(a),\pi^o(b)]=0,
 \qquad
 [[D,\pi(a)],\pi^o(b)]=0.
 \label{eq:order-zero-first}
\end{equation}
If $B_a:\Hh_0\to\Hh_1(\delta)$ is defined by
$B_ab=(\partial a)b$, then
\begin{equation}
 [D,\pi(a)]
 =\begin{pmatrix}0&-B_{a^*}^*\\B_a&0\end{pmatrix}.
 \label{eq:finite-commutator}
\end{equation}
In particular, for $a=a^*$,
\begin{equation}
 \|[D,\pi(a)]\|=\|B_a\|,
 \qquad
 \Ker\bigl(a\mapsto\|[D,\pi(a)]\|\bigr)=\Ker\partial\cap\A_{\rm sa}.
 \label{eq:finite-Lip}
\end{equation}
\end{theorem}

\begin{proof}
Faithfulness follows from the left regular summand.  Since all spaces are finite
dimensional, $D$ is self-adjoint with compact resolvent and all commutators are
bounded.  The grading identities in \eqref{eq:even-grading-axioms} follow directly
from the block-diagonal representation and the off-diagonal form of $D$.
The Leibniz rule gives
\[
 (\partial L_a-\lambda(a)\partial)b
 =\partial(ab)-a\partial b=(\partial a)b=B_ab,
\]
and taking adjoints yields \eqref{eq:finite-commutator}.  The reversal identities
make $J_1$ antiunitary, give $J_1^2=I$, and imply
$J_1\partial=\partial J_0$; hence the three signs in
\eqref{eq:KO-zero-signs}.  Moreover $J_0L_{b^*}J_0=R_b$ and
$J_1\lambda(b^*)J_1=\rho(b)$, so the opposite representation has the asserted
form.  Left and right actions commute.  Finally, the lower-left block of the
double commutator is $B_aR_b-\rho(b)B_a$, which vanishes because $B_a$ is right
$\A$-linear.  The upper-right block vanishes by taking the corresponding
adjoint identity for $a^*$ and $b^*$.  For self-adjoint $a$, the two
off-diagonal blocks are $B_a$ and $-B_a^*$, so both have norm $\|B_a\|$,
proving \eqref{eq:finite-Lip}.
\end{proof}

\begin{remark}
Here ``real even'' refers to the grading, reality, order-zero and first-order
axioms listed in Theorem~\ref{thm:finite-Hodge-Dirac}; orientability and
Poincar\'e duality are not asserted.  The real structure is the algebraic reversal
of the differential bimodule.  Additional physical interpretations of $J$ require
extra markings and are not used below.
\end{remark}

\begin{proposition}
\label{prop:no-scale}
Let $\delta\neq0$ be a finite real differential datum and let $s>0$.  Replacing
$\delta$ by $s\delta$ leaves $(\A,\tau)$ unchanged and, under the canonical
identification of the one-form spaces, sends
\begin{equation}
 D\longmapsto sD,
 \qquad
 L_D(a):=\|[D,\pi(a)]\|\longmapsto sL_D(a).
 \label{eq:scale-D}
\end{equation}
Whenever the Connes distance is finite, it scales as
$d_D\mapsto s^{-1}d_D$.  Hence the algebra and canonical trace do not determine an
absolute spectral metric scale.
\end{proposition}

\begin{proof}
The inner product induced by $s\delta$ is $s^2Q_\delta$, so the canonical unitary
from the scaled one-form space to the original one sends the derivation vector
$[\dd_u a]$ to $s[\dd_u a]$.  Both off-diagonal blocks of $D$ therefore acquire the
factor $s$.  The commutator and dual-distance scalings follow immediately.
\end{proof}

The spectral seminorm also gives a compact quantum metric space, provided that
zero differential means constant observable.  This is the finite-dimensional
form of the state-space criterion underlying Rieffel's theory
\cite{Rieffel1998,Rieffel1999}.

\begin{proposition}
\label{prop:finite-CQMS}
Suppose $\Ker\partial=\C\one$ and set
$L_D(a)=\|[D,\pi(a)]\|$.  On the real order-unit space $\A_{\rm sa}$,
$L_D$ is a continuous Lip-norm.  Explicitly, the distance
\begin{equation}
 \operatorname{mk}_{L_D}(\varphi,\psi)
 =\sup\{|\varphi(a)-\psi(a)|:
          a=a^*,\ L_D(a)\le1\}
 \label{eq:finite-state-metric}
\end{equation}
is finite on the state space $S(\A)$ and induces its weak-$*$ topology.
Moreover, the complex extension satisfies
\begin{equation}
 L_D(ab)\le\|a\|L_D(b)+\|b\|L_D(a),\qquad
 L_D(a^*)=L_D(a).
 \label{eq:Leibniz-Lip}
\end{equation}
Conversely, if $\Ker\partial\ne\C\one$, the distance in
\eqref{eq:finite-state-metric} is infinite for some pair of states, so it is
not a compact quantum metric on the whole state space.
\end{proposition}

\begin{proof}
The commutator is a linear map between finite-dimensional normed spaces;
hence $L_D$ is continuous.  The commutator product rule and the $*$-identity
give \eqref{eq:Leibniz-Lip}.  By Theorem~\ref{thm:finite-Hodge-Dirac},
the kernel on $\A_{\rm sa}$ is $\R\one$.
On $E_\tau=\{a=a^*:\tau(a)=0\}$ it follows that $L_D$ is a norm, and
finite-dimensional norm equivalence gives a constant $C_\delta$ such that
\begin{equation}
 \|a-\tau(a)\one\|\le C_\delta L_D(a),\qquad a=a^*.
 \label{eq:finite-Poincare-Lip}
\end{equation}
Thus $\{a\in E_\tau:L_D(a)\le1\}$ is norm compact.  Scalar multiples of
$\one$ do not affect differences of states, so the supremum in
\eqref{eq:finite-state-metric} can be taken over this compact set and is at
most $2C_\delta$.  Weak-$*$ convergence of states is uniform on a norm-compact
set, since states have norm one; it therefore implies convergence in
$\operatorname{mk}_{L_D}$.  Conversely, for each self-adjoint $a$,
\[
 |\varphi(a)-\psi(a)|\le L_D(a)
       \operatorname{mk}_{L_D}(\varphi,\psi),
\]
with the scalar case immediate.  Hence metric convergence implies weak-$*$
convergence and the metric separates states.

If the kernel is larger, star compatibility gives a non-scalar self-adjoint
$b\in\Ker\partial$.  States separate self-adjoint elements modulo scalars,
so some $\varphi,\psi$ have $\varphi(b)\ne\psi(b)$.  Every multiple $tb$
has zero seminorm, making their distance infinite.
\end{proof}

The quotient used here is $\A_{\rm sa}/\R\one$; the metric itself is on
$S(\A)$, not on a quotient of states.  Finite dimensionality is essential to
the automatic compactness argument: for an infinite-dimensional limit, a scalar
kernel alone does not replace total boundedness of the normalized Lipschitz ball.


\begin{proposition}
\label{prop:hodge-zero-mode}
For every differential datum of Definition~\ref{def:finite-real-differential}, the
Hodge operator of Theorem~\ref{thm:finite-Hodge-Dirac} satisfies
\begin{equation}
 D(\one,0)=0.
 \label{eq:hodge-zero-mode}
\end{equation}
Consequently the Hodge functor $\mathfrak H_{\rm fin}$ never produces an
invertible finite Dirac operator.  More generally, knowledge of the commutator
derivation does not determine the component of $D$ in the compatible commutant;
that residual ambiguity is characterized in Theorem~\ref{thm:commutant-fibre}.
\end{proposition}

\begin{proof}
The universal derivation satisfies $\partial\one=0$, hence the lower-left block of
$D(\one,0)$ vanishes; the upper-right block acts on the zero one-form component and
vanishes as well.  The second assertion is immediate from
Theorem~\ref{thm:commutant-fibre}.
\end{proof}

\subsection{The commutant fibre of a fixed derivation}

\begin{definition}[Marked finite spectral triple and commutant coordinates]
\label{def:marked-spectral-triple}
A marked finite spectral triple is
\[
 \mathsf S=(\A,\Hh,\pi,D,J,\gamma;\mathfrak m),
\]
where $\mathfrak m$ denotes any additional declared data preserved by
isomorphisms.  When a real structure is present, its Dirac sign is written
$JD=\varepsilon'DJ$, $\varepsilon'\in\{\pm1\}$; absent marks impose no
condition.  Define the compatible real commutant space
\begin{equation}
 \mathfrak c_{\mathsf S}
 =\{C\in\pi(\A)'_{\rm sa}:
 C\gamma+\gamma C=0,\quad JC=\varepsilon'CJ\},
 \label{eq:commutant-space}
\end{equation}
with the corresponding condition omitted when $\gamma$ or $J$ is absent.  An $m$-component \emph{real linear commutant-coordinate map} means
\[
 \alpha:\mathfrak c_{\mathsf S}\longrightarrow\R^m.
\]
Its components are the scalar commutant coordinates used below.
Let $P_{\mathfrak c_{\mathsf S}}$ be real Hilbert--Schmidt orthogonal projection
onto this space and put
\begin{equation}
 D_{\rm der}=D-P_{\mathfrak c_{\mathsf S}}D,
 \qquad
 q_{\mathsf S}=\dim_{\R}\mathfrak c_{\mathsf S}.
 \label{eq:centered-D}
\end{equation}
\end{definition}

For a fixed represented algebra and fixed grading and reality signs, the elementary
commutator identity identifies the remaining ambiguity exactly.  The following
statement records its affine structure and the sharp number of independent
real linear coordinates needed to recover the operator.

\begin{theorem}[Dirac implementers with a fixed commutator derivation]
\label{thm:commutant-fibre}
Among self-adjoint operators with the declared grading and reality signs that
implement the same commutator derivation $[D,\pi(\cdot)]$, the complete solution
set is
\begin{equation}
 D+\mathfrak c_{\mathsf S}.
 \label{eq:Dirac-affine}
\end{equation}
The centered operator $D_{\rm der}$ is the unique minimum-Hilbert--Schmidt-norm
representative of this affine fibre.  If
$C_1,\ldots,C_{q_{\mathsf S}}$ is a real Hilbert--Schmidt orthonormal basis of
$\mathfrak c_{\mathsf S}$ and
\begin{equation}
 \alpha_j(D)=\langle C_j,D\rangle_{\HS,\R},
 \qquad 1\le j\le q_{\mathsf S},
 \label{eq:commutant-coordinates}
\end{equation}
then
\begin{equation}
 D=D_{\rm der}
   +\sum_{j=1}^{q_{\mathsf S}}\alpha_j(D)C_j.
 \label{eq:commutant-reconstruction}
\end{equation}
Consequently $q_{\mathsf S}$ independent real linear commutant coordinates are sufficient to
recover every compatible implementer from its commutator derivation, and fewer
coordinates cannot distinguish the full affine fibre in general.  In particular, the
canonical derivation-only reconstruction is exactly the centered class
$P_{\mathfrak c_{\mathsf S}}D=0$.
\end{theorem}

\begin{proof}
If $D'$ has the same commutator derivation as $D$, then $C=D'-D$ commutes with
$\pi(\A)$.  Subtracting the grading and reality equations for $D'$ and $D$ gives
$C\in\mathfrak c_{\mathsf S}$.  Conversely every
$C\in\mathfrak c_{\mathsf S}$ preserves the commutators and all declared linear
sign relations, proving \eqref{eq:Dirac-affine}.  Orthogonal projection gives
\[
 \|D+C\|_{\HS}^2
 =\|D_{\rm der}\|_{\HS}^2
  +\|P_{\mathfrak c_{\mathsf S}}D+C\|_{\HS}^2,
\]
so $D_{\rm der}$ is the unique minimum-norm representative.  Expanding
$P_{\mathfrak c_{\mathsf S}}D$ in the orthonormal basis gives
\eqref{eq:commutant-reconstruction}.  By Definition~\ref{def:marked-spectral-triple}, any $m$-component real linear commutant-coordinate map is a
real linear map $\mathfrak c_{\mathsf S}\to\R^m$.  If
$m<q_{\mathsf S}$ the map has nonzero kernel, so two distinct compatible
implementers remain indistinguishable.  Taking the coordinate functionals in
\eqref{eq:commutant-coordinates} gives an injective commutant-coordinate map of size $q_{\mathsf S}$.
\end{proof}

\subsection{Functoriality and finite spectral compressions}

Let $\mathbf{Diff}_{\rm fin}$ have objects $(\A,\tau,\delta)$ as above and
morphisms the trace-preserving $*$-isomorphisms intertwining the differential
bimodule data.  The construction of Theorem~\ref{thm:finite-Hodge-Dirac} is
functorial: the algebra isomorphism induces unitaries on $L^2(\A,\tau)$ and on the
minimal one-form quotient, and these unitaries intertwine $\pi,D,J,\gamma$.
We denote this finite Hodge functor by $\mathfrak H_{\rm fin}$.

Reconstruction from commutator data retains the represented algebra, the
specified real structure and grading, and separating real linear coordinates
on the compatible commutant.  The Dirac operator is then recovered from these
data.  In the notation below, the superscript $\mathrm{cal}$ records this
choice of commutant coordinates.  For compact-resolvent triples, finite data
are obtained by spectral compression, with two distinct convergence
requirements.

\begin{definition}[Finite reconstruction and spectral-compression data]
\label{def:reconstruction-categories}
Let $\mathbf{Spec}_{\rm fin}$ be the category of finite marked spectral triples,
with unitary equivalences intertwining the represented algebra, $D$, and all
declared marks as isomorphisms.  Let $\mathbf{Spec}_{\rm cpt}^{\rm sep}$ denote
the corresponding separable compact-resolvent triples.  For every compact-resolvent triple
we fix, as part of the marked datum, a dense unital $*$-subalgebra
$\A_D\subseteq\A$ on which $[D,\pi(a)]$ is bounded.  In the finite case
$\A_D=\A$.

The reconstruction problem is considered in three source categories.
\begin{enumerate}[label=\textup{(C\arabic*)},leftmargin=3em]
\item In the \emph{finite reconstruction category}, the represented finite algebra and
      symmetry marks are retained together with the commutator derivation and
      separating real linear commutant coordinates on $\mathfrak c_{\mathsf S}$.
      The resulting reconstruction is denoted $\mathfrak S_{\rm fin}^{\rm cal}$.
\item In the \emph{strong spectral-compression category}, one uses increasing
      finite-rank spectral projections $P_n\uparrow I$ with $P_nD=DP_n$ and records
      the compressed matrices, product defects, commutators, marks, and finite
      commutant coordinates.  After extending compressed operators by zero on
      $P_n^\perp\Hh$, convergence means that for every $a,b\in\A_D$,
      $z\in\mathbb C\setminus\mathbb R$, and $\xi\in\Hh$,
      \begin{align}
      P_n\pi(a)P_n\xi&\longrightarrow\pi(a)\xi,\nonumber\\
      P_n[D,\pi(a)]P_n\xi&\longrightarrow[D,\pi(a)]\xi,\nonumber\\
      P_n(D-z)^{-1}P_n\xi&\longrightarrow(D-z)^{-1}\xi,\nonumber\\
      \bigl(P_n\pi(ab)P_n-P_n\pi(a)P_n\pi(b)P_n\bigr)\xi&\longrightarrow0.
      \label{eq:strong-compression-topology}
      \end{align}
      The declared bounded marks are required to converge in the same strong sense.
      The reconstruction is denoted $\mathfrak S_{\rm cpt}^{\rm strong,cal}$.
\item In the \emph{norm-multiplicative spectral-compression category}, the same
      data are required in addition to satisfy
      \begin{equation}
      \|P_n\pi(ab)P_n-P_n\pi(a)P_n\pi(b)P_n\|\longrightarrow0
      \qquad(a,b\in\A_D).
      \label{eq:norm-multiplicative-compression}
      \end{equation}
      The corresponding reconstruction is $\mathfrak S_{\rm cpt}^{\rm norm,cal}$.
\end{enumerate}
A compact-resolvent triple is called \emph{spectrally quasidiagonal} when there
are finite-rank symmetric spectral projections $P_n\uparrow I$ such that
\begin{equation}
 P_nD=DP_n,
 \qquad
 \|[P_n,\pi(a)]\|\longrightarrow0
 \label{eq:sqd}
\end{equation}
for every $a\in\A_D$, equivariantly for $J$ and $\gamma$ when these are marked.
Write $\mathbf{Spec}_{\rm cpt}^{\rm sqd}$ for this subcategory.  The source classes
are taken modulo unitary equivalences preserving the represented algebra,
$\A_D$, marks, commutant-coordinate functionals, and compression maps.
\end{definition}

Strong compression uses strong convergence of bounded operator data and
exists for every compact-resolvent triple.  Norm-multiplicative compression
also recovers multiplication in operator norm.  Skalski and Zacharias
\cite{SkalskiZacharias2009} study the relation between compact-type spectral
triples and quasidiagonal representations, including a Toeplitz example.
Our condition is more restrictive than quasidiagonality of the representation:
the approximating projections must also be spectral projections of the
prescribed $D$.  The next theorem identifies this additional requirement
precisely through the norm of the multiplication defect.

\begin{proposition}[Finite recovery from commutant coordinates]
\label{prop:finite-calibrated-recovery}
For a finite marked spectral triple $\mathsf S$, its represented algebra, declared
symmetry signs and commutator derivation determine the centered implementer
$D_{\rm der}$.  Exactly
$q_{\mathsf S}=\dim_{\mathbb R}\mathfrak c_{\mathsf S}$ independent real linear
commutant coordinates suffice, and are in general necessary, to recover the full
compatible Dirac operator.
\end{proposition}

\begin{proof}
Let
$\mathsf S_0=(\mathcal B,K,\pi_0,D_0,J_0,\gamma_0)$ be finite.  Retain the
represented algebra $\pi_0(\mathcal B)$, the declared marks, and the commutator
derivation
\[
 \delta_{D_0}(a)=[D_0,\pi_0(a)].
\]
Choose a real Hilbert--Schmidt orthonormal basis
$C_1,\ldots,C_{q_{\mathsf S_0}}$ of the compatible commutant and retain the
coordinates $\alpha_j(D_0)$ of \eqref{eq:commutant-coordinates}.  By
Theorem~\ref{thm:commutant-fibre}, the commutator data determine the centered
implementer $(D_0)_{\rm der}$ and these coordinates recover $D_0$ uniquely by
\eqref{eq:commutant-reconstruction}.  Thus the full generator need not be retained,
and the same theorem proves sharpness of the number of real linear coordinates.
This is a finite-dimensional linear reconstruction statement.
\end{proof}

The compression problem is different: finite-rank compressions of a representation
are generally operator systems, not representations of the original algebra.
In particular, $a\mapsto P_n\pi(a)P_n$ need not be multiplicative, and the compressed
matrices are not asserted to inherit every real first-order axiom.  This is the
operator-system distinction relevant to spectral truncations
\cite{vanSuijlekom2021,Leimbach2025}.  The next theorem concerns the matrix data and
topologies explicitly specified in Definition~\ref{def:reconstruction-categories}.

\begin{theorem}[Strong and norm-multiplicative spectral compressions]
\label{thm:spectral-compressions}
Every separable compact-resolvent spectral triple admits finite-rank symmetric
spectral projections $P_n\uparrow I$ satisfying the strong convergence conditions
\eqref{eq:strong-compression-topology}.  For
$D_n=P_nD|_{P_n\Hh}$ and $a_n=P_n\pi(a)|_{P_n\Hh}$ one has
\begin{equation}
 [D_n,a_n]=P_n[D,\pi(a)]|_{P_n\Hh},\qquad a\in\A_D.
 \label{eq:compressed-commutator-identity}
\end{equation}
For a sequence of such projections, norm-multiplicativity
\eqref{eq:norm-multiplicative-compression} holds if and only if
$\|[P_n,\pi(a)]\|\to0$ for every $a\in\A_D$.  Consequently a
norm-multiplicative spectral-compression sequence exists exactly for spectrally
quasidiagonal triples.  This class is a proper subclass of the separable
compact-resolvent triples.
\end{theorem}

\begin{proof}
Let $(\A,\Hh,\pi,D)$ be separable with compact resolvent and choose increasing
symmetric spectral projections
$P_n=\mathbf 1_{[-R_n,R_n]}(D)$, $R_n\uparrow\infty$.  Then
$P_nD=DP_n$ and $P_n\to I$ strongly.  The same commutation relation gives
\eqref{eq:compressed-commutator-identity} by expanding the two compressed products.
Symmetric cutoffs preserve the declared $J$ and $\gamma$ signs, and their bounded
compressions converge strongly as well.  For $a,b\in\A_D$,
\begin{equation}
 P_n\pi(ab)P_n-P_n\pi(a)P_n\pi(b)P_n
 =P_n\pi(a)(I-P_n)\pi(b)P_n,
 \label{eq:compression-defect}
\end{equation}
which tends strongly to zero because $(I-P_n)\to0$ strongly and the remaining
factors are uniformly bounded.  Since $[D,\pi(a)]$ is bounded,
$P_n[D,\pi(a)]P_n\to[D,\pi(a)]$ strongly.  Likewise
$P_n(D-z)^{-1}P_n\to(D-z)^{-1}$ strongly for $z\notin\mathbb R$.
Thus \eqref{eq:strong-compression-topology} holds for the canonical spectral
compressions, proving the strong-compression assertion.

If \eqref{eq:sqd} holds, then from \eqref{eq:compression-defect}
\[
 \|P_n\pi(a)(I-P_n)\pi(b)P_n\|
 \le \|\pi(a)\|\,\|[P_n,\pi(b)]\|,
\]
so multiplication is recovered in norm.  Conversely, applying
\eqref{eq:norm-multiplicative-compression} to $(a^*,a)$ gives
\[
 \|P_n\pi(a)^*(I-P_n)\pi(a)P_n\|
 =\|(I-P_n)\pi(a)P_n\|^2\longrightarrow0.
\]
Applying the same argument to $a^*$ controls the opposite off-diagonal block;
hence $\|[P_n,\pi(a)]\|\to0$.  This proves the equivalence with spectral quasidiagonality.

Properness is witnessed by the Toeplitz representation on $\ell^2(\N_0)$ with
number operator $Ne_k=ke_k$, unilateral shift $Se_k=e_{k+1}$, and dense unital
$*$-algebra $\A_D=\operatorname{alg}(I,S,S^*)$.  Here
$[N,S]=S$ and $[N,S^*]=-S^*$, so the commutators are bounded.  Every finite-rank
spectral projection of $N$ is diagonal in the number basis, and every nontrivial
finite diagonal projection has $\|[P,S]\|=1$.  This compact-resolvent triple lies
in the strong-compression class but not in the norm-multiplicative class.
\end{proof}

\begin{corollary}[Reconstruction classes and their essential images]
\label{thm:essential-image}
For the reconstruction classes in Definition~\ref{def:reconstruction-categories},
\begin{equation}
 \begin{aligned}
 \EssIm(\mathfrak S_{\rm fin}^{\rm cal})
   &=\mathbf{Spec}_{\rm fin},\\
 \EssIm(\mathfrak S_{\rm cpt}^{\rm strong,cal})
   &=\mathbf{Spec}_{\rm cpt}^{\rm sep},\\
 \EssIm(\mathfrak S_{\rm cpt}^{\rm norm,cal})
   &=\mathbf{Spec}_{\rm cpt}^{\rm sqd}
   \subsetneq\mathbf{Spec}_{\rm cpt}^{\rm sep}.
 \end{aligned}
 \label{eq:essential-image}
\end{equation}
\end{corollary}

\begin{proof}
The finite equality is Proposition~\ref{prop:finite-calibrated-recovery}.
The compact-resolvent equalities are Theorem~\ref{thm:spectral-compressions}, with
the finite commutant coordinates retained as in the source definitions.
\end{proof}

\begin{remark}
The essential images in Corollary~\ref{thm:essential-image} refer to the
specified reconstruction data.  Finite recovery uses the commutator derivation
and separating commutant coordinates.  The Hodge construction alone is more
restrictive: it always has the unit zero mode of
Proposition~\ref{prop:hodge-zero-mode}.  For spectral compressions the
additional distinction is whether multiplication is recovered strongly or
in operator norm.
\end{remark}


% =============================================================================

\section{Metric realization on stationary graphs}
\label{sec:metric}
% =============================================================================

The finite reconstruction results identify compatible operator data.  To test their
metric content, we now compare the Connes distance with a metric obtained without
using the Dirac operator.  Stationary graph dynamics provides a transparent model:
its symmetric jump covariance fixes the candidate metric normalization, while its
Hodge operator supplies the commutator seminorm.  The graph construction belongs
to the established Hodge--Dirac framework \cite{Davies1993,Requardt2002}; the
continuum comparison is the question addressed here.

\subsection{Stationary graphs and the Hodge--Dirac operator}

\begin{definition}[Stationary finite Markov graph]
\label{def:supp-renewal-graph}
Let $V$ be finite, let $m_x>0$ with $\sum_xm_x=1$, and let $k_{xy}\ge0$ for
$x\ne y$.  Put
\begin{equation}
 (\mathcal Lf)(x)=\sum_{y\ne x}k_{xy}(f(y)-f(x)),
 \qquad q_{xy}=m_xk_{xy},
 \label{eq:supp-renewal-generator}
\end{equation}
and assume stationarity
\begin{equation}
 \sum_y(q_{xy}-q_{yx})=0
 \qquad(x\in V).
 \label{eq:supp-stationarity}
\end{equation}
Define
\begin{equation}
 c_{xy}=\frac{q_{xy}+q_{yx}}2=c_{yx},
 \qquad
 j_{xy}=\frac{q_{xy}-q_{yx}}2=-j_{yx}.
 \label{eq:supp-conductance-current}
\end{equation}
Assume the undirected graph on the pairs with $c_{xy}>0$ is connected.
\end{definition}

Set
\begin{equation}
 \Hh_0=\ell^2(V,m),
 \qquad
 \Hh_1=\ell^2(E^{\rm or},c),
 \qquad
 (\partial f)(x,y)=f(y)-f(x),
 \label{eq:supp-graph-carriers}
\end{equation}
where both orientations are retained and
$\norm{\omega}_1^2=\sum_{x,y}c_{xy}|\omega(x,y)|^2$.  Let
\begin{equation}
 (\lambda(a)\omega)(x,y)=a(y)\omega(x,y),
 \qquad
 (\rho(a)\omega)(x,y)=a(x)\omega(x,y).
 \label{eq:supp-graph-actions}
\end{equation}
Then $\partial(fg)=\lambda(f)\partial g+\rho(g)\partial f$.  The associated
Hodge operator is
\begin{equation}
 D_{m,c}=\begin{pmatrix}0&\partial^*\\\partial&0\end{pmatrix}.
 \label{eq:supp-graph-Dirac}
\end{equation}
The represented algebra and real-even structure are
\begin{equation}
 \begin{gathered}
 \pi(a)=M_a\oplus\lambda(a),\qquad
 \gamma=\begin{pmatrix}I&0\\0&-I\end{pmatrix},\qquad J=J_0\oplus J_1,\\
 (J_0f)(x)=\overline{f(x)},\qquad
 (J_1\omega)(x,y)=-\overline{\omega(y,x)}.
 \end{gathered}
 \label{eq:graph-real-even-structure}
\end{equation}
The minus sign is required by the derivation convention:
$J_1\partial f=\partial\overline f$.  Symmetry of $c$ makes $J_1$ antiunitary,
and $J_1\lambda(a^*)J_1=\rho(a)$.  Thus these are precisely the KO-dimension-zero
signs and opposite actions of Theorem~\ref{thm:finite-Hodge-Dirac}.

\begin{theorem}[Graph Hodge square and exact Lipschitz seminorm]
\label{thm:supp-graph-square}
Let $\mathcal L^\dagger$ be the adjoint of $\mathcal L$ in $\ell^2(V,m)$ and
let $\mathcal L_s=(\mathcal L+\mathcal L^\dagger)/2$.  Then
\begin{equation}
 {D_{m,c}^2|_{\Hh_0}=\partial^*\partial=-2\mathcal L_s.}
 \label{eq:supp-graph-square}
\end{equation}
For every real $f\in C(V)$,
\begin{equation}
 {
 L_{m,c}(f)^2
 =\norm{[D_{m,c},f]}^2
 =\max_{x\in V}\frac1{m_x}
   \sum_yc_{xy}|f(y)-f(x)|^2.}
 \label{eq:supp-graph-Lip}
\end{equation}
Thus, with the standard symmetric carr\'e-du-champ normalization
\begin{equation}
 \Gamma_s(f)(x)=\frac1{2m_x}\sum_yc_{xy}|f(y)-f(x)|^2,
 \label{eq:supp-carre-du-champ}
\end{equation}
one has $L_{m,c}(f)^2=2\|\Gamma_s(f)\|_\infty$.  The kernel of $L_{m,c}$ is the constants, so the Connes formula gives a finite
metric on $V$.  On the reversible branch, $\mathcal L=\mathcal L_s$ and
$D_{m,c}^2|_{\Hh_0}=-2\mathcal L$.
\end{theorem}

\begin{proof}
The adjoint generator has rates $q_{yx}/m_x$, so $\mathcal L_s$ has rates
$c_{xy}/m_x$.  Therefore
\[
 -\ip{f}{\mathcal L_sf}_0
 =\frac12\sum_{x,y}c_{xy}|f(y)-f(x)|^2
 =\frac12\norm{\partial f}_1^2.
\]
Polarization gives \eqref{eq:supp-graph-square}.  The lower-left commutator
block satisfies
\[
 (B_fg)(x,y)=(f(y)-f(x))g(x).
\]
Hence
\[
 \norm{B_fg}_1^2
 =\sum_x|g(x)|^2\sum_yc_{xy}|f(y)-f(x)|^2.
\]
Its norm relative to $\sum_xm_x|g(x)|^2$ is the maximum in
\eqref{eq:supp-graph-Lip}.  Connectedness makes $\partial f=0$ equivalent to
constancy.
\end{proof}

The connectedness assumption has a state-space consequence as well as a
pointwise one.  Proposition~\ref{prop:finite-CQMS} implies that
$(C(V)_{\rm sa},L_{m,c})$ is a compact quantum metric space.  The exact normalization
\eqref{eq:supp-graph-Lip}--\eqref{eq:supp-carre-du-champ} relates the squared
commutator seminorm to the maximal local carr\'e-du-champ of the symmetric Markov
generator; the continuum result below is
the uniform convergence of the resulting $A_3$ Connes metric to the independently
reconstructed Riemannian metric.  Its states are
probability distributions on $V$, and the spectral metric on those distributions
induces their usual weak-$*$ topology.  This conclusion does not require the
Hodge operator itself to be invertible; it is the scalar kernel of the
\emph{commutator seminorm} that matters.  The following large one-form kernel is
therefore an issue for the operator continuum limit, not for the finite Lip-norm.

\begin{proposition}
\label{cth:supp-extensive-kernel}
Let the connected underlying graph have $n$ vertices and $e$ undirected edges,
and retain both orientations in $\Hh_1$.  Then
\begin{equation}
 {\dim\Ker D_{m,c}=2e-n+2.}
 \label{eq:supp-extensive-kernel}
\end{equation}
On a periodic $A_3$ root graph, $e=6n$, so the kernel has dimension $11n+2$.
Consequently, any compatible cofinal limit that retains all these zero modes
has infinite-dimensional kernel and cannot have compact resolvent.
\end{proposition}

\begin{proof}
Connectedness gives $\dim\Ker\partial=1$ and $\rank\partial=n-1$.  Since the
directed one-form space has dimension $2e$,
$\dim\Ker\partial^*=2e-n+1$.  Because
$\Ker D=(\Ker\partial)\oplus(\Ker\partial^*)$, the formula follows.
\end{proof}

\subsection{The periodic \texorpdfstring{$A_3$}{A3} lattice and its covariance normalization}

Let
\begin{equation}
 W_0=\left\{x\in\R^4:\sum_{i=1}^4x_i=0\right\},
 \qquad
 \Lambda_{A_3}=W_0\cap\Z^4,
 \label{eq:supp-A3-lattice}
\end{equation}
and let
$\Phi_{A_3}=\{e_i-e_j:i\ne j\}$ be the twelve oriented roots.

\begin{lemma}[$A_3$ tight-frame identity]
\label{lem:supp-A3-frame}
On $W_0$,
\begin{equation}
 {\sum_{\alpha\in\Phi_{A_3}}\alpha\alpha^{\mathsf T}=8I_{W_0}.}
 \label{eq:supp-A3-frame}
\end{equation}
Equivalently,
$\sum_{i\ne j}(v_i-v_j)^2=8\norm v^2$ for $v\in W_0$.
\end{lemma}

\begin{proof}
Expand the left side.  Each $v_i^2$ occurs six times in the two diagonal sums,
while
$\sum_{i\ne j}v_iv_j=(\sum_iv_i)^2-\sum_iv_i^2=-\sum_iv_i^2$.
The total is $8\sum_i v_i^2$.  Polarization gives the operator identity.
\end{proof}

Put
\begin{equation}
 M=W_0/\Lambda_{A_3},
 \qquad
 V_h=h\Lambda_{A_3}/\Lambda_{A_3},
 \qquad h=d^{-1},\quad d\in\N,\ d\ge3,
 \label{eq:supp-A3-torus}
\end{equation}
and assign each root edge the reversible rate and vertex mass
\begin{equation}
 k_{\alpha,h}=\frac1{8h^2},
 \qquad
 m_h(x)=h^3,
 \qquad
 c_{\alpha,h}=\frac h8.
 \label{eq:supp-A3-data}
\end{equation}


We identify $W_0$ with its dual using the ambient Euclidean pairing.  For a
translation-invariant jump process with displacements $v_\alpha$ and rates
$k_\alpha$, the infinitesimal covariance tensor is
$Q=\sum_\alpha k_\alpha v_\alpha\otimes v_\alpha$.  When it is positive definite,
we use $Q$ as the inverse metric tensor, so that the candidate metric is
$g=Q^{-1}$.  This convention corresponds to a diffusion generator with principal
part $\frac12 Q:\nabla^2$ and fixes the factor of two independently of the Hodge
square.

\begin{proposition}[Jump covariance and the independent metric]
\label{prop:A3-covariance}
For the rates and displacements in \eqref{eq:supp-A3-data},
\begin{equation}
 Q_h=\sum_{\alpha\in\Phi_{A_3}}k_{\alpha,h}
                  (h\alpha)\otimes(h\alpha)
     =\frac18\sum_{\alpha\in\Phi_{A_3}}\alpha\otimes\alpha
     =I_{W_0}.
 \label{eq:A3-jump-covariance}
\end{equation}
Thus $g=Q_h^{-1}$ is the ambient flat metric on $W_0/\Lambda_{A_3}$, independent
of $h$.  Its geodesic distance is
\begin{equation}
 d_g(x,y)=\min_{\ell\in\Lambda_{A_3}}
                  \|\widetilde x-\widetilde y+\ell\|,
 \label{eq:A3-independent-distance}
\end{equation}
where $\widetilde x,\widetilde y$ are arbitrary lifts to $W_0$.
\end{proposition}

\begin{proof}
Substitute $k_{\alpha,h}=1/(8h^2)$ and apply
Lemma~\ref{lem:supp-A3-frame}.  The quotient formula for the flat geodesic
distance gives \eqref{eq:A3-independent-distance} and is independent of the lifts.
No Dirac operator or commutator seminorm enters this computation.
\end{proof}

Integrals on $M$ below use normalized Haar measure, equivalently the flat
Riemannian volume divided by its total volume.  Each periodic fundamental mesh
cell then has measure $h^3$, consistently with the vertex masses in
\eqref{eq:supp-A3-data}.  This normalization does not change $g$ or its gradient.

\begin{theorem}[$A_3$ Hodge identities and smooth-function consistency]
\label{thm:supp-A3-finite}
The graph spectral triple has
\begin{equation}
 {
 L_h(f)^2
 =\max_{x\in V_h}\frac1{8h^2}
  \sum_{\alpha\in\Phi_{A_3}}
  |f(x+h\alpha)-f(x)|^2,}
 \label{eq:supp-A3-Lip}
\end{equation}
\begin{equation}
 {D_h^2|_{C(V_h)}=-2\mathcal L_h,}
 \qquad
 (\mathcal L_hf)(x)
 =\frac1{8h^2}\sum_{\alpha\in\Phi_{A_3}}
   (f(x+h\alpha)-f(x)).
 \label{eq:supp-A3-square}
\end{equation}
For every $F\in C^2(M)$,
\begin{equation}
 \left|L_h(F|_{V_h})-\norm{\nabla F}_{L^\infty(M)}\right|
 \le C_{A_3}h\norm F_{C^2(M)}.
 \label{eq:supp-A3-Lip-convergence}
\end{equation}
Moreover $\mathcal L_hF\to\frac12\Delta F$ uniformly for $F\in C^4(M)$.
\end{theorem}

\begin{proof}
The first two identities follow from \Cref{thm:supp-graph-square} and
\eqref{eq:supp-A3-data}.  Taylor expansion gives
\[
 \frac{F(x+h\alpha)-F(x)}h
 =\alpha\cdot\nabla F(x)+r_{\alpha,h}(x),
 \qquad |r_{\alpha,h}(x)|\le Ch\norm F_{C^2}.
\]
Apply the reverse triangle inequality in $\ell^2(\Phi_{A_3})$ and
\Cref{lem:supp-A3-frame}, then maximize over the mesh.  Since the mesh is a
$C h$-net and $\nabla F$ has Lipschitz constant bounded by $\|F\|_{C^2}$, the
maximum of $|\nabla F|$ on the mesh differs from its supremum on $M$ by at most
$C h\|F\|_{C^2}$.  This gives \eqref{eq:supp-A3-Lip-convergence}.  For the generator, the linear terms cancel
by central symmetry.  The quadratic term is
$\frac1{16}\sum_\alpha(\alpha\otimes\alpha):\nabla^2F=\frac12\Delta F$;
the cubic terms cancel and the remainder is $O(h^2\norm F_{C^4})$.
\end{proof}

\subsection{Uniform convergence of the Connes distance}
\label{subsec:A3-distance}

Consistency on smooth functions does not by itself control the supremum in
the Connes distance formula.  We first construct a fixed family of periodic
interpolants and then prove compactness of the normalized discrete Lipschitz
balls.  The interpolation estimate gives a uniform Lipschitz bound; the
quadratic edge constraint will improve the limiting bound to the sharp
constant one.

\begin{lemma}[A periodic shape-regular $A_3$ triangulation]
\label{lem:A3-periodic-triangulation}
Let $b_i=e_i-e_4$ for $1\le i\le3$, and define the linear isomorphism
$T:\mathbb R^3\to W_0$ by $T u_i=b_i$, where $u_1,u_2,u_3$ is the standard
basis of $\mathbb R^3$.  For a permutation $\sigma\in S_3$, put
\begin{equation}
 K_\sigma=\operatorname{conv}\{0,u_{\sigma(1)},
 u_{\sigma(1)}+u_{\sigma(2)},u_1+u_2+u_3\}.
 \label{eq:A3-reference-simplex}
\end{equation}
The simplices $hT(k+K_\sigma)$, $k\in\mathbb Z^3$, form a periodic
shape-regular triangulation, with a shape-regularity constant independent
of $h$.  For $h=d^{-1}$, $d\ge3$, it descends to a simplicial mesh on $M$
with vertex set $V_h$.  The endpoints of every mesh edge can be joined by
at most three root steps.  If $I_h f$ is the continuous piecewise-affine interpolant of
$f:V_h\to\mathbb R$, then
\begin{equation}
 \|\nabla I_h f\|_{L^\infty(M)}\le\sqrt{24}\,L_h(f).
 \label{eq:A3-interpolation-bound}
\end{equation}
The cells $Q_{h,x}=x+hT[0,1)^3$, interpreted on the torus, partition $M$
up to boundaries, have normalized Haar measure $h^3$, and have diameter
at most $C h$.  If $\widehat f_h$ equals $f(x)$ on $Q_{h,x}$, then
\begin{equation}
 \|\widehat f_h-I_h f\|_\infty\le C h L_h(f).
 \label{eq:A3-cell-interpolation-bound}
\end{equation}
\end{lemma}

\begin{proof}
Every $x\in\Lambda_{A_3}$ can be written uniquely as
$x=\sum_{i=1}^3x_i b_i$, so $T\mathbb Z^3=\Lambda_{A_3}$.
The six simplices $K_\sigma$ partition the unit cube: the simplex indexed
by $\sigma$ is the region
$1\ge t_{\sigma(1)}\ge t_{\sigma(2)}\ge t_{\sigma(3)}\ge0$.
On a common face of adjacent cubes the induced ordering of the remaining
two coordinates agrees, so the translated subdivisions fit face to face.
Applying the fixed isomorphism $T$ preserves this compatibility.  There
are only six nondegenerate reference simplices, and scaling by $h$ leaves
the ratio of diameter to inradius unchanged.  This proves uniform shape
regularity.  Translations by $\Lambda_{A_3}$ correspond to translations
by $d\mathbb Z^3$ before scaling, which gives the periodic quotient mesh.

Consecutive vertices in the defining chain of each simplex differ by
$h b_{\sigma(i)}$.  These are root edges, and any pair of its vertices
is joined by a subchain of at most three such edges.  By
\eqref{eq:supp-A3-Lip}, each consecutive difference in the values of $f$
has absolute value at most $\sqrt8\,hL_h(f)$.  On one simplex let
$v=\nabla I_h f$, which is constant there.  The three consecutive
differences therefore give
$|\langle v,b_i\rangle|\le\sqrt8\,L_h(f)$ for $i=1,2,3$.
Since $T^*T=I_3+\one_3\one_3^{\mathsf T}$ has eigenvalues $1,1,4$,
the smallest singular value of $T^*:W_0\to\mathbb R^3$ is one.  Hence
\[
 |v|^2\le|T^*v|^2
 =\sum_{i=1}^3|\langle v,b_i\rangle|^2\le24L_h(f)^2,
\]
which proves \eqref{eq:A3-interpolation-bound}.
Here $\one_3=(1,1,1)^{\mathsf T}$.

The image under $T$ of the unit cube is a fundamental parallelepiped of
$\Lambda_{A_3}$.  Its $h$-scaled translates have normalized measure $h^3$
and uniformly bounded diameter divided by $h$.  On $Q_{h,x}$, integrate
the piecewise-constant gradient of $I_h f$ along the segment from $x$ to
the evaluation point.  The preceding gradient bound gives
\eqref{eq:A3-cell-interpolation-bound}.
\end{proof}

\begin{lemma}
\label{lem:supp-A3-compactness}
Let $h_n\downarrow0$ and let $f_n:V_{h_n}\to\R$ satisfy
$L_{h_n}(f_n)\le1$ and $f_n(x_{0,n})=0$ for $x_{0,n}\to x_0$.  Then a
subsequence of the interpolants $I_{h_n}f_n$ of
Lemma~\ref{lem:A3-periodic-triangulation} converges uniformly to
$F\in W^{1,\infty}(M)$ with
\begin{equation}
 |\nabla F|\le1
 \quad\text{almost everywhere}.
 \label{eq:supp-A3-limit-Lip}
\end{equation}
\end{lemma}

\begin{proof}
From \eqref{eq:supp-A3-Lip}, for every oriented root $\alpha$ and every vertex,
\begin{equation}
 |f_n(x+h_n\alpha)-f_n(x)|\le\sqrt8\,h_n.
 \label{eq:A3-edge-bound}
\end{equation}
Use the triangulation of Lemma~\ref{lem:A3-periodic-triangulation} and
put $F_n=I_{h_n}f_n$.  Formula~\eqref{eq:A3-interpolation-bound} gives
$\|\nabla F_n\|_\infty\le\sqrt{24}$, uniformly in $n$.
Together with $F_n(x_{0,n})=0$ and compactness of $M$, this bounds
$\|F_n\|_\infty$ uniformly.  Arzel\`a--Ascoli yields, after passing to a
subsequence, uniform convergence $F_n\to F$ for a Lipschitz function $F$.
The task is now to improve this preliminary gradient bound to one.

For every $\alpha\in\Phi_{A_3}$ define the forward difference quotient on vertices
by
\[
 g_{n,\alpha}(x)
 =\frac{f_n(x+h_n\alpha)-f_n(x)}{h_n},
\]
and extend it as a piecewise-constant function on the cells
$Q_{h_n,x}$ of Lemma~\ref{lem:A3-periodic-triangulation}.
At every cell based at a lattice vertex,
\begin{equation}
 \frac18\sum_{\alpha\in\Phi_{A_3}}|g_{n,\alpha}|^2\le1.
 \label{eq:A3-vector-ball}
\end{equation}
Thus the vector $g_n=(g_{n,\alpha})_\alpha$ is bounded in
$L^\infty(M;\R^{12})$.  Pass to a weak-$*$ convergent subsequence
$g_n\stackrel{*}{\rightharpoonup}g$.  The set
\[
 K=\left\{z\in\R^{12}:\frac18\sum_\alpha|z_\alpha|^2\le1\right\}
\]
is closed and convex.  Since each $g_n$ takes values in $K$, separation of $K$ by
linear functionals shows that the weak-$*$ limit also satisfies $g(x)\in K$ for
almost every $x$.

It remains to identify the components of $g$.  Let $\varphi\in C^\infty(M)$ and
sample it on the lattice.  Periodic discrete summation by parts gives
\begin{align*}
 h_n^3\sum_{x\in V_{h_n}}g_{n,\alpha}(x)\varphi(x)
 &=-h_n^3\sum_{x\in V_{h_n}}f_n(x)
   \frac{\varphi(x)-\varphi(x-h_n\alpha)}{h_n}.
\end{align*}
To justify the passage to integrals, let $\widehat f_n$, $\varphi_n$ and
$\psi_{n,\alpha}$ be the cellwise-constant extensions of, respectively,
$f_n(x)$, $\varphi(x)$ and
$(\varphi(x)-\varphi(x-h_n\alpha))/h_n$.
The cell diameter bound, Taylor's formula and
\eqref{eq:A3-cell-interpolation-bound} give
\begin{align*}
 \|\widehat f_n-F_n\|_\infty&\le C h_n,\\
 \|\varphi_n-\varphi\|_\infty&\le C h_n\|\varphi\|_{C^1},\\
 \|\psi_{n,\alpha}-\alpha\cdot\nabla\varphi\|_\infty
 &\le C h_n\|\varphi\|_{C^2}.
\end{align*}
Since each cell has normalized measure $h_n^3$, the summation-by-parts
identity is exactly
\[
 \int_M g_{n,\alpha}\varphi_n
 =-\int_M\widehat f_n\psi_{n,\alpha}.
\]
The displayed estimates and the uniform bounds on $g_{n,\alpha}$ and
$F_n$ allow replacement of the cellwise test functions by their smooth
limits with an $O(h_n)$ error.  Weak-$*$ convergence of $g_{n,\alpha}$
and uniform convergence of $F_n$ therefore yield
\[
 \int_M g_\alpha\varphi
 =-\int_M F\,\alpha\cdot\nabla\varphi.
\]
Hence $g_\alpha=\alpha\cdot\nabla F$ distributionally, and therefore almost
everywhere.  Combining this identity with \eqref{eq:A3-vector-ball} and the
tight-frame identity \eqref{eq:supp-A3-frame} gives
\[
 |\nabla F|^2
 =\frac18\sum_\alpha|\alpha\cdot\nabla F|^2
 =\frac18\sum_\alpha|g_\alpha|^2\le1
\]
almost everywhere.
\end{proof}

\begin{lemma}[Uniform smoothing of the torus distance]
\label{lem:A3-uniform-distance-smoothing}
There are constants $\varepsilon_0>0$ and $C_M<\infty$ such that, for
every $y\in M$ and $0<\varepsilon<\varepsilon_0$, there is a function
$F_{y,\varepsilon}\in C^\infty(M)$ satisfying
\begin{equation}
 \|F_{y,\varepsilon}-d_g(\cdot,y)\|_\infty\le\varepsilon,
 \qquad
 \|\nabla F_{y,\varepsilon}\|_\infty\le1,
 \label{eq:A3-mollified-distance}
\end{equation}
and
\begin{equation}
 \sup_{y\in M}\|F_{y,\varepsilon}\|_{C^2(M)}
 \le C_M\varepsilon^{-1}.
 \label{eq:A3-uniform-smoothing-C2}
\end{equation}
All constants are independent of the centre $y$.
\end{lemma}

\begin{proof}
Choose a nonnegative smooth function $\eta$ supported in the Euclidean
unit ball of $W_0$, with integral one, and put
$\eta_\varepsilon(u)=\varepsilon^{-3}\eta(u/\varepsilon)$.
Let $r_y(\widetilde x)=d_g(\widetilde x+\Lambda_{A_3},y)$ be the periodic
lift of the distance function.  For any lift $\widetilde x$ of $x\in M$,
define
\[
 F_{y,\varepsilon}(x)=\int_{W_0}
       \eta_\varepsilon(u)r_y(\widetilde x-u)\,\dd u,
\]
where $\dd u$ is Euclidean volume on $W_0$.  Periodicity makes this
independent of the lift.  The function $r_y$ is one-Lipschitz, so
\[
 |F_{y,\varepsilon}(x)-d_g(x,y)|
 \le\int_{W_0}\eta_\varepsilon(u)|u|\,\dd u\le\varepsilon.
\]
Convolution preserves the one-Lipschitz bound, giving the gradient
estimate in \eqref{eq:A3-mollified-distance}.  Equivalently,
$\nabla F_{y,\varepsilon}=\eta_\varepsilon*\nabla r_y$ with
$|\nabla r_y|\le1$ almost everywhere.
Differentiating this identity once more gives, componentwise,
$\nabla^2F_{y,\varepsilon}=(\nabla\eta_\varepsilon)*\nabla r_y$.
Consequently
\[
 \|\nabla^2F_{y,\varepsilon}\|_\infty
 \le C\|\nabla\eta_\varepsilon\|_{L^1(W_0)}
 \le C\varepsilon^{-1}.
\]
Also $0\le F_{y,\varepsilon}\le\operatorname{diam}(M)$.
Taking $\varepsilon_0<1$ proves \eqref{eq:A3-uniform-smoothing-C2}.
The bounds are uniform in $y$ because the distance functions are
translations of one another on the flat torus.
\end{proof}

\begin{theorem}[Connes-distance convergence on the periodic $A_3$ lattice]
\label{thm:supp-A3-distance}
Let $d_h$ be the pure-state Connes distance defined by
\eqref{eq:supp-A3-Lip}, and let $d_g$ be the flat geodesic distance of the
metric determined by the jump covariance in
Proposition~\ref{prop:A3-covariance}.  Then
\begin{equation}
 \sup_{x,y\in V_h}|d_h(x,y)-d_g(x,y)|\longrightarrow0.
 \label{eq:supp-A3-distance}
\end{equation}
\end{theorem}

\begin{proof}
We prove the lower and upper bounds uniformly in the vertices.

For the lower bound, fix $0<\varepsilon<\varepsilon_0$ and use the
functions supplied by Lemma~\ref{lem:A3-uniform-distance-smoothing}.
The smooth-function consistency estimate in
Theorem~\ref{thm:supp-A3-finite} gives
\[
 L_h(F_{y,\varepsilon}|_{V_h})\le1+C h/\varepsilon,
\]
uniformly in $y$.  Dividing by this factor makes the sampled function
admissible in the Connes dual formula.  Formula~\eqref{eq:A3-mollified-distance}
and the uniform diameter bound then give
\begin{equation}
 \begin{split}
 d_h(x,y)
 &\ge\frac{F_{y,\varepsilon}(x)-F_{y,\varepsilon}(y)}
              {1+C h/\varepsilon}\\
 &\ge\frac{d_g(x,y)-2\varepsilon}{1+C h/\varepsilon}
 \ge d_g(x,y)-C'\left(\varepsilon+\frac h\varepsilon\right).
 \end{split}
 \label{eq:A3-uniform-distance-lower-bound}
\end{equation}
The constants are independent of $x,y,h$ and $\varepsilon$ in the
specified range.  First let $h\downarrow0$ and then
$\varepsilon\downarrow0$ to obtain the uniform lower bound.

For the upper bound, suppose the uniform upper estimate failed.  Then for some
$\varepsilon>0$ there would be $h_n\downarrow0$ and vertices $x_n,y_n$ such that
\[
 d_{h_n}(x_n,y_n)\ge d_g(x_n,y_n)+\varepsilon.
\]
Because the finite Lipschitz ball modulo constants is compact, the dual supremum is
attained; choose $f_n$ with $L_{h_n}(f_n)\le1$, normalize
$f_n(y_n)=0$, and arrange
\begin{equation}
 f_n(x_n)\ge d_g(x_n,y_n)+\varepsilon.
 \label{eq:A3-upper-contradiction}
\end{equation}
After a subsequence, compactness of $M$ gives $x_n\to x$ and $y_n\to y$.
Lemma~\ref{lem:supp-A3-compactness}, applied with base points $y_n$, gives
uniform convergence of the corresponding interpolants to a function $F$ with
$|\nabla F|\le1$ almost everywhere.  The uniform edge Lipschitz bounds imply
$F(y)=0$ and $f_n(x_n)\to F(x)$.  Continuity of $d_g$ gives
$d_g(x_n,y_n)\to d_g(x,y)$.  Passing to the limit in
\eqref{eq:A3-upper-contradiction} yields
\[
 F(x)-F(y)\ge d_g(x,y)+\varepsilon,
\]
contradicting the fact that $|\nabla F|\le1$ implies that $F$ is
one-Lipschitz for the flat Riemannian distance.  This proves the uniform upper
bound and hence \eqref{eq:supp-A3-distance}.
\end{proof}


Each finite $A_3$ approximant is a compact quantum metric space, but the theorem
above compares only pure states, namely vertices.  A quadratic graph commutator
seminorm need not equal the ordinary Lipschitz seminorm of its pure-state metric
\cite{Rieffel1999}.  Consequently, the displayed distance convergence is not by
itself a quantum Gromov--Hausdorff, propinquity, or spectral-propinquity theorem
\cite{Rieffel2004,Latremoliere2022,FarsiLatremolierePacker2024}.
The stronger approximation results of
\cite{LandryLapidusLatremoliere2021,vanSuijlekom2021} control additional state-space
or operator data.  Lemma~\ref{lem:supp-A3-compactness} supplies the scalar
precompactness needed here; an upgrade would also require an explicit comparison
of the full Lip-norm structures.  The extensive one-form kernel remains a
separate obstruction to retaining the full Hodge operator in a compact-resolvent
limit.  The operator theorem of Miranda--Parra \cite{MirandaParra2023}
concerns a different differential complex with higher-degree cochains;
it is complementary to this pure-state metric limit.

\begin{remark}[Comparison with endpoint triples]
\label{subsec:endpoint-pointer}
Section~\ref{app:endpoint-metric} gives an independent construction on
arbitrary closed connected Riemannian manifolds.  A two-point spectral block
on each short mesh edge has the graph Lipschitz seminorm exactly, so its
Connes distance is the shortest-path metric.  That direct approximation
should be distinguished from the quadratic, covariance-normalized Hodge
seminorm used in Theorem~\ref{thm:supp-A3-distance}.
\end{remark}

\section{Stationary-current fibres and affinity coordinates}
\label{sec:inverse}
% =============================================================================

The graph Hodge--Dirac construction depends only on symmetric jump data and
therefore forgets the stationary current.  We first describe this fibre as a
cycle-current polytope.  The principal result of this section then identifies
its strictly positive relative interior with $H^1(G;\mathbb R)$ through the
affinity class, giving a global coordinate for the directed generator at
fixed metric data.  The section ends with the behaviour of currents and
entropy production under graph quotients.

\subsection{Directed currents at fixed metric geometry}

Fix a connected undirected graph $G=(V,E)$, positive masses $m_x$, and positive
conductances $c_e$.  Choose one orientation of each undirected edge and let
$B:\R^E\to\R^V$ be the incidence matrix.  An oriented current is extended to
the reverse orientation by $j_{\bar e}=-j_e$.

\begin{definition}[Cycle-current polytope and its positive interior]
\label{def:supp-current-polytope}
Define the closed cycle-current polytope
\begin{equation}
 \overline{\mathcal J}(G,c)
 =\{j\in\Ker B:\abs{j_e}\le c_e\text{ for every }e\in E\},
 \label{eq:supp-current-polytope-closed}
\end{equation}
and its relative interior
\begin{equation}
 \mathcal J(G,c)
 =\{j\in\Ker B:\abs{j_e}<c_e\text{ for every }e\in E\}.
 \label{eq:supp-current-polytope}
\end{equation}
The closed polytope allows vanishing directed rates.  The open fibre
$\mathcal J(G,c)$ consists precisely of currents for which both directional fluxes
$c_e\pm j_e$ are strictly positive on every retained edge; it is the natural domain
of the affinity coordinate $\operatorname{arctanh}(j/c)$ used below.
\end{definition}

The decomposition $q=c+j$ is the usual separation of traffic and net current.
For comparison, the stationary coordinates in \cite{YangDill2026} have node
probability $\pi_x=m_x$, edge traffic $\tau_{xy}=2c_{xy}$ and net flux
$J_{xy}=2j_{xy}$, so their formula
$k_{xy}=(\tau_{xy}+J_{xy})/(2\pi_x)$ is exactly
$k_{xy}=(c_{xy}+j_{xy})/m_x$.  The next theorem records what this established
parametrization means for a fixed graph Hodge--Dirac triple.

\begin{theorem}[Stationary Markov realization fibre of a metric graph triple]
\label{thm:supp-graph-realization}
Let
$\mathfrak T(G,m,c)=(C(V),\Hh_{m,c},D_{m,c},J,\gamma)$ be the metric graph
spectral triple of \Cref{thm:supp-graph-square}.  Then:
\begin{enumerate}[label=\textup{(R\arabic*)},leftmargin=3em]
\item it has the reversible Markov realization
      \begin{equation}
      k_{xy}^{\rm rev}=\frac{c_{xy}}{m_x};
      \label{eq:supp-reversible-realization}
      \end{equation}
\item every $j\in\overline{\mathcal J}(G,c)$ gives a stationary Markov realization
      \begin{equation}
      q_{xy}=c_{xy}+j_{xy},
      \qquad
      k_{xy}=\frac{c_{xy}+j_{xy}}{m_x};
      \label{eq:supp-current-realization}
      \end{equation}
      currents in $\mathcal J(G,c)$ are exactly those for which both directions on
      every retained edge have positive flux;
\item conversely, every stationary Markov generator with masses $m$ and
      symmetric conductances $c$ is uniquely of the form
      \eqref{eq:supp-current-realization} with
      $j\in\overline{\mathcal J}(G,c)$;
\item every realization in the fibre has the same metric spectral triple.
\end{enumerate}
Thus, up to unobserved parallel-channel refinements,
\begin{equation}
 {
 \RReal_{\rm graph}(\mathfrak T(G,m,c))
 \cong\overline{\mathcal J}(G,c),}
 \label{eq:supp-realization-fibre}
\end{equation}
while $\mathcal J(G,c)$ is its strictly positive relative interior.
\end{theorem}

\begin{proof}
For \eqref{eq:supp-reversible-realization}, the stationary flux is
$q_{xy}=c_{xy}=q_{yx}$, so detailed balance holds.  For
\eqref{eq:supp-current-realization}, nonnegativity follows from
$|j_{xy}|\le c_{xy}$; strict inequality is equivalent to positive flux in both
directions.  Moreover
\[
 \sum_y(q_{xy}-q_{yx})=2\sum_yj_{xy}=0
\]
because $Bj=0$, so $m$ is stationary.  Conversely, for a stationary continuous-time Markov
generator define $j=(q-q^{\mathsf T})/2$.  Stationarity gives $Bj=0$, positivity gives
the edge inequalities, and $q=c+j$ is unique.  The graph spectral triple uses
only $m$ and $(q+q^{\mathsf T})/2=c$.
\end{proof}

\begin{corollary}
\label{cor:supp-fibre-dimension}
For connected $G$, both the closed polytope
$\overline{\mathcal J}(G,c)$ and its relative interior $\mathcal J(G,c)$ have affine
dimension
\begin{equation}
 {|E|-|V|+1=b_1(G).}
 \label{eq:supp-fibre-dimension}
\end{equation}
The reversible realization is the only stationary realization exactly when $G$
is a tree.  If
$b_1(G)>0$, the metric passage to spectral data is not faithful.
\end{corollary}

\begin{proof}
The incidence matrix of a connected graph has rank $|V|-1$, so
$\dim\Ker B=|E|-|V|+1$.  The strict inequalities cut out an open neighborhood
of zero in this cycle space.  A connected graph has zero first Betti number
exactly when it is a tree.
\end{proof}

The stationary entropy-production functional used below is the standard
network expression for continuous-time Markov chains
\cite{Schnakenberg1976}.

\begin{proposition}
\label{prop:supp-entropy}
For $j\in\mathcal J(G,c)$, the steady entropy production is
\begin{equation}
 {
 \sigma(j)
 =2\sum_{e\in E}j_e
   \log\frac{c_e+j_e}{c_e-j_e}\ge0.}
 \label{eq:supp-entropy}
\end{equation}
It vanishes exactly at $j=0$.  On the boundary of
$\overline{\mathcal J}(G,c)$ the same formula is interpreted in the extended sense,
with value $+\infty$ when a positive-conductance edge has a vanishing directional
flux.
\end{proposition}

\begin{proof}
Start from
$\frac12\sum_{x,y}(q_{xy}-q_{yx})\log(q_{xy}/q_{yx})$, pair the two
orientations of each edge, and substitute $q_e=c_e+j_e$,
$q_{\bar e}=c_e-j_e$.  Each summand is nonnegative because $j_e$ and the
logarithm have the same sign; it vanishes only for $j_e=0$.
\end{proof}

\begin{example}[A stationary family on the three-cycle]
On an oriented three-cycle with equal conductance $c$, every constant current
$j_e=s$, $|s|<c$, is divergence free.  All values of $s$ give one metric
spectral triple, whereas
\[
 \sigma(s)=6s\log\frac{c+s}{c-s}
\]
varies continuously from zero and diverges as $|s|\uparrow c$.
\end{example}

Thus whenever $b_1(G)>0$, the map from stationary generators
with fixed $(G,m,c)$ to the metric graph triple is noninjective on a
$b_1(G)$-dimensional open set.  Quotienting by the finite automorphism group
preserving $(m,c)$ does not remove this positive-dimensional ambiguity.  This
consequence will be refined below by replacing the current coordinate with an
intrinsic affinity class.

\subsection{Cohomological coordinates of the stationary-current fibre}

For $j\in\mathcal J(G,c)$ define the real affinity
\begin{equation}
 a_e=\operatorname{arctanh}\frac{j_e}{c_e}
 =\frac12\log\frac{c_e+j_e}{c_e-j_e}.
 \label{eq:supp-affinity}
\end{equation}
Exact one-cochains are $\Ran B^{\mathsf T}$, so $a$ has a class
$[a]\in H^1(G;\R)=\R^E/\Ran B^{\mathsf T}$.

\begin{lemma}
\label{lem:supp-zero-affinity}
If $j\in\mathcal J(G,c)$ is divergence free and $[a]=0$, then $j=0$.
\end{lemma}

\begin{proof}
Write $a=B^{\mathsf T}\phi$.  Since
$j_e=c_e\tanh((B^{\mathsf T}\phi)_e)$ and $Bj=0$,
\[
 0=\langle\phi,Bj\rangle
 =\sum_ec_e(B^{\mathsf T}\phi)_e
       \tanh((B^{\mathsf T}\phi)_e).
\]
Every summand is nonnegative and vanishes only when the edge difference is
zero.  Connectedness makes $\phi$ constant, hence $a=j=0$.
\end{proof}

The preceding variational theory selects a nonlinear coclosed representative
of a cohomology class \cite[Proposition~V.7]{PardoGuerraThapaWashburn2026}.
Here $\psi(t)=\log\cosh t$ has $\psi'(t)=\tanh t$ and
$\psi''(t)=\operatorname{sech}^2t>0$.  Its conductance-weighted Euler
equation is exactly the stationarity condition $B(c\tanh a)=0$.
The theorem identifies these representatives with all currents satisfying
$|j_e|<c_e$ and establishes the smooth global coordinate needed to recover
the directed Markov generator.  We include the variational proof with the
weights, positivity domain and smooth dependence made explicit.

\begin{theorem}[Affinity coordinates for a fixed-traffic Markov fibre]
\label{thm:supp-nonlinear-Hodge}
The map
\begin{equation}
 \mathfrak A:\mathcal J(G,c)\longrightarrow H^1(G;\R),
 \qquad
 j\longmapsto[\operatorname{arctanh}(j/c)]
 \label{eq:supp-affinity-map}
\end{equation}
is a diffeomorphism.  More explicitly, for a class $\alpha$ and representative
$a_0$, define on mean-zero vertex potentials
\begin{equation}
 \mathcal F_\alpha(\phi)
 =\sum_{e\in E}c_e
  \log\cosh(a_{0,e}+(B^{\mathsf T}\phi)_e).
 \label{eq:supp-convex-affinity-functional}
\end{equation}
This functional is coercive and strictly convex.  Its unique minimizer
$\phi_\alpha$ gives
\begin{equation}
 a_\alpha=a_0+B^{\mathsf T}\phi_\alpha,
 \qquad
 j_{\alpha,e}=c_e\tanh(a_{\alpha,e}),
 \qquad Bj_\alpha=0,
 \label{eq:supp-affinity-inverse}
\end{equation}
and this is the unique current with affinity class $\alpha$.
\end{theorem}

\begin{proof}
If $E=\varnothing$, connectedness gives $|V|=1$ and both spaces are singletons,
so the assertion is immediate.  Assume henceforth that $E\ne\varnothing$.
Write $V_0=\{\phi\in\mathbb R^V:\sum_x\phi_x=0\}$ and
$c_* = \min_e c_e>0$.  Connectedness gives
$\|\phi\|\le C\|B^{\mathsf T}\phi\|$ on $V_0$.  The inequality
$\log\cosh t\ge |t|-\log2$ therefore gives
\[
 \mathcal F_\alpha(\phi)
 \ge c_*\|B^{\mathsf T}\phi\|_1
        -\sum_e c_e|a_{0,e}|-(\log2)\sum_e c_e
 \longrightarrow\infty
\]
as $\|\phi\|\to\infty$ in $V_0$.  Its Hessian is
\[
 B\,\diag\!\left(c_e\operatorname{sech}^2
  (a_{0,e}+(B^{\mathsf T}\phi)_e)\right)B^{\mathsf T},
\]
which is positive definite on $V_0$: all edge weights are positive and
$\Ker B^{\mathsf T}$ consists of constants.  Thus the functional is coercive and
strictly convex and has a unique minimizer.  The Euler equation is
\[
 0=\sum_ec_e\tanh(a_{\alpha,e})(B^{\mathsf T}\psi)_e
   =\langle Bj_\alpha,\psi\rangle\qquad(\psi\in V_0).
\]
Since $Bj_\alpha$ has zero vertex sum, it follows that $Bj_\alpha=0$.
Also $|\tanh a_{\alpha,e}|<1$, so $j_\alpha\in\mathcal J(G,c)$.

Replacing $a_0$ by $a_0+B^{\mathsf T}\chi$, with $\chi\in V_0$, replaces the
minimizer by $\phi_\alpha-\chi$ and leaves $a_\alpha$ and $j_\alpha$ unchanged.
The construction therefore depends only on the class $\alpha$.
Conversely, if $j\in\mathcal J(G,c)$ and
$a=\operatorname{arctanh}(j/c)$, stationarity makes zero a critical point of
\eqref{eq:supp-convex-affinity-functional} with representative $a$.
Strict convexity makes it the unique minimizer.  Hence the two constructions are
mutual inverses.

Finally choose a linear section of
$\mathbb R^E\twoheadrightarrow H^1(G;\mathbb R)$ and use it to select $a_0$
smoothly.  The stationarity equation, as a map to $V_0$, has derivative in
$\phi$ equal to the displayed positive-definite Hessian.  The implicit-function
theorem gives smooth local dependence on $\alpha$; uniqueness makes these local
solutions agree globally.  The forward map is smooth because componentwise
$\operatorname{arctanh}(j/c)$ is smooth on the open current polytope.
This proves the diffeomorphism.
\end{proof}

\begin{corollary}
\label{cor:supp-faithful-modular}
For fixed $(G,m,c)$, the graph triple together with its affinity class
\begin{equation}
 \mathfrak S_{\rm aff}(j)
 =\bigl(C(V),\Hh_{m,c},D_{m,c},J,\gamma;[a]\bigr)
 \label{eq:supp-modular-enrichment}
\end{equation}
determines the complete stationary generator through
\[
 [a]\longmapsto a_\alpha
 \longmapsto j=c\tanh a_\alpha
 \longmapsto q=c+j
 \longmapsto k_{xy}=q_{xy}/m_x.
\]
Thus the affinity class determines the directed generator throughout the strictly positive fibre.
\end{corollary}

\begin{proof}
Apply \Cref{thm:supp-nonlinear-Hodge} and
\Cref{thm:supp-graph-realization}.
\end{proof}

The entropy production is the current--affinity pairing
\begin{equation}
 {\sigma(j)=4\sum_{e\in E}j_ea_e,}
 \label{eq:supp-entropy-modular-pairing}
\end{equation}
and is unchanged by $a\mapsto a+B^{\mathsf T}\phi$ because
$\langle j,B^{\mathsf T}\phi\rangle=\langle Bj,\phi\rangle=0$.

\subsection{Coarse graining of stationary currents}
\label{subsec:current-coarse-graining}

The affinity coordinate reconstructs each current at fixed conductances.
Under a graph quotient, however, the natural operation is to sum the
directional fluxes, and hence to sum currents and conductances.  The
following result records this operation and the associated entropy bound.

Write $B_G=B:\mathbb R^E\to\mathbb R^V$ for the incidence matrix and
$\mathcal Z(G)=\Ker B_G$.  A vertex quotient $q:V\twoheadrightarrow\bar V$
induces maps $Q_V$ and $Q_E$ by summing vertex and oriented edge coefficients over
fibres.

\begin{theorem}[Graph quotients and entropy contraction]
\label{thm:supp-current-chain}
The incidence maps satisfy
\begin{equation}
{B_{\bar G}Q_E=Q_VB_G,}
\label{eq:supp-current-chain}
\end{equation}
so $Q_E\mathcal Z(G)\subseteq\mathcal Z(\bar G)$.  If
$\bar c_{\bar e}=\sum_{e\mapsto\bar e}c_e$, then
$Q_E\mathcal J(G,c)\subseteq\mathcal J(\bar G,\bar c)$.  When the vertex fibres are
connected and the coarse graph retains every inter-fibre edge class, the induced
map on cycle spaces is onto and
\begin{equation}
0\longrightarrow\Ker(Q_E|_{\mathcal Z(G)})
\longrightarrow\mathcal Z(G)
\xrightarrow{Q_E}\mathcal Z(\bar G)
\longrightarrow0,
\label{eq:supp-cycle-exact}
\end{equation}
with
\begin{equation}
\dim\Ker(Q_E|_{\mathcal Z(G)})=b_1(G)-b_1(\bar G).
\label{eq:supp-cycle-kernel-dim}
\end{equation}
The stationary entropy production satisfies
\begin{equation}
\sigma_{\bar G}(Q_Ej)\le\sigma_G(j).
\label{eq:supp-entropy-contraction}
\end{equation}
\end{theorem}

\begin{proof}
For a coarse vertex, the divergence of the summed inter-fibre current equals the
sum of fine divergences; internal fine edges cancel in pairs.  This is exactly
\eqref{eq:supp-current-chain}.  The conductance inequality
$|(Q_Ej)_{\bar e}|\le\sum_{e\mapsto\bar e}|j_e|<\bar c_{\bar e}$ gives the fibre
map.  Connected fibres allow every coarse cycle to be lifted and completed inside
fibres, proving surjectivity.  The dimension formula follows from the Betti-number
formula for connected graphs.  Finally, on every coarse edge the log-sum inequality
for the aggregated forward and reverse fluxes gives entropy contraction; summing
over coarse edges yields \eqref{eq:supp-entropy-contraction}.
\end{proof}

The surjectivity assertion concerns the linear cycle spaces.  It does not
assert surjectivity of the bounded current polytopes.  The graph quotient
therefore defines a map of admissible current fibres without identifying
every coarse current with an admissible fine current.

Section~\ref{app:edge-affinity} develops the edge-ratio and cellular Hodge
interpretations of the same affinity class, including its curvature and
holonomy data.

\section{Independent endpoint metric approximation}
\label{app:endpoint-metric}

The periodic $A_3$ metric theorem uses the Hodge--Dirac operator associated with a
stationary graph and therefore has a quadratic local commutator seminorm.  The
following construction is independent of that choice.  It assigns one elementary
spectral block to each sufficiently short edge of a dense mesh.  Edge-based
triples recovering graph geodesic distance are standard; see, for example, the
auxiliary edge construction in \cite{Besnard2021}.  Here we record the resulting
uniform approximation on a closed manifold.  The commutator seminorm is exactly
the ordinary graph Lipschitz seminorm, so the comparison is explicit on all pure
states.

Let $M$ be closed and connected.  For each $h\downarrow0$, let $V_h\subset M$ be an $h$-net,
so every point of $M$ lies within distance $h$ of a vertex.  Choose a mesoscopic
radius $\rho_h\downarrow0$ with $h/\rho_h\to0$ and, for all sufficiently small
$h$, $\rho_h$ below the injectivity radius.  Join every pair
$x,y\in V_h$ with $d_g(x,y)\le\rho_h$ and assign that edge the length
\begin{equation}
 \ell_h(x,y)=d_g(x,y),
 \label{eq:supp-metric-edge-length}
\end{equation}
or, more generally, a positive length satisfying
$\sup_{E_h}|\ell_h/d_g-1|\to0$.  For an unoriented edge $e=\{x,y\}$ set
\begin{equation}
 \Hh_e=\C^2,\qquad
 \pi_e(f)=\begin{pmatrix}f(x)&0\\0&f(y)\end{pmatrix},\qquad
 D_e=\frac1{\ell_h(x,y)}\begin{pmatrix}0&1\\1&0\end{pmatrix}.
 \label{eq:supp-endpoint-block}
\end{equation}
Let
\begin{equation}
 \Hh_h^{\rm met}=\bigoplus_{e\in E_h}\Hh_e,\qquad
 \pi_h=\bigoplus_{e\in E_h}\pi_e,\qquad
 D_h^{\rm met}=\bigoplus_{e\in E_h}D_e.
 \label{eq:supp-endpoint-triple}
\end{equation}
This is the finite endpoint metric triple used below.

\begin{theorem}[Endpoint spectral triples and geodesic convergence]
\label{thm:supp-global-metric}
Under the preceding mesh assumptions, for $f\in C(V_h)$,
\begin{equation}
{
L_h^{\rm met}(f)
=\|[D_h^{\rm met},\pi_h(f)]\|
=\max_{\{x,y\}\in E_h}\frac{|f(x)-f(y)|}{\ell_h(x,y)}.}
\label{eq:supp-edge-lip}
\end{equation}
The pure-state Connes distance is exactly the graph shortest-path metric.  On every
closed Riemannian manifold,
\begin{equation}
{
\sup_{x,y\in V_h}
|d_h^{\rm met}(x,y)-d_g(x,y)|\longrightarrow0.}
\label{eq:supp-global-metric-convergence}
\end{equation}
For smooth $f$, $L_h^{\rm met}(f|_{V_h})\to\|\nabla_gf\|_\infty$.
\end{theorem}

\begin{proof}
The commutator of one endpoint block has off-diagonal entry
$(f(y)-f(x))/\ell_h(x,y)$, so the operator norm gives
\eqref{eq:supp-edge-lip}.  Telescoping along graph paths gives one inequality in the
Connes dual formula; the distance-to-a-fixed-vertex function has seminorm one and
gives equality.

When $\ell_h=d_g$ on edges, every graph path is a chain of geodesic segments,
so the triangle inequality gives $d_g\le d_h^{\rm met}$; the same conclusion
holds with an $o(1)$ multiplicative error for the allowed approximate lengths.
Conversely, fix $x,y\in V_h$ and partition a minimizing geodesic from $x$ to $y$
into pieces of length at most $\rho_h/2$.  Replace each interior partition point
by a vertex of the $h$-net.  For small $h$, consecutive replacement vertices are
at distance at most $\rho_h/2+2h<\rho_h$ and are therefore connected.  The number
of pieces is $O(1+d_g(x,y)/\rho_h)$, so the accumulated replacement error is
$O(h+h/\rho_h)$ uniformly on the compact manifold.  Allowing the final
partition remainder gives, uniformly,
\[
 d_h^{\rm met}(x,y)
 \le d_g(x,y)+C\left(\rho_h+\frac h{\rho_h}\right)+o(1).
\]
Hence \eqref{eq:supp-global-metric-convergence} follows.

For smooth $f$, the upper bound on $L_h^{\rm met}(f|_{V_h})$ follows from the
geodesic mean-value estimate.  For the reverse bound, choose $p$ where
$|\nabla f|$ is within $o(1)$ of its maximum, move a distance $\rho_h/2$ from
$p$ in the gradient direction, and approximate both endpoints by vertices of
the $h$-net.  Taylor expansion and $h/\rho_h\to0$ give the matching lower bound.
Thus $L_h^{\rm met}(f|_{V_h})\to\|\nabla_gf\|_\infty$.
\end{proof}

For these endpoint metric triples the finite Lip-norm is exactly the
ordinary Lipschitz seminorm of the graph shortest-path distance.  The metric
on $S(C(V_h))$, identified with probability measures on $V_h$, is therefore
the Kantorovich transport distance for $d_h^{\rm met}$
\cite{Rieffel1999}.  If $p,q$ are probability measures supported on $V_h$,
the same transport plans give
\[
 \left|\operatorname{mk}_{L_h^{\rm met}}(p,q)-W_{1,d_g}(p,q)\right|
 \le\sup_{x,y\in V_h}|d_h^{\rm met}(x,y)-d_g(x,y)|.
\]
Here $W_{1,d_g}$ denotes the transport distance between the same measures
viewed on $M$.  Unlike the general quadratic graph seminorm, this particular
endpoint construction thus yields an immediate comparison on all states,
not only on vertices.

\section{Edge affinities and cellular Hodge decomposition}
\label{app:edge-affinity}

This section records two auxiliary descriptions of the affinity used in
Section~\ref{sec:inverse}: forward/reverse edge-density ratios and their
cellular Hodge decomposition.  They relate the stationary-current coordinate
to transport, local curvature and global holonomy.

\subsection{Relative edge-density data}

Let $E^{\rm or}$ contain both orientations of every edge, with reversal
$e\mapsto\bar e$, source $s(e)$, terminal $t(e)=s(\bar e)$, and positive
directed stationary flux $q_e$.  On $\Hh_E=\ell^2(E^{\rm or})$, define
\begin{equation}
 (J_E\xi)(e)=\overline{\xi(\bar e)},
 \qquad
 (\Delta_q\xi)(e)=\frac{q_e}{q_{\bar e}}\xi(e)=e^{2a_e}\xi(e),
 \qquad
 S_q=J_E\Delta_q^{1/2}.
 \label{eq:edge-relative-modular-data}
\end{equation}
Using the same left/right endpoint convention as in
\eqref{eq:supp-graph-actions}, put
$(\lambda_E(f)\xi)(e)=f(t(e))\xi(e)$ and
$(\rho_E(f)\xi)(e)=f(s(e))\xi(e)$.
The involution $J_E$ is auxiliary edge reversal; it is not the signed
one-form real structure in \eqref{eq:graph-real-even-structure}.

\begin{theorem}[Edge reversal and relative density ratios]
\label{thm:edge-relative-modular-identities}
The edge data \eqref{eq:edge-relative-modular-data} satisfy
\begin{equation}
 {
 J_E^2=I,
 \qquad
 J_E\Delta_qJ_E=\Delta_q^{-1},
 \qquad
 S_q^2=I,}
 \label{eq:edge-reversal-identities}
\end{equation}
\begin{equation}
 {
 J_E\lambda_E(f^*)J_E=\rho_E(f),
 \qquad
 S_q\lambda_E(f^*)S_q^{-1}=\rho_E(f),
 \qquad
 \log\Delta_q(e)=2a_e.}
 \label{eq:edge-opposite-identities}
\end{equation}
The reversible branch is exactly $\Delta_q=I$.
\end{theorem}

\begin{proof}
Reversal preserves counting measure, so $J_E$ is antiunitary and squares to
the identity.  Since $a_{\bar e}=-a_e$, conjugating the diagonal multiplier
$e^{2a_e}$ by $J_E$ gives its inverse; hence $S_q^2=I$.  Reversal exchanges
the terminal and source of an edge, proving the opposite-action identities.
The logarithm statement is the definition of $a_e$.
\end{proof}

\begin{remark}[Scope of the modular terminology]
Here $\Delta_q$ is the multiplication operator given by the density ratio
of the forward and reversed edge measures.  No von Neumann algebra with a
cyclic separating vector is specified.  The displayed reversal identities
are therefore not asserted to construct a Tomita--Takesaki modular operator.
\end{remark}

For an oriented cycle $C$,
\begin{equation}
 \operatorname{Hol}_{\Delta}(C)
 =\prod_{e\in C}\left(\frac{q_e}{q_{\bar e}}\right)^{\epsilon_{C,e}}
 =\exp\left(2\sum_{e\in C}\epsilon_{C,e}a_e\right).
 \label{eq:supp-modular-cycle-holonomy}
\end{equation}

\subsection{Cellular curvature and gauge classes}

Let
\begin{equation}
 C^0\xrightarrow{d_0}C^1\xrightarrow{d_1}C^2
 \xrightarrow{d_2}C^3
 \label{eq:supp-cellular-Hodge-complex}
\end{equation}
be a finite real Hilbert complex and let
$\mathscr H^1=\Ker d_1\cap\Ker d_0^*$.

\begin{proposition}
\label{thm:supp-modular-Hodge}
For $a\in C^1$, define
\begin{equation}
 h_a=P_{\mathscr H^1}a,
 \qquad
 F_a=d_1a.
 \label{eq:supp-modular-hF}
\end{equation}
Then $a$ has the unique orthogonal decomposition
\begin{equation}
 {a=d_0\phi+h_a+d_1^*\psi.}
 \label{eq:supp-modular-Hodge}
\end{equation}
The three orthogonal components are unique; the potentials are unique only after
requiring $\phi\perp\Ker d_0$ and $\psi\perp\Ker d_1^*$.
Its coexact part is reconstructed from the curvature by
\begin{equation}
 P_{\Ran d_1^*}a
 =d_1^*(d_1d_1^*)^\dagger F_a.
 \label{eq:supp-modular-coexact}
\end{equation}
Consequently
\begin{equation}
 C^1/\Ran d_0\longrightarrow\mathscr H^1\oplus\Ran d_1,
 \qquad
 [a]_{\rm gauge}\longmapsto(h_a,F_a)
 \label{eq:supp-modular-gauge-map}
\end{equation}
is injective and onto the displayed product.  A canonical representative is
\begin{equation}
 a_{h,F}=h+d_1^*(d_1d_1^*)^\dagger F.
 \label{eq:supp-modular-representative}
\end{equation}
Moreover $d_2F_a=0$ and
\begin{equation}
 \norm a^2
 =\norm{d_0\phi}^2+\norm{h_a}^2
 +\ip{F_a}{(d_1d_1^*)^\dagger F_a}.
 \label{eq:supp-modular-norm}
\end{equation}
\end{proposition}

\begin{proof}
Finite Hodge decomposition gives
$C^1=\Ran d_0\oplus\mathscr H^1\oplus\Ran d_1^*$, proving
\eqref{eq:supp-modular-Hodge}.  Since $d_1$ annihilates the first two
summands, $F_a=d_1d_1^*\psi$; Moore--Penrose inversion on $\Ran d_1$ gives
\eqref{eq:supp-modular-coexact}.  If two cochains have the same $(h,F)$, their
difference has zero harmonic projection and lies in $\Ker d_1$.  Its coexact
part then has zero norm, so the difference is exact.  This proves injectivity;
\eqref{eq:supp-modular-representative} proves surjectivity.  Orthogonality
gives the norm identity, and $d_2d_1=0$ gives the Bianchi identity.
\end{proof}

\begin{corollary}
\label{cor:supp-modular-trichotomy}
For a cellular affinity $a$:
\begin{enumerate}[label=\textup{(M\arabic*)},leftmargin=3em]
\item if $F_a\ne0$, the edge-ratio transport has local face curvature;
\item if $F_a=0$ but $h_a\ne0$, it is locally flat with nontrivial global positive
      holonomy;
\item if $F_a=h_a=0$, then $a=d_0\phi$ and the edge-ratio transport is trivialized by a
      positive vertex gauge.
\end{enumerate}
On the periodic three-torus, the global flat sector has dimension three.
\end{corollary}

\begin{theorem}
\label{thm:supp-face-holonomy}
For every oriented two-cell $p$,
\begin{equation}
 {
 \prod_{e\subset\partial p}
 \left(\frac{q_e}{q_{\bar e}}\right)^{\epsilon_{p,e}}
 =\exp(2F_a(p)).}
 \label{eq:supp-face-holonomy}
\end{equation}
Thus the edge-ratio transport defines a flat positive line transport on the
cell complex exactly when its curvature vanishes on every two-cell.
\end{theorem}

\begin{proof}
By the cellular coboundary formula,
$F_a(p)=\sum_e\epsilon_{p,e}a_e$.  Substitute
$q_e/q_{\bar e}=e^{2a_e}$ and multiply.
\end{proof}

\section{Operational reconstruction, realization and recognition}
\label{app:operational-recognition}

This section develops an operational source for the finite algebra
in Section~\ref{sec:finite}.  Complete intervention statistics first give a
future-separated history $C^*$-algebra through causal Choi tomography,
a prefix-minimal coherent realization, and a contextual quotient.  Finite
word-Gram data then reconstruct its algebraic operations.  The converse
realization and recognition results preserve representation multiplicities
within a boundary-relative convention.  The construction recovers the
algebraic input; the differential and its scale remain independent data.

Quantum combs and process tensors provide the multitime realization framework
\cite{ChiribellaDArianoPerinotti2009,BisioDArianoPerinottiChiribella2011,PollockEtAl2018};
future-equivalence quotients provide the predictive-state background
\cite{shalizicrutchfield,TraversCrutchfield2014,MarzenCrutchfield2017}.

\subsection{Reconstruction from complete intervention data}
\label{subsec:operational-reconstruction}

The causal Choi constraints and sequential-dilation ingredients used here are
standard in quantum-comb/process-tensor theory
\cite{ChiribellaDArianoPerinotti2009,PollockEtAl2018}; the future-equivalence
quotient is the finite typed analogue of minimal predictive-state constructions
\cite{shalizicrutchfield,TraversCrutchfield2014,MarzenCrutchfield2017}.  The new
object used below is the future-separated represented $C^*$-history algebra and
its finite flat word-Gram reconstruction.

Fix a finite cutoff $X$.  The operational input consists of a finite typed
event grammar with finite reachable histories.  If a history $h$ ends at cut type
$x$, let $\mathsf F_{X,x}$ be the finite admissible future words beginning at $x$.
The conditional event law assigns normalized one-step probabilities and therefore
cylinder probabilities
\begin{equation}
 p_X(w\mid h)
 =\prod_{j=1}^{k}p_X(a_j\mid ha_1\cdots a_{j-1}),
 \qquad w=a_1\cdots a_k,
\label{eq:ncg-chain-law}
\end{equation}
with $p_X(\epsilon_x\mid h)=1$.  No persistent state space, algebra, Hilbert
space, metric, or Dirac operator is assumed at this stage.

Two histories of the same cut type are \emph{future equivalent} when no declared
finite continuation distinguishes them:
\begin{equation}
 h\sim_X h'
 \quad\Longleftrightarrow\quad
 p_X(w\mid h)=p_X(w\mid h')
 \quad\text{for every }w\in\mathsf F_{X,x}.
\label{eq:ncg-future-equivalence}
\end{equation}
Histories of different cut type are never identified.  The relation is a typed
\emph{partial} right congruence on reachable branches.  Indeed, if $h\sim_Xh'$
and $a:x\to y$ satisfies $p_X(a\mid h)>0$, then the one-letter future gives
$p_X(a\mid h)=p_X(a\mid h')>0$, and for every continuation $v$,
\[
 p_X(v\mid ha)
 =\frac{p_X(av\mid h)}{p_X(a\mid h)}
 =\frac{p_X(av\mid h')}{p_X(a\mid h')}
 =p_X(v\mid h'a).
\]
Thus $ha\sim_Xh'a$.  No update state is assigned to a zero-probability extension;
its conditional continuation law is operationally undefined and is not needed in
the reachable predictive state space.

Define the minimal predictive state space by
\begin{equation}
 Z_{X,x}^{\min}:=\mathsf H_{X,x}/\!\sim_X,
 \qquad
 q_{X,x}:\mathsf H_{X,x}\twoheadrightarrow Z_{X,x}^{\min}.
\label{eq:ncg-predictive-state}
\end{equation}
For an admissible event $a:x\to y$, define
\begin{equation}
 \kappa_{X,a}([h])=p_X(a\mid h),
 \qquad
 \operatorname{Dom}(T_{X,a})
 =\{[h]:\kappa_{X,a}([h])>0\},
 \qquad
 T_{X,a}([h])=[ha].
\label{eq:ncg-predictive-update}
\end{equation}
Future equivalence makes $\kappa_{X,a}$ well defined on every class and
$T_{X,a}$ well defined on its displayed domain.  Thus persistence is reconstructed
from the conditional future law rather than introduced as an independent
coordinate.  Moreover, this quotient is universal among reachable deterministic
predictive realizations, where updates are required only on positive-probability
branches.  Equality of two realization states forces equality of all their future
cylinder probabilities, and therefore there is a unique surjection
\begin{equation}
 \pi_{X,x}:S_{X,x}\twoheadrightarrow Z_{X,x}^{\min}
\label{eq:ncg-predictive-universal}
\end{equation}
intertwining branch probabilities and all defined updates.  Hence $Z_X^{\min}$ is
the unique coarsest reachable deterministic predictive realization in this
partial-transition sense.

To reconstruct the operator-algebraic history object, assume in addition that the
finite multitime intervention tables are specified on tomographically complete
finite operator frames, are affine under randomization and additive under coarse
graining, and obey normalization, prefix compatibility and causal-discard
conditions.  Positivity here means positivity of the reconstructed multitime Choi
tensors, as required by completely positive process realization; nonnegative entries
in a product-frame table alone are not a substitute for this condition.  These
finite conditions will be stated explicitly in
\eqref{eq:ncg-comb-causality}.

For intervention slot $k$, let
\[
 \mathsf E_k=\operatorname{Herm}(H_k^{\rm out}\otimes H_k^{\rm in})
\]
be the real Choi-coordinate space.  Physical completely positive branches span
$\mathsf E_k$: a Hermitian element is the difference of its positive and negative
parts, after separately rescaling those parts so that their output partial traces
are bounded by the identity.  Choose physical bases $M_{k,\alpha}$ and
Hilbert--Schmidt duals $D_{k,\alpha}$.  Write $\mathfrak p_X$ for the resulting
multilinear intervention functional.  From the $N$-slot probability table define
\begin{equation}
 R_X^{(N)}
 =\sum_{\alpha_1,\ldots,\alpha_N}
 \mathfrak p_X(M_{1,\alpha_1},\ldots,M_{N,\alpha_N})
 D_{N,\alpha_N}\otimes\cdots\otimes D_{1,\alpha_1}.
 \label{eq:ncg-comb-tomography}
\end{equation}

We use the following standard finite-dimensional process realization as an input
\cite{ChiribellaDArianoPerinotti2009,BisioDArianoPerinottiChiribella2011,PollockEtAl2018}.  It is stated in the present
frame notation to make the subsequent quotient construction unambiguous.

\begin{proposition}[Standard finite process realization]
\label{thm:ncg-comb-tomography}
\label{thm:ncg-minimal-comb-dilation}
The complete finite intervention table determines a unique frame-independent
Hermitian tensor $R_X^{(N)}$ representing every multilinear intervention
probability.  It represents a deterministic causal process exactly when its prefixes
satisfy
\begin{equation}
 R_X^{(k)}\succeq0,
 \qquad
 \operatorname{Tr}_{H_k^{\rm out}}R_X^{(k)}
 =I_{H_k^{\rm in}}\otimes R_X^{(k-1)},
 \qquad R_X^{(0)}=1.
 \label{eq:ncg-comb-causality}
\end{equation}
Moreover, with $A_{X,0}=\mathbb C$ and $A_{X,k}$ the conjugate support of
$R_X^{(k)}$, there is a coherent sequential realization with
$\dim A_{X,k}=\rank R_X^{(k)}$ and isometries
\begin{equation}
 V_{X,k}:H_k^{\rm in}\otimes A_{X,k-1}
 \longrightarrow H_k^{\rm out}\otimes A_{X,k}
 \label{eq:ncg-minimal-isometries}
\end{equation}
whose sequential Choi link gives $R_X^{(N)}$.  Here \emph{support minimality}
means that the cumulative isometry through each cut is a minimal Stinespring
dilation of the corresponding prefix channel.  Any two such coherent sequential
realizations, with the same external port identifications and $U_{X,0}=I$, are
related by unique memory unitaries $U_{X,k}$ satisfying
\begin{equation}
 V'_{X,k}=(I\otimes U_{X,k})V_{X,k}(I\otimes U_{X,k-1}^*).
 \label{eq:ncg-minimal-unitary-covariance}
\end{equation}
Consequently the reduced primitive branch maps, their Hilbert--Schmidt adjoints, all
mixed history words, and their word Grams are transported by canonical unitary
equivalence.
\end{proposition}

\begin{proof}
Because the physical frame is a basis of the real Hermitian Choi space, its tensor
products determine the multilinear functional uniquely.  Expansion in the dual
frame gives \eqref{eq:ncg-comb-tomography}; uniqueness makes the resulting tensor
independent of the frame.  The positive causal-prefix criterion is the standard
comb realization theorem \cite{ChiribellaDArianoPerinotti2009}.  Existence of a
sequential isometric realization with
$\dim A_{X,k}=\rank R_X^{(k)}$ is
\cite[Theorem~2]{BisioDArianoPerinottiChiribella2011}.
We give the compatibility argument for uniqueness explicitly.

Suppress $X$ and put
\[
 I_k=\bigotimes_{j=1}^k H_j^{\rm in},\qquad
 O_k=\bigotimes_{j=1}^k H_j^{\rm out}.
\]
Let $W_k:I_k\to O_k\otimes A_k$ be the cumulative isometry through cut $k$,
and let $W'_k:I_k\to O_k\otimes A'_k$ be that of a second realization.
Both give the channel with Choi matrix $R^{(k)}$ after tracing out the last
memory.  Prefix minimality means
\begin{equation}
 \Span\{(\langle\eta|\otimes I)W_k\xi:
           \eta\in O_k,\ \xi\in I_k\}=A_k.
 \label{eq:ncg-prefix-minimality}
\end{equation}
Minimal Stinespring uniqueness \cite{Stinespring1955} therefore gives a unique
unitary $U_k:A_k\to A'_k$ with
\begin{equation}
 W'_k=(I_{O_k}\otimes U_k)W_k.
 \label{eq:ncg-prefix-unitary}
\end{equation}
Uniqueness follows directly from the spanning condition
\eqref{eq:ncg-prefix-minimality}, not from a choice of purification coordinates.

The sequential factorizations, with the natural tensor permutations understood,
are
\[
 W_k=(I_{O_{k-1}}\otimes V_k)(W_{k-1}\otimes I_{H_k^{\rm in}})
\]
and the analogous identity for $W'_k$.  Substituting
\eqref{eq:ncg-prefix-unitary} at cuts $k$ and $k-1$ shows that
\[
 V'_k(I_{H_k^{\rm in}}\otimes U_{k-1})
 -(I_{H_k^{\rm out}}\otimes U_k)V_k
\]
annihilates all vectors obtained by contracting the earlier output legs of
$W_{k-1}$, tensored with arbitrary vectors of $H_k^{\rm in}$.
These vectors span $H_k^{\rm in}\otimes A_{k-1}$ by
\eqref{eq:ncg-prefix-minimality}.  The displayed difference is therefore zero,
which proves \eqref{eq:ncg-minimal-unitary-covariance} on the whole memory domain.
For $k=1$ the same argument uses $A_0=\mathbb C$ and $U_0=I$.

The maps $X\mapsto U_kXU_k^*$ are unitaries on the Hilbert--Schmidt memory
spaces.  They intertwine the branch maps obtained from the fixed external
interventions, and hence also their Hilbert--Schmidt adjoints, mixed words and
word Grams.
\end{proof}

\begin{remark}[Scope of sequential uniqueness]
The uniqueness assertion concerns coherent isometric realizations that are minimal
at every prefix.  It is not a uniqueness theorem for arbitrary mixed-state,
measured-memory or environment-discarding implementations with a small physical
memory.  The canonical process representation below uses precisely the
prefix-minimal coherent class.
\end{remark}

For each typed boundary object $\xi$, put
$\mathscr X_{X,\xi}=\mathcal B(A_{X,\xi})$ with its Hilbert--Schmidt inner
product and let $\mathscr X_X=\bigoplus_\xi\mathscr X_{X,\xi}$.  If a primitive
branch $a:\xi\to\eta$ has reduced Schr\"odinger superoperator $T_{X,a}$, write
\[
 \rho_X(a)=P_\eta T_{X,a}P_\xi
\]
and adjoin a formal reverse $a^\sharp$ represented by
$\rho_X(a)^*_{\rm HS}$.  The reverse is an algebraic adjoint and need not be an
executable physical channel.  Let $\mathfrak F_X^{\rm op}$ be the finite-object
free path $*$-algebra on these primitive branches and their formal reverses.

\begin{theorem}
\label{thm:ncg-process-history-representation}
The represented process algebra
\begin{equation}
 \A_X^{\rm proc}
 :=\mathfrak F_X^{\rm op}/\Ker\rho_X
 \cong C^*\!\left(\rho_X(a):a\text{ primitive}\right)
 \subseteq\mathcal B(\mathscr X_X)
 \label{eq:ncg-process-algebra}
\end{equation}
is a finite-dimensional $C^*$-algebra.  Its represented $*$-isomorphism class is
independent of the tomographic frame, Kraus coordinates, purification coordinates,
and the chosen support-minimal sequential realization.
\end{theorem}

\begin{proof}
Typed concatenation is sent to block composition, object units to orthogonal
projections, and $\sharp$ to Hilbert--space adjoint, so $\rho_X$ is a
$*$-homomorphism and its kernel is a two-sided $*$-ideal.  The image is
finite-dimensional.  The unitary covariance of
Proposition~\ref{thm:ncg-minimal-comb-dilation} conjugates the complete word
representation, giving the asserted represented $*$-isomorphism.
\end{proof}

The standard dilation theorems ensure that finite CP branches can be implemented
on finite enlarged Hilbert spaces, but an operator corner is not automatically a
multiplicative representation of its parent algebra.  The following elementary
test, which will also motivate the separation between the operational input
and the algebraic essential-image theorem above, makes this distinction explicit
\cite{Stinespring1955,Holevo1982}.

\begin{proposition}
\label{prop:ncg-compression-defect}
Every finite family of finite-dimensional completely positive trace-nonincreasing
maps admits a simultaneous finite unitary completion after adjoining finite
ancillas and a finite control register.  Moreover, let
$\mathcal B\subseteq\mathcal B(H)$ be a finite represented unital $C^*$-algebra,
let $V:K\to H$ be an isometry, put $P=VV^*$ and $Q=I-P$, and set
$\kappa(b)=V^*bV$.  Then
\begin{equation}
 \kappa(ab)-\kappa(a)\kappa(b)=V^*aQbV,
 \qquad
 \kappa(a^*a)-\kappa(a)^*\kappa(a)=(QaV)^*(QaV)\succeq0.
 \label{eq:ncg-compression-square-defect}
\end{equation}
Consequently $\kappa$ is a $*$-homomorphism on a unital $*$-subalgebra
$\mathcal C\subseteq\mathcal B$ if and only if $P$ reduces $\mathcal C$.
\end{proposition}

\begin{proof}
Choose Stinespring contractions for the finitely many CP maps, dilate each
contraction to its Julia unitary, and take their finite controlled direct sum.
For the compression identity, insert $I=P+Q$ between $a$ and $b$.  The second
formula is the case $(a^*,a)$.  If $P$ reduces $\mathcal C$, the defect vanishes
and compression is multiplicative and $*$-preserving.  Conversely, if the
restriction is a $*$-homomorphism, applying the positive square defect to $c$
and $c^*$ gives $QcP=Qc^*P=0$ for every $c\in\mathcal C$, hence $[P,c]=0$.
\end{proof}

The process algebra removes coordinate ambiguity in the minimal sequential
dilation, but it can still contain algebraic directions that no admitted future
experiment detects.  The next quotient removes exactly those observationally null
directions.  Thus support minimality and future separation solve two different
minimality problems.

Let $\mathscr R_X$ denote the finite family of scalar readout functionals
generated by the declared future continuations and the minimal predictive record
reads.  Define
\begin{equation}
 \mathcal N_X^{\rm fut}
 =\{q\in\A_X^{\rm proc}:
 R(uqv)=R(uq^*v)=0\text{ for every }R\in\mathscr R_X
 \text{ and }u,v\in\A_X^{\rm proc}\}.
 \label{eq:ncg-contextual-null}
\end{equation}
If $\pi:\A_X^{\rm proc}\twoheadrightarrow B$ is a $C^*$-quotient through which
every $R\in\mathscr R_X$ factors as $R=\widetilde R\circ\pi$, call $B$
\emph{future separated} when
\begin{equation}
 \widetilde R(ubv)=\widetilde R(ub^*v)=0
 \quad\text{for all admitted }\widetilde R\text{ and }u,v\in B
 \quad\Longrightarrow\quad b=0.
 \label{eq:ncg-future-separated-definition}
\end{equation}

\begin{theorem}[Universal future-separated history algebra]
\label{thm:ncg-contextual-envelope}
The space $\mathcal N_X^{\rm fut}$ is a two-sided $*$-ideal and
\begin{equation}
 \A_X^{\mathrm{hist}}
 :=\A_X^{\rm proc}/\mathcal N_X^{\rm fut}
 \label{eq:ncg-history-algebra}
\end{equation}
is future separated and carries every admitted readout.  It is universal with these
properties: if a $C^*$-quotient $B$ of $\A_X^{\rm proc}$ carries every admitted
readout and is future separated in the sense of
\eqref{eq:ncg-future-separated-definition}, then the quotient maps identify
$B$ canonically with $\A_X^{\rm hist}$.
\end{theorem}

\begin{proof}
For $q\in\mathcal N_X^{\rm fut}$ and $c,d\in\A_X^{\rm proc}$,
$R(u(cqd)v)=R((uc)q(dv))=0$ in every admitted context; applying the same argument
to adjoints makes the space $*$-stable.  Hence it is a closed two-sided $*$-ideal.
Every admitted readout vanishes on $\mathcal N_X^{\rm fut}$ and therefore factors
through $\A_X^{\rm hist}$.  By construction, an element of the quotient annihilated
in all two-sided readout contexts has a representative in $\mathcal N_X^{\rm fut}$,
so the quotient is future separated.

Now let $\pi:\A_X^{\rm proc}\twoheadrightarrow B$ carry every admitted readout.  If
$q\in\Ker\pi$, then every admitted two-sided readout context vanishes on $q$ and
$q^*$, hence $\Ker\pi\subseteq\mathcal N_X^{\rm fut}$.  Thus there is a canonical
surjection $B\twoheadrightarrow\A_X^{\rm hist}$.  Conversely, if
$q\in\mathcal N_X^{\rm fut}$, then $\pi(q)$ is annihilated by every induced
two-sided readout context in $B$; future separation gives $\pi(q)=0$.  Hence
$\mathcal N_X^{\rm fut}\subseteq\Ker\pi$, the two kernels coincide, and the
canonical surjection is a $*$-isomorphism.
\end{proof}

On every finite-dimensional $C^*$-algebra there is a canonical normalized regular
trace.  For
$\A_X^{\mathrm{hist}}\cong\bigoplus_\lambda M_{n_\lambda}(\C)$, set
\begin{equation}
 \widehat\tau_{X,\mathrm{reg}}(a)
 =\frac{\Tr_{\C}(L_a)}{\dim_{\C}\A_X^{\mathrm{hist}}}
 =\frac{\sum_\lambda n_\lambda\Tr(a_\lambda)}
        {\sum_\lambda n_\lambda^2}.
\label{eq:ncg-grand-regular-trace}
\end{equation}
If $[w]$ denotes the represented class of a finite mixed history word, define
\begin{equation}
 \mathbb K_{X,r}(v,w)
 =\widehat\tau_{X,\mathrm{reg}}\bigl([v]^*[w]\bigr),
 \qquad |v|,|w|\le r.
\label{eq:ncg-word-gram}
\end{equation}
The regular trace is faithful, so the null space of $\mathbb K_{X,r}$ is exactly
the word-relation space at depth $r$.

\begin{theorem}[Finite flat reconstruction of the history algebra]
\label{thm:ncg-flat-history-reconstruction}
There is a least finite depth $r_X$ for which
\begin{equation}
 \rank\mathbb K_{X,r_X}=\rank\mathbb K_{X,r_X+1}.
 \label{eq:ncg-flat-rank}
\end{equation}
At this depth the represented words of length at most $r_X$ span all of
$\A_X^{\rm hist}$, and the panel through depth $r_X+1$ determines the unit,
multiplication, involution, and every represented mixed history.  More generally,
if two finite source-generated history realizations have equal depth-$(r+1)$ word
Grams and are both flat at depth $r$, then there is a unique unital
$*$-isomorphism between them carrying every source word to its namesake.
\end{theorem}

\begin{proof}
The word spans form an increasing sequence of subspaces of the finite-dimensional
algebra, so their dimensions stabilize.  At the first rank equality the depth-$r$
and depth-$(r+1)$ spans coincide.  Left multiplication by every primitive source
preserves the common span; since it contains the unit, it contains every generated
word and hence the whole algebra.  The additional word layer records left
multiplication by each primitive source on a spanning family, while reversing word order
and taking formal adjoints recovers the involution.  For two realizations with the same flat
Gram panel, the map sending one represented word to the namesake word in the
second is well defined and isometric because the Gram kernels have the same null
space and inner products.  Flatness makes it onto; the one-step word identities
show that it preserves multiplication and involution.  Uniqueness follows because
the saturated words span.
\end{proof}

\begin{proposition}
\label{prop:ncg-grand-tensor-reconstruction}
Under the finite operational assumptions above, the complete intervention
statistics determine, up to their natural isomorphisms,
\begin{enumerate}[label=\textup{(\roman*)},leftmargin=2.8em]
\item the minimal predictive state space $Z_X^{\min}$, its branch probabilities, and its partial update maps;
\item the support-minimal process representation and the universal future-separated
      history algebra $\A_X^{\rm hist}$;
\item the normalized regular trace and the positive word-Gram hierarchy;
\item a finite flat depth $r_X$ at which the history algebra, multiplication and
      involution are reconstructible from the word panel.
\end{enumerate}
Thus the finite history algebra and its source representation are reconstructed
from the operational table without an independently specified hidden-state or
Hilbert-space realization.
\end{proposition}

\begin{proof}
The predictive quotient and its universal property were proved above.
Proposition~\ref{thm:ncg-comb-tomography} reconstructs the positive causal process;
Proposition~\ref{thm:ncg-minimal-comb-dilation} fixes its support-minimal sequential
class; Theorems~\ref{thm:ncg-process-history-representation} and
\ref{thm:ncg-contextual-envelope} give the canonical future-separated history
algebra; and Theorem~\ref{thm:ncg-flat-history-reconstruction} proves finite
saturation and reconstruction of its algebraic operations.  The normalized regular
trace is intrinsic to the resulting finite-dimensional $C^*$-algebra and therefore
so are the word Grams.  These constructions are invariant under the natural
unitary and represented $*$-isomorphisms, proving the claim.
\end{proof}

The proposition isolates the finite operational hypotheses used for the predictive
application.  The spectral reconstruction in Section~\ref{sec:finite} depends only on the resulting
finite $C^*$-algebra, trace, and differential datum.

\subsection{Realization and recognition with fixed boundary contexts}
\label{subsec:operational-recognition-detail}

The physical memory-dimension minimality used in this subsection is distinct from
the coherent prefix-minimality used for the canonical process realization above.
The two notions serve different purposes and are not identified.

The preceding reconstruction starts from a given finite intervention table.  We
now prove a converse realization statement for the represented algebra.  The
category is deliberately boundary-relative: preparations and terminal readouts are
external contexts and are not adjoined as internal generators.  This convention
is substantive, as shown below.  Measure-and-prepare maps belong to the standard
theory of entanglement-breaking channels \cite{HorodeckiShorRuskai2003}.

\begin{theorem}[Finite operational realization with multiplicities]
\label{thm:ncg-operational-realization}
Let $\pi:\A\to\mathcal B(K)$ be a faithful unital representation of a
finite-dimensional $C^*$-algebra.  Then there are a finite physical memory $H$,
a finite family of normalized completely positive instruments, boundary
preparations, and terminal readouts such that the reducing source representation of
the internal history algebra is unitarily equivalent to $(\A,K,\pi)$.  In
particular all representation multiplicities are preserved.
\end{theorem}

\begin{proof}
After a unitary change of coordinates write
\[
 K=\bigoplus_{\lambda=1}^s\C^{n_\lambda}\otimes\C^{m_\lambda},
 \qquad
 \pi(\A)=\bigoplus_{\lambda=1}^sM_{n_\lambda}(\C)\otimes I_{m_\lambda}.
\]
Take $H$ with orthonormal basis $e_{\lambda,a,r}$ and put
$p_{\lambda,a,r}=|e_{\lambda,a,r}\rangle\langle e_{\lambda,a,r}|$.  For
$1\le a,b\le n_\lambda$ define
\[
 L_{ab,r}^\lambda
 =|e_{\lambda,a,r}\rangle\langle e_{\lambda,b,r}|,
 \qquad
 \Phi_{ab}^\lambda(X)
 =\sum_{r=1}^{m_\lambda}L_{ab,r}^\lambda X(L_{ab,r}^\lambda)^*.
\]
Each map is CP and trace nonincreasing, and direct multiplication of the
rank-one Kraus operators gives
\[
 \Phi_{ab}^\lambda\Phi_{cd}^\mu
 =\delta_{\lambda\mu}\delta_{bc}\Phi_{ad}^\lambda,
 \qquad
 (\Phi_{ab}^\lambda)^*_{\rm HS}=\Phi_{ba}^\lambda.
\]
Let $\mathsf P(X)=\sum_{\lambda,a,r}p_{\lambda,a,r}Xp_{\lambda,a,r}$ be the
dephasing channel.  For each success branch use the complementary failure branch
\[
 \Psi_b^\lambda=\mathsf P-\Phi_{bb}^\lambda.
\]
The pair $(\Phi_{ab}^\lambda,\Psi_b^\lambda)$ is a normalized two-outcome
instrument, and adjoining the failure branches adds no generator beyond the
matrix-unit algebra and the identity.  The diagonal Hilbert--Schmidt subspace
\[
 \mathcal V=\operatorname{span}\{p_{\lambda,a,r}\}
\]
is reducing for every branch and every Hilbert--Schmidt adjoint.  Under the
unitary $e_{\lambda,a,r}\mapsto p_{\lambda,a,r}$, the restricted matrix units
act as
$E_{ab}^{(\lambda)}\otimes I_{m_\lambda}$.  Boundary preparations
$p_{\lambda,a,r}$ and the corresponding terminal effects separate all source
matrix coefficients, so the future-null ideal vanishes on this source
representation.  Hence the future-separated internal history representation is
exactly $\pi(\A)$ with the stated multiplicities.
\end{proof}

\begin{lemma}[Rank-one forcing from a deterministic transition panel]
\label{lem:ncg-rank-one-forcing}
Let $\rho_1,\ldots,\rho_N$ be density matrices on a finite Hilbert space $H$ that
are perfectly distinguished by effects $E_1,\ldots,E_N$, so that
$\Tr(E_i\rho_j)=\delta_{ij}$.  Then
\begin{equation}
 \sum_{j=1}^N\rank\rho_j\le\dim H.
 \label{eq:perfect-discrimination-rank}
\end{equation}
In particular, if $\dim H=N$, the $\rho_j$ are rank-one states with mutually
orthogonal supports.  Let $\Phi$ be completely positive and suppose that, on these
rank-one states, a retained one-step panel prescribes
\begin{equation}
 \Phi(\rho_j)=w_j\rho_{\sigma(j)},\qquad w_j>0,
 \label{eq:deterministic-panel-map}
\end{equation}
for $j$ in a domain $J$, while $\Tr\Phi(\rho_j)=0$ for $j\notin J$.  If
$\sigma:J\to\{1,\ldots,N\}$ is injective, then the diagonal Hilbert--Schmidt span
$\mathcal V=\Span\{\rho_j\}$ is invariant under both $\Phi$ and
$\Phi^*_{\rm HS}$.
\end{lemma}

\begin{proof}
Perfect discrimination implies that the supports of the $\rho_j$ are mutually
orthogonal: $\Tr(E_i\rho_i)=1$ and $0\le E_i\le I$ force $E_i=I$ on
$\supp\rho_i$, whereas $\Tr(E_i\rho_j)=0$ forces $E_i=0$ on
$\supp\rho_j$ for $j\ne i$.  The rank bound follows.  At equality
$\dim H=N$, every support is one dimensional; write
$\rho_j=|e_j\rangle\langle e_j|$.

Choose Kraus operators $K_\nu$ for $\Phi$.  If $j\notin J$, then
$0=\Tr\Phi(\rho_j)=\sum_\nu\|K_\nu e_j\|^2$, hence every
$K_\nu e_j$ vanishes.  If $j\in J$, the positive operator
$\sum_\nu|K_\nu e_j\rangle\langle K_\nu e_j|$ has one-dimensional support
$\mathbb C e_{\sigma(j)}$, so every $K_\nu e_j$ lies in that line.  Thus
$\Phi$ preserves $\mathcal V$.  For the Hilbert--Schmidt adjoint,
\[
 \langle e_k,\Phi^*(\rho_i)e_j\rangle
 =\sum_\nu
   \overline{\langle e_i,K_\nu e_k\rangle}
   \langle e_i,K_\nu e_j\rangle.
\]
A nonzero term requires simultaneously $\sigma(k)=i=\sigma(j)$; injectivity of
$\sigma$ therefore forces $j=k$.  Hence $\Phi^*(\rho_i)$ is diagonal in the same
basis and $\mathcal V$ is reducing.
\end{proof}

\begin{proposition}[Finite recognition from the one-step panel]
\label{prop:ncg-operational-recognition}
For the realization in Theorem~\ref{thm:ncg-operational-realization}, retain the
finite preparation--branch--readout panel consisting of the perfect boundary
discrimination data and the one-step success and complementary failure
transitions, including the zero probabilities outside each branch domain.  Any realization of this
panel whose total physical memory dimension is minimal has a future-separated
source representation unitarily equivalent to $(\A,K,\pi)$.
\end{proposition}

\begin{proof}
There are $N=\dim K$ retained source preparations.  By
Lemma~\ref{lem:ncg-rank-one-forcing}, any realization needs physical memory
dimension at least $N$.  In a dimension-minimal realization equality holds, so the
preparations are rank-one states
$p_i=|e_i\rangle\langle e_i|$ on an orthonormal basis, uniquely up to a unitary and
phase choices that disappear at the level of the projections.

Fix the retained success branch indexed by $(\lambda,a,b)$.  Its domain
contains all $(\lambda,b,r)$, $1\le r\le m_\lambda$, with successor
$(\lambda,a,r)$; the panel records zero probabilities outside this domain.
The deterministic outputs and the zero probabilities put it in the situation of
Lemma~\ref{lem:ncg-rank-one-forcing}.  For a fixed branch the successor rule is a
partial injection, so the diagonal source space
\[
 \mathcal V=\Span\{p_{\lambda,a,r}\}
\]
is reducing for that branch and its Hilbert--Schmidt adjoint.  On $\mathcal V$ the
retained transition probabilities determine the partial matrix unit
$E_{ab}^{(\lambda)}\otimes I_{m_\lambda}$.  The complementary failure branches
have a partial-identity successor rule on the remaining basis states, so the same
lemma makes $\mathcal V$ reducing for them as well.  Their restrictions already
belong to the generated block algebra.  Composition of the retained one-step
branches therefore recovers
\[
 \bigoplus_\lambda M_{n_\lambda}(\mathbb C)\otimes I_{m_\lambda}
\]
with exactly the prescribed multiplicities.

Finally, the retained terminal readouts perfectly distinguish the basis projections
and, after composition with the matrix-unit branches, separate every matrix
coefficient of this represented algebra.  An element annihilated by all two-sided
retained readout contexts is therefore zero on the source representation, so the
contextual future-null ideal vanishes there.  The resulting future-separated source
representation is consequently unitarily equivalent to $(\A,K,\pi)$.
\end{proof}

\begin{proposition}
\label{prop:ncg-boundary-enlargement}
Let $N=\dim K\ge2$ and let $p_1,\ldots,p_N$ be the distinguished source
preparations.  If the reset maps
\[
 \mathcal R_u(X)=\Tr(X)p_u
\]
are adjoined to the internal algebra, then on the source span their products with
Hilbert--Schmidt adjoints generate $M_N(\C)$.  Thus a proper block algebra
$\bigoplus_\lambda M_{n_\lambda}(\C)\otimes I_{m_\lambda}$ can be destroyed by
promoting boundary preparations to internal generators.
\end{proposition}

\begin{proof}
On the diagonal source basis the reset $\mathcal R_u$ is the rank-one map
$|e_u\rangle\langle\mathbf 1|$.  Hence
$\mathcal R_u\mathcal R_v^*=N|e_u\rangle\langle e_v|$.  These are all matrix
units of $M_N(\C)$.
\end{proof}


\begin{theorem}
\label{thm:ncg-enriched-closure}
Let $\rho_1,\ldots,\rho_r$ be linearly independent reproducible conditional
memory states and put $Lx=\sum_jx_j\rho_j$, $G=L^*L$.  Let $R$ be the retained
boundary response map and set
\[
 H_\alpha=L^*\Phi_\alpha L,
 \quad K_\alpha=(\Phi_\alpha L)^*(\Phi_\alpha L),
 \quad B=RL,
 \quad C_\alpha=R\Phi_\alpha L.
\]
Assume $B$ has full row rank.  If, for every branch $\alpha$,
\[
 K_\alpha-H_\alpha^*G^{-1}H_\alpha=0,
 \qquad
 C_\alpha(I-B^\dagger B)=0,
\]
then $\mathcal V=\Ran L$ is forward invariant and, with
$A_\alpha=C_\alpha B^\dagger$,
\[
 R\Phi_\alpha X=A_\alpha RX\quad(X\in\mathcal V),
 \qquad
 R\Phi_wL=A_wB
\]
for every finite word $w$.  Hence finitely many successor-state Gram and response
measurements certify the full forward-response quotient.  These enriched overlap
measurements are additional resources and are not inferred from the original
boundary probability table.
\end{theorem}

\begin{proof}
The orthogonal projection onto $\mathcal V$ is $P_\mathcal V=LG^{-1}L^*$, so
\[
 K_\alpha-H_\alpha^*G^{-1}H_\alpha
 =((I-P_\mathcal V)\Phi_\alpha L)^*((I-P_\mathcal V)\Phi_\alpha L).
\]
The first condition is therefore equivalent to forward invariance and gives
$\Phi_\alpha L=LG^{-1}H_\alpha$.  The second says precisely that $C_\alpha$
annihilates $\ker B$, hence factors uniquely through the response coordinates:
$C_\alpha=A_\alpha B$.  Combining these identities and iterating proves the
word formula.  The overlaps entering $G,H_\alpha,K_\alpha$ can, when the whole
memory is experimentally accessible, be measured on two copies using the swap
identity $\Tr(\mathsf F(X\otimes Y))=\Tr(XY)$ \cite{EkertEtAl2002}.
\end{proof}

\begin{corollary}
\label{cor:ncg-operational-calibrated-triples}
For every finite marked spectral triple, its represented algebra component admits
the operational realization of Theorem~\ref{thm:ncg-operational-realization}.
After independently supplying its commutator derivation, declared marks, and the
sharp real-linear commutant coordinates of Theorem~\ref{thm:commutant-fibre}, the full
compatible Dirac operator is reconstructed.  No claim is made that the ordinary
operational probabilities alone determine the differential datum or its scale.
\end{corollary}

\section{Positive Stieltjes memory and static-preserving cutoffs}
\label{app:memory-details}

This section studies the inverse data lost by a static reduction.
Eliminating auxiliary degrees of
freedom produces a resolvent self-energy, whereas its value at zero retains
only a static Schur correction.  The distinction between these two data is
described using classical Naimark dilation and matrix-valued
Herglotz--Stieltjes theory
\cite{Holevo1982,GesztesyTsekanovskii2000,KreinNudelman1977}.
We record the minimal realization relative to the effective space, moment
determinacy and indeterminacy, static-preserving cutoffs, and high-energy
contact limits.  These results distinguish the full dynamical memory from
the static correction retained by the effective operator.

\subsection{Positive measures, self-energies, and inverse moments}
\label{subsec:dynamic-memory}

\begin{definition}[Positive operator-valued tail measure]
\label{def:supp-POVM}
Let $H$ be finite dimensional.  A bounded positive operator-valued tail
measure is a weakly countably additive map
\[
 \mu:\mathfrak B((0,\infty))\to\mathcal B(H),
 \qquad
 \mu(E)\succeq0,
 \qquad
 \|\mu((0,\infty))\|<\infty.
\]
Its Stieltjes self-energy is
\begin{equation}
 {
 \Sigma_\mu(z)
 =\int_{(0,\infty)}
  \frac{1}{\lambda-z}\,\dd\mu(\lambda),
 \qquad z\in\mathbb C\setminus[0,\infty).}
 \label{eq:supp-Stieltjes}
\end{equation}
\end{definition}

By the standard Naimark dilation and matrix-valued Herglotz--Stieltjes realization
theory \cite{Holevo1982,GesztesyTsekanovskii2000,KreinNudelman1977}, every such
measure has a source-minimal realization
\begin{equation}
 \Sigma_\mu(z)=B_\mu(C_\mu-z)^{-1}B_\mu^*,
 \qquad C_\mu\ge0,
 \label{eq:supp-Naimark-realization}
\end{equation}
unique up to a unitary intertwining $C_\mu$ and $B_\mu^*$.  We use this classical
realization theorem as an input rather than reproving it.  The inverse-fibre questions
below concern which features of $\mu$ survive passage to the static self-energy and how the
remaining dynamical coordinate behaves under coarse graining.

\begin{theorem}[Static self-energy, inverse moments, and positive Euclidean memory]
\label{thm:supp-dynamic-jets}
For $\operatorname{Im}z>0$,
\begin{equation}
 \operatorname{Im}\Sigma_\mu(z)
 =(\operatorname{Im}z)
  \int\frac{1}{|\lambda-z|^2}\,\dd\mu(\lambda)
 \succeq0.
 \label{eq:supp-Herglotz-sign}
\end{equation}
The finite static self-energy
\[
 Q_0=\lim_{s\downarrow0}\Sigma_\mu(-s)
\]
exists exactly when the inverse first moment is bounded, and then
\begin{equation}
 {Q_0=\int\lambda^{-1}\,\dd\mu(\lambda).}
 \label{eq:supp-static-inverse-moment}
\end{equation}
If the corresponding inverse moment is finite, the $r$-th boundary jet is
\begin{equation}
 {
 \Sigma_{\mu,-}^{(r)}(0)
 =r!\int\lambda^{-(r+1)}\,\dd\mu(\lambda).}
 \label{eq:supp-dynamic-jet}
\end{equation}
The Euclidean memory kernel
\begin{equation}
 {
 K_\mu(t)=\int e^{-t\lambda}\,\dd\mu(\lambda)
 =B_\mu e^{-tC_\mu}B_\mu^*}
 \label{eq:supp-Euclidean-memory}
\end{equation}
is positive and
\[
 \Sigma_\mu(-s)
 =\int_0^\infty e^{-st}K_\mu(t)\,\dd t.
\]
\end{theorem}

\begin{proof}
The scalar identity
$\operatorname{Im}(\lambda-z)^{-1}
=(\operatorname{Im}z)|\lambda-z|^{-2}$ gives the sign after positive
operator integration.  Monotone convergence along $z=-s$ gives
\eqref{eq:supp-static-inverse-moment}.  Differentiating
$(\lambda-z)^{-1}$ gives the higher inverse moments.  The spectral calculus
gives the second equality in \eqref{eq:supp-Euclidean-memory}; Tonelli's
theorem applied to scalarizations gives the Laplace relation.
\end{proof}

\begin{remark}[Positive memory is not full OS positivity]
Positivity of $K_\mu(t)$ is a two-point positive memory statement.  It does not
by itself supply a Euclidean field law, a positive-time product algebra, or
Osterwalder--Schrader positivity for reflected words.
\end{remark}

The Stieltjes transform is the full dynamical coordinate.  Static data or
finitely many moments need not determine it.  The next subsection gives
sufficient determinacy hypotheses and an example in which even all positive
and inverse integer moments fail to determine the self-energy.

\subsection{Moment determinacy and indeterminacy}

\begin{theorem}
\label{thm:supp-memory-determinacy}
If $\supp\mu$ is compact, the complete nonnegative moment sequence
$M_n=\int\lambda^n\,\dd\mu(\lambda)$ determines $\mu$ uniquely.  The same
conclusion holds on unbounded support if
$\int e^{a\lambda}\,\dd\Tr\mu(\lambda)<\infty$ for some $a>0$.
\end{theorem}

\begin{proof}
On compact support, equality of all moments gives equality against all
polynomials; Stone--Weierstrass gives equality against every continuous
function, and polarization gives equality of the operator measures.

Under an exponential moment, each scalarized Laplace transform
$\int e^{-t\lambda}\,\dd\langle\xi,\mu\eta\rangle$ extends holomorphically to
a half-plane containing the origin.  Its Taylor coefficients are the moments.
Equality of all moments therefore gives equality of the Laplace transforms
near zero and hence throughout the connected half-plane.  Uniqueness of the
Laplace transform and polarization recover $\mu$.
\end{proof}

\begin{remark}[Scope of the memory classification]
The results in this subsection concern positive operator-valued Stieltjes memory.
Indefinite generalized-Nevanlinna or Pontryagin-space realizations require
additional realization theory and are outside the scope of this section; the finite rational extension is developed separately in Section~\ref{app:rational-krein}.
\end{remark}

\begin{proposition}
\label{cth:supp-moment-indeterminate}
For $|\varepsilon|\le1$ define
\begin{equation}
 {
 \dd\mu_\varepsilon(x)
 =
 \left[1+\varepsilon\sin(2\pi\log x)\right]
 \frac{e^{-(\log x)^2/2}}{x\sqrt{2\pi}}\,\dd x.}
 \label{eq:supp-lognormal-family}
\end{equation}
Then for every integer $n\in\mathbb Z$,
\begin{equation}
 {
 \int_0^\infty x^n\,\dd\mu_\varepsilon(x)=e^{n^2/2}}
 \label{eq:supp-lognormal-moments}
\end{equation}
independently of $\varepsilon$.  Nevertheless
$\varepsilon\ne\varepsilon'$ implies that the Stieltjes transforms differ for
some, hence an open set of, $z\in\mathbb C\setminus[0,\infty)$.  Consequently
all positive moments, all inverse integer moments, and every finite Hankel
panel can agree while the dynamic self-energies are inequivalent.
\end{proposition}

\begin{proof}
Set $y=\log x$.  The perturbation of the $n$-th moment is
\[
 \frac1{\sqrt{2\pi}}
 \int_\mathbb R
 e^{ny-y^2/2}\sin(2\pi y)\,\dd y
 =
 \operatorname{Im}
 \exp\!\left(\frac{(n+2\pi i)^2}{2}\right)
 =
 e^{(n^2-4\pi^2)/2}\sin(2\pi n)=0
\]
for every $n\in\mathbb Z$.  The unperturbed lognormal moment is
$e^{n^2/2}$.  Distinct $\varepsilon$ give distinct measures.

If a finite signed measure on $[0,\infty)$ had identically zero Stieltjes
transform on the negative real axis, the identity
\[
 \int\frac{1}{x+s}\,\dd\sigma(x)
 =
 \int_0^\infty e^{-st}
 \left(\int e^{-tx}\,\dd\sigma(x)\right)\dd t
\]
and two applications of uniqueness of the Laplace transform would force
$\sigma=0$.  Hence distinct measures have distinct Stieltjes transforms.
\end{proof}

\subsection{Memory cutoff and composition}

A cutoff removes high-energy memory while incorporating its static
contribution into the effective operator.  The following identities
separate exact preservation at zero energy from the error at nonzero
spectral parameter.

For an effective operator $A$ and positive memory measure $\mu$, write
\begin{equation}
F_{A,\mu}(z)=A-z-\Sigma_\mu(z),
\qquad
Q_{>\Lambda}=\int_{(\Lambda,\infty)}\lambda^{-1}\,\dd\mu(\lambda),
\label{eq:supp-memory-head}
\end{equation}
and define
\begin{equation}
\RG_\Lambda(A,\mu)
=\left(A-Q_{>\Lambda},\mu|_{(0,\Lambda]}\right).
\label{eq:supp-memory-RG}
\end{equation}

\begin{theorem}[Static-preserving memory cutoffs]
\label{thm:supp-memory-RG}
Assume the static inverse moment
$Q_0=\int\lambda^{-1}\,\dd\mu(\lambda)$ is finite, and interpret $F_{A,\mu}(0)$
as the limit from the negative real axis.  If
$(A_\Lambda,\mu_\Lambda)=\RG_\Lambda(A,\mu)$, then
\begin{equation}
{F_{A_\Lambda,\mu_\Lambda}(0)=F_{A,\mu}(0).}
\label{eq:supp-memory-static-preserved}
\end{equation}
For $z\in\mathbb C\setminus[0,\infty)$ with $|z|\le r<\Lambda$,
\begin{equation}
{
\|F_{A,\mu}(z)-F_{A_\Lambda,\mu_\Lambda}(z)\|
\le\frac r{\Lambda-r}\|Q_{>\Lambda}\|.}
\label{eq:memory-cutoff-error}
\end{equation}
Moreover, if $0<\Lambda_2\le\Lambda_1$,
\begin{equation}
\RG_{\Lambda_2}\circ\RG_{\Lambda_1}=\RG_{\Lambda_2}.
\label{eq:supp-memory-semigroup}
\end{equation}
\end{theorem}

\begin{proof}
Split the Stieltjes integral at $\Lambda$.  At $z=0$ the removed high-energy part is
exactly $Q_{>\Lambda}$ and is canceled by the effective local correction.  Away from zero,
\[
\frac1{\lambda-z}-\frac1\lambda
=\frac z{\lambda(\lambda-z)},
\]
whose norm is at most
$r[(\Lambda-r)\lambda]^{-1}$ for $\lambda>\Lambda$.  Integrating gives
\eqref{eq:memory-cutoff-error}.  Two successive cutoffs simply split the same
measure into disjoint shells; their inverse-moment local corrections add, proving the
semigroup law.
\end{proof}

\subsection{High-energy escape and contact limits}

The same static correction can also survive when the memory measure escapes
to infinite energy.  The limiting object then contains an instantaneous
term, as the following two results show.

\begin{theorem}[High-energy memory escape and instantaneous contact terms]
\label{thm:supp-contact-escape}
Let $\mu_n$ be positive tail measures supported in $[m,\infty)$ for one
$m>0$ and define
\[
 \dd\nu_n(\lambda)=\lambda^{-1}\dd\mu_n(\lambda),
 \qquad
 Q_n=\nu_n([m,\infty))=\Sigma_n(0).
\]
If $\sup_n\Tr Q_n<\infty$, every sequence has a subsequence and a positive
operator-valued measure on the one-point compactification of the form
\[
 \nu_\infty=\nu_{\rm fin}+Q_{\rm esc}\delta_\infty,
 \qquad Q_{\rm esc}\succeq0,
\]
such that, locally uniformly off $[m,\infty)$,
\begin{equation}
 {
 \Sigma_n(z)\longrightarrow
 Q_{\rm esc}
 +\int_{[m,\infty)}
  \frac{\lambda}{\lambda-z}\,\dd\nu_{\rm fin}(\lambda).}
 \label{eq:supp-contact-selfenergy}
\end{equation}
The Euclidean memory measures satisfy
\begin{equation}
 {
 K_n(t)\,\dd t
 \longrightarrow
 Q_{\rm esc}\delta_0+K_{\rm fin}(t)\,\dd t.}
 \label{eq:supp-contact-kernel}
\end{equation}
\end{theorem}

\begin{proof}
The matrix entries of $\nu_n$ are uniformly bounded finite measures on the
compact space $[m,\infty]$.  Weak-* compactness and polarization give a
positive operator-valued limit, which decomposes into finite and point-at-
infinity parts.  For $z\notin[m,\infty)$,
$g_z(\lambda)=\lambda/(\lambda-z)$ extends continuously to infinity with
$g_z(\infty)=1$, proving \eqref{eq:supp-contact-selfenergy}.
For $\varphi\in C_c([0,\infty))$,
\[
 h_\varphi(\lambda)
 =\int_0^\infty\varphi(t)\lambda e^{-\lambda t}\,\dd t
\]
extends continuously with $h_\varphi(\infty)=\varphi(0)$; integration against
$\nu_n$ gives \eqref{eq:supp-contact-kernel}.
\end{proof}

\begin{proposition}
\label{cth:supp-contact-static}
For $Q\succeq0$ and $\lambda_n\to\infty$, let
\[
 \mu_n=\lambda_nQ\,\delta_{\lambda_n}.
\]
Then
\begin{equation}
 {
 \Sigma_n(0)=Q,\qquad
 \Sigma_n(z)=\frac{\lambda_n}{\lambda_n-z}Q\longrightarrow Q,}
 \label{eq:supp-contact-static}
\end{equation}
whereas $\lambda_ne^{-\lambda_nt}Q\to0$ for every $t>0$ and converges
distributionally to $Q\delta_0$.  The same static local correction is
therefore compatible with an intrinsic instantaneous term or with memory escaping to
infinite energy.
\end{proposition}

\begin{proof}
The resolvent identities are immediate, and the exponential kernels form the
standard positive approximate identity at the origin.
\end{proof}

\section{Marked monodromy and componentwise coarse graining}
\label{app:marked-monodromy}

Locally isomorphic spectral triples need not admit a compatible global
trivialization.  This section records the standard graph-holonomy
description of that ambiguity and the condition for descent under graph
contraction.  It also distinguishes the holonomy of a full marked triple
from the information detected by its one-dimensional characters.

\subsection{Flat bundles of marked spectral triples}

For a marked spectral triple $\mathsf S$ in the sense of
Definition~\ref{def:marked-spectral-triple}, let
\begin{equation}
 G_{\mathsf S}=\operatorname{Aut}(\mathsf S)
 \label{eq:supp-marked-aut}
\end{equation}
be the compact group of unitary automorphisms of the marked triple preserving all declared
marks and the chosen algebra automorphism convention.  In this subsection the
auxiliary marks may include orientation, modular, calibration, frame, or
determinant-line data.

\begin{proposition}
\label{thm:supp-marked-bundle}
Let $\Gamma$ be connected, let each vertex triple be marked-isomorphic to
$\mathsf S$, and let each oriented edge carry a marked unitary isomorphism
$U_e$ with $U_{\bar e}=U_e^{-1}$.  Then:
\begin{enumerate}[label=\textup{(B\arabic*)},leftmargin=3em]
\item after choosing a base vertex and spanning tree, chord holonomies define a
      homomorphism
      \begin{equation}
       \rho_U:\pi_1(\Gamma)\longrightarrow G_{\mathsf S};
       \label{eq:supp-monodromy}
      \end{equation}
\item changing vertex frames conjugates $\rho_U$ by one element of
      $G_{\mathsf S}$;
\item a global marked trivialization compatible with all edge identifications exists if and only if $\rho_U$ is
      trivial;
\item in general the isomorphism class of the flat marked spectral bundle is
      \begin{equation}
       {
       [\rho_U]\in
       \operatorname{Hom}(\pi_1(\Gamma),G_{\mathsf S})/G_{\mathsf S}.}
       \label{eq:supp-monodromy-class}
      \end{equation}
\end{enumerate}
Since $\pi_1(\Gamma)$ is free of rank $b_1(\Gamma)$, the bundle is determined
by $b_1(\Gamma)$ elements of $G_{\mathsf S}$ modulo simultaneous conjugation.
\end{proposition}

\begin{proof}
Use the edge maps along a spanning tree to identify every vertex triple with
the base triple.  Each chord gives the ordered product around its fundamental
cycle, an element of $G_{\mathsf S}$.  These cycle values extend uniquely to a
homomorphism from the free fundamental group.  Changing the base frame
conjugates every cycle value by the same automorphism.  Trivial cycle values
make the tree identification path independent and hence give one global marked
frame.  Conversely, a global frame makes every cycle product telescope.  This
is precisely the clutching classification
\eqref{eq:supp-monodromy-class} on a graph.
\end{proof}

\begin{corollary}
\label{cor:supp-twisted-bundle}
If the local triples are marked-isomorphic and the edge identifications satisfy the declared compatibility relations, nontrivial $[\rho_U]$ obstructs a global marked trivialization compatible with those identifications.
\end{corollary}

\begin{theorem}
\label{thm:character-holonomy-detection}
Let
$\chi_1,\ldots,\chi_m:G_{\mathsf S}\to U(1)$ be determinant, charge, or other
complex line characters and let
$\epsilon_1,\ldots,\epsilon_r:G_{\mathsf S}\to\{\pm1\}$ be real Pfaffian-sign
characters.  Trivial frame holonomy implies triviality of all induced line
holonomies.  The converse holds for all such graph bundles exactly when
\begin{equation}
 {
 \bigcap_j\Ker\chi_j\cap\bigcap_k\Ker\epsilon_k=\{1\}.}
 \label{eq:supp-character-faithfulness}
\end{equation}
\end{theorem}

\begin{proof}
Every character sends the identity to the identity.  If
\eqref{eq:supp-character-faithfulness} holds and all character holonomies are
trivial, every cycle value of $\rho_U$ lies in the displayed intersection and
is therefore the identity.  Conversely, if the intersection contains
$g\ne1$, assign $g$ as monodromy on a one-cycle graph.  Every listed character
is trivial although the marked bundle is not.
\end{proof}

\begin{proposition}
\label{cth:supp-det-no-frame}
Let a marked spectral triple have scalar algebra, zero Dirac operator, and a
two-dimensional multiplicity space so that its marked automorphism group contains
$SU(2)$.  On a one-cycle graph choose monodromy
\begin{equation}
 g=\begin{pmatrix}e^{i\theta}&0\\0&e^{-i\theta}\end{pmatrix},
 \qquad \theta\notin2\pi\Z.
 \label{eq:supp-SU2-monodromy}
\end{equation}
Then the determinant-line holonomy is one although the marked Hilbert bundle
has nontrivial monodromy.
\end{proposition}

\begin{proof}
Every $SU(2)$ element has determinant one, while the displayed matrix is not
the identity.  Apply \Cref{thm:supp-marked-bundle,thm:character-holonomy-detection}.
\end{proof}

\subsection{Descent under graph contraction}

Let $\rho:\pi_1(\Gamma,v_0)\to G_{\mathsf S}$ be marked monodromy and let a graph
contraction $q:\Gamma\to\bar\Gamma$ induce a surjection $q_*$ on fundamental
groups.

\begin{theorem}[Exact marked-monodromy descent]
\label{thm:supp-monodromy-descent}
There is a unique $\bar\rho:\pi_1(\bar\Gamma)\to G_{\mathsf S}$ satisfying
\begin{equation}
\rho=\bar\rho\circ q_*
\label{eq:supp-monodromy-factor}
\end{equation}
if and only if
\begin{equation}
{\Ker q_*\subseteq\Ker\rho.}
\label{eq:supp-monodromy-kernel}
\end{equation}
The conjugacy class $[\bar\rho]$ is gauge independent and functorial under further
contractions.  A determinant or Pfaffian character may satisfy its own descended
kernel condition even when \eqref{eq:supp-monodromy-kernel} fails, so line descent
does not imply descent of the full marked spectral triple.
\end{theorem}

\begin{proof}
If \eqref{eq:supp-monodromy-factor} holds, every element of $\Ker q_*$ lies in
$\Ker\rho$.  Conversely, if \eqref{eq:supp-monodromy-kernel} holds, define
$\bar\rho(q_*\gamma)=\rho(\gamma)$.  The kernel condition makes the definition
independent of the chosen lift.  Surjectivity of $q_*$ gives uniqueness and the
quotient-group universal property gives functoriality.  The character statement
follows by applying the same argument to a nonfaithful one-dimensional
representation of $G_{\mathsf S}$.
\end{proof}

\subsection{Componentwise coarse graining}

Current aggregation, memory cutoff, and monodromy descent act on different
kinds of data.  The following corollary combines their respective
composition laws while retaining the separate admissibility hypotheses
for currents, positive memory and marked holonomy.

\begin{corollary}[Componentwise coarse graining of the combined inverse data]
\label{thm:supp-fibre-RG}
Collect
\begin{equation}
\mathfrak U(G,c;A,\mu;\rho)
=\bigl(\mathcal J(G,c),(A,\mu),[\rho]\bigr).
\label{eq:combined-inverse-data}
\end{equation}
For a connected-fibre graph quotient, a memory threshold $\Lambda$, a finite static
inverse moment as in Theorem~\ref{thm:supp-memory-RG}, and a compatible isometric
effective-space compression $V$, define
\begin{equation}
{
\mathfrak U\longmapsto
\left(
Q_E\mathcal J(G,c),
\bigl(V^*(A-Q_{>\Lambda})V,
      V^*\mu|_{(0,\Lambda]}V\bigr),
[\bar\rho]
\right).}
\label{eq:supp-U-map}
\end{equation}
When the monodromy kernel condition holds, the three components are canonical.
For composable graph quotients and nested memory thresholds they compose in the
corresponding componentwise sense.  For effective-space compression we require
nested isometries: if $V_1:H_1\to H_0$ and $V_2:H_2\to H_1$, the composite
compression is defined by $V_1V_2:H_2\to H_0$.  Along these maps entropy production
cannot increase; the zero-energy Schur complement is preserved before effective-space
compression; minimal positive-memory dilation rank cannot increase; and marked
monodromy descends exactly when collapsed cycles lie in its kernel.
\end{corollary}

\begin{proof}
The current and entropy statements are \Cref{thm:supp-current-chain}; the memory
statement and cutoff composition are \Cref{thm:supp-memory-RG}; and the monodromy
part is \Cref{thm:supp-monodromy-descent}.  For nested effective-space isometries,
\[
 V_2^*V_1^*(\cdot)V_1V_2=(V_1V_2)^*(\cdot)(V_1V_2),
\]
and compression of the operator-valued measure obeys the same identity.  Thus the
compressed memory coordinate also composes under the stated compatibility.  A
compression of a positive minimal memory-dilation space cannot increase its source
rank.
\end{proof}

% =============================================================================

\clearpage
\phantomsection
\addcontentsline{toc}{part}{Part II. Further mathematical extensions}
\section*{Part II. Further mathematical extensions}
\label{part:extensions}

The following sections extend, rather than replace, the preceding core.
The spin construction retains its local analytic estimates, the unconditional
flat periodic result, and the conditional global stable-atlas theorem.
An explicit positive dilation proof complements the shorter realization input
used in Part~I.  Rational indefinite memory and joint/completed realization
spaces are treated in their own declared categories.

\section{Analytic details for compact-spin spectralization}
\label{app:compact-spin}
% =============================================================================

The periodic $A_3$ construction in Section~\ref{sec:metric} supplies the explicit metric model.  This section records the analytic estimates used in the compact-spin convergence criterion.  The chartwise construction is verified below; the global theorem additionally assumes the atlas stability conditions of Definition~\ref{def:stable-spin-atlas}.

\subsection{A positive Renewal stencil for an arbitrary local inverse metric}

Let $d\ge2$ and put
\begin{equation}
\mathcal R_d^+
=\{e_i:1\le i\le d\}
\cup\{e_i+e_j,e_i-e_j:1\le i<j\le d\}.
\label{eq:supp-general-stencil}
\end{equation}
There are exactly $d^2$ unoriented directions.

\begin{lemma}[Positive $d^2$-direction decomposition]
\label{lem:supp-positive-decomposition}
Let $A=(A_{ij})\in\operatorname{Sym}_d(\mathbb R)$ and suppose that for some
$\eta>0$,
\begin{equation}
|A_{ij}|<\eta\quad(i\ne j),
\qquad
A_{ii}>(d-1)\eta.
\label{eq:supp-stencil-window}
\end{equation}
Define
\begin{equation}
w_i=A_{ii}-(d-1)\eta,
\qquad
w_{ij}^{\pm}=\frac12(\eta\pm A_{ij}).
\label{eq:supp-stencil-weights}
\end{equation}
Then every coefficient is strictly positive and
\begin{equation}
{
A=\sum_iw_ie_ie_i^{\mathsf T}
+\sum_{i<j}\left[
 w_{ij}^+(e_i+e_j)(e_i+e_j)^{\mathsf T}
+w_{ij}^-(e_i-e_j)(e_i-e_j)^{\mathsf T}
\right].}
\label{eq:supp-positive-decomposition}
\end{equation}
\end{lemma}

\begin{proof}
Positivity follows immediately from \eqref{eq:supp-stencil-window}.  In the
right-hand side of \eqref{eq:supp-positive-decomposition}, the $(i,j)$ off-diagonal
entry is $w_{ij}^+-w_{ij}^-=A_{ij}$.  Every pair containing $i$ contributes
$w_{ij}^++w_{ij}^-=\eta$ to the $i$th diagonal, so the diagonal coefficient is
$w_i+(d-1)\eta=A_{ii}$.
\end{proof}

A linear coordinate normalization makes $g^{-1}$ equal to $I_d$ at the centre of a
chart.  By continuity, after shrinking the chart, the inverse metric remains inside
one window \eqref{eq:supp-stencil-window}.  Compactness gives a finite good cover by
such charts.

Let $\rho=\sqrt{\det g}$ and write the decomposition coefficients of $g^{-1}$ as
$w_r(x)$, $r\in\mathcal R_d^+$.  Put $a_r(x)=\rho(x)w_r(x)$.  On
$Q_h=h\mathbb Z^d/\mathbb Z^d$ take
\begin{equation}
m_h(x)=h^d\rho(x),
\qquad
c_h(x,r)=\frac{h^{d-2}}4\bigl(a_r(x)+a_r(x+hr)\bigr).
\label{eq:supp-general-conductance}
\end{equation}
The two directed rates are the conductance divided by the corresponding endpoint
mass.

\begin{theorem}[Reversible coframe Renewal process]
\label{thm:supp-general-renewal-process}
The packet \eqref{eq:supp-general-conductance} is reversible with stationary mass
$m_h$.  Its predictable coordinate bracket is
\begin{equation}
{
\mathcal B_h(x)
=\frac1{4\rho(x)}
\sum_{r\in\mathcal R_d^+}
\bigl[a_r(x+hr)+2a_r(x)+a_r(x-hr)\bigr]rr^{\mathsf T},}
\label{eq:supp-general-bracket}
\end{equation}
and satisfies
\begin{equation}
\|\mathcal B_h-g^{-1}\|_{L^\infty}\le Ch^2.
\label{eq:supp-general-bracket-error}
\end{equation}
For every $f\in C^4$,
\begin{equation}
{\mathcal L_hf=\frac12\Delta_gf+O(h^2\|f\|_{C^4})}
\label{eq:supp-general-generator}
\end{equation}
uniformly on the chart, and the Dirichlet forms converge to
\begin{equation}
-\langle f_h,\mathcal L_hf_h\rangle_{m_h}
\longrightarrow
\frac12\int g^{ij}\partial_if\,\partial_j\overline f\,\rho\,\dd x.
\label{eq:supp-general-form-limit}
\end{equation}
\end{theorem}

\begin{proof}
Symmetry of $c_h$ gives detailed balance.  The bracket follows by summing the two
orientations of every direction through a vertex.  Taylor expansion of
$a_r(x\pm hr)$ cancels odd terms and, together with
\eqref{eq:supp-positive-decomposition}, gives
\eqref{eq:supp-general-bracket-error}.  Expanding
$f(x\pm hr)$ to fourth order, the first-order terms cancel and the second-order
terms contract with the same tensor decomposition.  The derivative of the physical
density and of the conductances produces the divergence-form first-order part of
$\Delta_g$, yielding \eqref{eq:supp-general-generator}.  Summation by parts gives
the finite Dirichlet form; the coefficient convergence and Riemann sums give
\eqref{eq:supp-general-form-limit}.
\end{proof}

\begin{proposition}[The metric bracket does not select the spin marking]
\label{cth:supp-metric-no-coframe}
The positive bracket determines $g^{-1}$ but not an oriented orthonormal coframe, a
spin lift, or a global spin structure.  A smooth rotation
$R(x)\in\mathrm{SO}(d)$ changes the coframe and spin connection without changing
$g^{-1}$ or the scalar Renewal generator, and a nontrivial $\mathbb Z_2$ spin
holonomy is invisible to every local metric bracket.
\end{proposition}

\subsection{Density-symmetric spin Dirac and finite Wilson packet}

Let $S_d$ be a complex spin module, let $\gamma_a$ be Hermitian Clifford matrices,
and let $E_a{}^j$ be the inverse oriented orthonormal coframe.  Set
\begin{equation}
c^j=E_a{}^j\gamma_a,
\qquad
c^jc^k+c^kc^j=2g^{jk}I,
\label{eq:supp-clifford-coeff}
\end{equation}
and let $\Omega_j$ be the Levi--Civita spin connection.  On
$L^2(Q,S_d;\rho\dd x)$,
\begin{equation}
D_g=-i\sum_jc^j(\partial_j+\Omega_j).
\label{eq:supp-geometric-dirac}
\end{equation}

\begin{lemma}[Exact density-symmetric Dirac identity]
\label{lem:supp-density-symmetric}
With $U_\rho\psi=\rho^{1/2}\psi$ and
$\nabla_j=\partial_j+\Omega_j$,
\begin{equation}
{
U_\rho D_gU_\rho^{-1}
=-\frac{i}{2}\sum_{j=1}^d
\bigl(M_{c^j}\nabla_j+\nabla_jM_{c^j}\bigr).}
\label{eq:supp-density-symmetric}
\end{equation}
\end{lemma}

\begin{proof}
Metric compatibility of Clifford multiplication gives
\[
\partial_jc^k+[\Omega_j,c^k]+\Gamma^k{}_{j\ell}c^\ell=0.
\]
Setting $k=j$ and using
$\Gamma^j{}_{j\ell}=\partial_\ell\log\rho$ gives
\[
\partial_jc^j+[\Omega_j,c^j]+c^j\partial_j\log\rho=0.
\]
Conjugating \eqref{eq:supp-geometric-dirac} by $\rho^{1/2}$ and applying this
contracted identity yields \eqref{eq:supp-density-symmetric}.
\end{proof}

Introduce one fixed two-dimensional normal factor and put
\begin{equation}
\widehat S_d=S_d\otimes\mathbb C^2,
\quad
\widehat c^j=c^j\otimes\sigma_1,
\quad
\Gamma_\perp=I\otimes\sigma_2,
\quad
\widehat\gamma=I\otimes\sigma_3.
\label{eq:supp-doubled-clifford}
\end{equation}
Let $T_{j,h}^{\Omega}$ be exact inverse spin parallel transport followed by a lattice
shift and define
\begin{equation}
P_{j,h}^{\Omega}
=\frac{T_{j,h}^{\Omega}-(T_{j,h}^{\Omega})^*}{2ih},
\qquad
W_h^{\Omega}
=\frac1{2h}\sum_{j=1}^d
\bigl(2I-T_{j,h}^{\Omega}-(T_{j,h}^{\Omega})^*\bigr).
\label{eq:supp-covariant-differences}
\end{equation}
The local doubled Wilson operator is the density-conjugated version of
\begin{equation}
\widetilde D_{g,h}^{\rm W}
=\frac12\sum_j
\bigl(M_{\widehat c_h^j}P_{j,h}^{\Omega}
      +P_{j,h}^{\Omega}M_{\widehat c_h^j}\bigr)
+\varpi\Gamma_\perp W_h^{\Omega}.
\label{eq:supp-general-Wilson}
\end{equation}

For the local estimates, extend the coefficients smoothly to a fixed periodic
coordinate box, preserving uniform ellipticity, and use the density-conjugated
$\ell^2$ norm with cell weight $h^d$.  Such an extension is an auxiliary local
calculation, not a boundary condition on a chart of $M$.  Let $S_{j,h}$ be the
ordinary coordinate shift.  In a smooth spin frame,
$T_{j,h}^{\Omega}-S_{j,h}=O(h)$ in operator norm.  Define
\begin{equation}
 \|u\|_{1,h}^2=\|u\|_h^2+
       \sum_j\left\|h^{-1}(S_{j,h}-I)u\right\|_h^2.
 \label{eq:discrete-H1-norm}
\end{equation}
The analogous covariant norm is uniformly equivalent to this one, since the
connection coefficients are bounded.  All constants below may depend on the
fixed atlas and finitely many coefficient bounds, but not on $h$.

\begin{theorem}[Frozen symbol, unique species, and uniform graph estimate]
\label{thm:supp-general-Wilson-ellipticity}
For fixed $\varpi>0$, the frozen principal symbol of
\eqref{eq:supp-general-Wilson} is
\begin{equation}
 q_{h,x_0}(\theta)=\frac1h\left[
 \sum_j\widehat c^j(x_0)\sin\theta_j
 +\varpi\Gamma_\perp\sum_j(1-\cos\theta_j)\right].
 \label{eq:supp-general-frozen-symbol}
\end{equation}
It satisfies
\begin{equation}
 q_{h,x_0}(\theta)^2=\frac1{h^2}\left[
 g^{jk}(x_0)\sin\theta_j\sin\theta_k
 +\varpi^2\Bigl(\sum_j(1-\cos\theta_j)\Bigr)^2\right]I.
 \label{eq:supp-general-frozen-square}
\end{equation}
Its only zero is $\theta=0$ on the Brillouin torus, within the fixed doubled spin module.  On the fixed smooth
periodic extension, and for sections localized in the interior of the original
chart, there is a cutoff-independent bound
\begin{equation}
 \|u\|_{1,h}\le C\bigl(\|\widetilde D_{g,h}^{\rm W}u\|_h+\|u\|_h\bigr).
 \label{eq:supp-general-Garding}
\end{equation}
The symmetric placement of coefficients makes the finite operator self-adjoint;
the fixed doubled Clifford convention supplies the grading.  A declared
real structure is retained when its compatibility is supplied as in
Definition~\ref{def:stable-spin-atlas}.
\end{theorem}

\begin{proof}
The Clifford relations eliminate the mixed terms in the square.  Put
$s(\theta)=\sum_j(1-\cos\theta_j)$ and
$\ell_h(\theta)=2h^{-2}s(\theta)$.  Uniform ellipticity and the elementary
identity
\[
 \sum_j\sin^2\theta_j+s(\theta)^2
 =2s(\theta)+2\sum_{i<j}(1-\cos\theta_i)(1-\cos\theta_j)
 \ge2s(\theta)
\]
give constants $0<c\le C<\infty$ such that
\begin{equation}
 c\ell_h(\theta)I\preceq q_{h,x_0}(\theta)^2
            \preceq C\ell_h(\theta)I.
 \label{eq:frozen-uniform-ellipticity}
\end{equation}
For the upper bound use $s\le2d$ and $\sum_j\sin^2\theta_j\le2s$.
The lower bound controls all forward differences, including modes near the
nonzero Brillouin corners.  Fourier transformation proves the frozen graph
estimate.

To pass to variable coefficients, choose a partition with supports of fixed
small diameter $r$, independent of $h$.  On one support the difference of the
variable principal part and its frozen value is bounded by
$C\omega_c(r)\|u\|_{1,h}+C_r\|u\|_h$, where $\omega_c$ is the modulus of
continuity of the coframe coefficients.  Derivatives of coefficients and the
replacement of ordinary shifts by smooth spin links contribute only the
second term.  Choose $r$ once so that the first term is absorbed by the frozen
estimate.  For every fixed smooth cutoff $\chi$,
\begin{equation}
 \|[\widetilde D_{g,h}^{\rm W},M_{\chi,h}]\|
          \le C\|\nabla\chi\|_\infty.
 \label{eq:local-cutoff-commutator}
\end{equation}
Indeed, a shift commutator contains $\chi(x+he_j)-\chi(x)=O(h)$,
which cancels the $h^{-1}$ factor in both the Dirac and Wilson differences.
Apply the local estimate to $\chi u$, sum the squared bounds over a fixed
finite partition with $\sum\chi^2=1$, and use the same discrete product rule
for the forward differences in \eqref{eq:discrete-H1-norm}.  This proves
\eqref{eq:supp-general-Garding}.  The fixed localization scale is important:
no cutoff derivative diverging as $h\to0$ is absorbed into a uniform constant.
\end{proof}

Let $\mathcal J_h^0$ be isometric piecewise-constant reconstruction and take
$S_h=(\mathcal J_h^0)^*$, the cell-average sampling map in the chosen density
and spin frame.  This avoids a dimension-dependent point-evaluation assumption
on Sobolev functions.  A piecewise-affine reconstruction $I_h$ on a fixed
shape-regular refinement of the coordinate cells satisfies
\begin{equation}
 \|I_hu\|_{H^1}\le C\|u\|_{1,h},\qquad
 \|I_hu-\mathcal J_h^0u\|_{L^2}\le Ch\|u\|_{1,h}.
 \label{eq:local-reconstruction-estimates}
\end{equation}
These inequalities follow cell by cell from the affine interpolation formula
and scaling to a unit cell.  Thus it is $I_hu$, not the discontinuous
piecewise-constant reconstruction itself, that has a uniform $H^1$ bound.

\begin{lemma}[Smooth-core consistency and the Wilson term]
\label{lem:supp-general-core}
For every smooth doubled spinor supported in the coordinate interior,
\begin{equation}
 \left\|\mathcal J_h^0\widetilde D_{g,h}^{\rm W}S_h\psi
       -U_\rho\widehat D_gU_\rho^{-1}\psi\right\|_{L^2}
       \le Ch\|\psi\|_{H^3}.
 \label{eq:supp-general-core}
\end{equation}
In particular, on this core,
\begin{equation}
 \|W_h^\Omega S_h\psi\|_h\le Ch\|\psi\|_{H^2}.
 \label{eq:Wilson-core-smallness}
\end{equation}
This is a core estimate, not an operator-norm smallness assertion on the full
lattice space.
\end{lemma}

\begin{proof}
Taylor's formula with integral remainder, applied to spin-parallel translated
cell averages, gives the covariant centered derivative $-i\nabla_j$ and the
second-difference expansion $W_h^\Omega=-(h/2)\sum_j\nabla_j^2$ on a smooth
core, with the corresponding Sobolev remainders.  Cell averaging and
reconstruction add $O(h)$ errors controlled by one further derivative.
Symmetric multiplication by the smooth coefficients and
Lemma~\ref{lem:supp-density-symmetric} give
\eqref{eq:supp-general-core}; applying the second-difference estimate alone
gives \eqref{eq:Wilson-core-smallness}.
\end{proof}

The two roles of the Wilson term must be distinguished.  On smooth sections it
vanishes as $h\to0$, whereas at nonzero lattice corners its size is of order
$h^{-1}$ and removes the spurious low-energy species.  It is therefore not a
bounded perturbation tending to zero in operator norm.  This distinction is
also central in rigorous continuum-limit studies of lattice Dirac operators
\cite{nielsenninomiya1981a,nielsenninomiya1981b,Nakamura2024}.  Moreover, for any bounded sequence $u_h$ and fixed smooth
$\psi$, self-adjointness gives
\[
 |\langle W_h^\Omega u_h,S_h\psi\rangle_h|
 =|\langle u_h,W_h^\Omega S_h\psi\rangle_h|
 \le Ch\|u_h\|_h\|\psi\|_{H^2}.
\]
Hence the regularizer leaves no additional distributional term in the limiting
Dirac equation.  On a closed manifold the intended self-adjoint realization is
the closure of $D_g$ on smooth spinors, with domain $H^1(M;S)$
\cite{lawson-michelsohn}; no chart-boundary domain is introduced.

\subsection{Atlas compatibility and the global resolvent argument}
\label{subsec:atlas-compatibility}

Local consistency and a positive overlap Gram do not, by themselves, imply
global norm-resolvent convergence.  The following formulation states the
required global approximation properties explicitly, so that the analytic
conclusion is separate from a choice of mesh and overlap realization.

\begin{definition}[Compatible stable spin discretization]
\label{def:stable-spin-atlas}
Let $\widehat D=D_{M,\mathfrak s}\otimes\sigma_1$ on
$\mathscr H=L^2(M;S\otimes\C^2)$, where $M$ is closed.  A compatible stable
spin discretization is a sequence of finite Hilbert spaces $\Hh_h$, self-adjoint
operators $D_h$, and isometries $W_h:\Hh_h\to\mathscr H$ satisfying:
\begin{enumerate}[label=\textup{(A\arabic*)},leftmargin=3em]
\item $W_hW_h^*\to I$ strongly.  There are discrete first-order norms
$\|\cdot\|_{1,h}$ and reconstructions $I_h:\Hh_h\to H^1(M;S\otimes\C^2)$
with
\begin{equation}
 \|I_hu\|_{H^1}\le C\|u\|_{1,h},\qquad
 \|I_hu-W_hu\|\le\varepsilon_h\|u\|_{1,h},\quad
 \varepsilon_h\longrightarrow0.
 \label{eq:atlas-reconstruction}
\end{equation}
\item The assembled, rather than merely chartwise, operators satisfy
\begin{equation}
 \|u\|_{1,h}\le C(\|D_hu\|_h+\|u\|_h).
 \label{eq:atlas-global-stability}
\end{equation}
\item For every smooth doubled spinor $\psi$ there are samples $S_h\psi$ with
\begin{equation}
 W_hS_h\psi\longrightarrow\psi,\qquad
 W_hD_hS_h\psi\longrightarrow\widehat D\psi.
 \label{eq:atlas-global-core}
\end{equation}
The global samples use the transition functions of the fixed spin bundle on
shared degrees of freedom, including coordinate and density changes.  When
real and grading structures are asserted, their finite actions intertwine
these reconstructions and converge to the declared continuum actions.
\end{enumerate}
\end{definition}

Here is the cancellation that a proposed atlas assembly must respect.  Choose
fixed smooth real cutoffs $\chi_\alpha$ supported in chart interiors with
$\sum_\alpha\chi_\alpha^2=1$.  After transporting local spinors and their
density factors to the same bundle, the continuum first-order identity is
\begin{equation}
 \sum_\alpha\chi_\alpha\widehat D\chi_\alpha=\widehat D,
 \qquad
 \sum_\alpha\chi_\alpha[\widehat D,\chi_\alpha]
       =\tfrac12[\widehat D,\sum_\alpha\chi_\alpha^2]=0.
 \label{eq:atlas-cutoff-cancellation}
\end{equation}
Thus the apparent order-one derivatives of the cutoffs cancel; one must not
assert that they individually tend to zero.  For a finite-range Wilson
operator on a common coordinate mesh the analogous localization identity is
\begin{equation}
 \sum_\alpha M_{\chi_\alpha,h}W_h^\Omega M_{\chi_\alpha,h}-W_h^\Omega
 =\tfrac12\sum_\alpha
 [M_{\chi_\alpha,h},[W_h^\Omega,M_{\chi_\alpha,h}]],
 \label{eq:Wilson-double-commutator}
\end{equation}
and its norm is at most $Ch\sum_\alpha\|\nabla\chi_\alpha\|_\infty^2$.
Each matrix entry contains two cutoff differences and only one factor $h^{-1}$.
Fixed collars keep all such finite-range stencils away from auxiliary chart
boundaries for sufficiently small $h$.

On a shared geometric link $x\to y$, spin transition matrices satisfy the
exact covariance identity
\[
 U_{\alpha\beta}(x)\,T^\beta_{xy}\,U_{\alpha\beta}(y)^{-1}
       =T^\alpha_{xy}.
\]
It removes frame mismatch on that link.  It does not identify different
coordinate-grid links or supply a stable interpolation between unrelated grids.
For independently meshed charts, such interpolation and its first-order error
bounds must be verified in \eqref{eq:atlas-reconstruction}--
\eqref{eq:atlas-global-core}.  In particular, positivity of a Gram matrix is
not a substitute for the assembled graph estimate
\eqref{eq:atlas-global-stability}.  The local Wilson estimates above establish
the local ingredients; the existence of a global assembly satisfying all these
properties is an additional construction requirement.

\begin{theorem}[Norm-resolvent convergence for compatible stable atlases]
\label{thm:supp-compact-spin-resolvent}
For a compatible stable spin discretization in the sense of
Definition~\ref{def:stable-spin-atlas}, and every $z\notin\R$,
\begin{equation}
 \left\|W_h(D_h-z)^{-1}W_h^*-(\widehat D-z)^{-1}\right\|
       \longrightarrow0.
 \label{eq:supp-global-norm-resolvent}
\end{equation}
The same argument on a fixed smooth periodic coordinate extension gives
\begin{equation}
 \left\|\mathcal J_h^0(D_{g,h}^{\rm W}-z)^{-1}(\mathcal J_h^0)^*
       -(\widehat D_g-z)^{-1}\right\|\longrightarrow0.
 \label{eq:supp-local-norm-resolvent}
\end{equation}
This local statement is about that self-adjoint periodic extension, not an
unspecified self-adjoint operator on an open chart.  Isolated bounded-energy
spectral projections converge, and compatible real and grading structures
are retained as specified in Definition~\ref{def:stable-spin-atlas}.
\end{theorem}

\begin{proof}
Write $R_h(z)=W_h(D_h-z)^{-1}W_h^*$ and
$R(z)=(\widehat D-z)^{-1}$.  Self-adjointness gives
\begin{equation}
 \|(D_h-z)^{-1}\|\le|\operatorname{Im}z|^{-1},\qquad
 \|D_h(D_h-z)^{-1}\|\le1+|z|\,|\operatorname{Im}z|^{-1}.
 \label{eq:uniform-resolvent-bounds}
\end{equation}
For $\|f\|\le1$ put $u_h=(D_h-z)^{-1}W_h^*f$.  The global graph estimate
bounds $\|u_h\|_{1,h}$ uniformly.  Therefore $I_hu_h$ lies in a bounded
$H^1$ set, while $\|W_hu_h-I_hu_h\|\to0$ uniformly on this resolvent unit
ball.  Rellich compactness proves collective compactness of the embedded
resolvents along the refining sequence.  This step uses the reconstruction
estimate and does not assign an $H^1$ norm to a piecewise-constant function.

For fixed $f$, let $W_hu_h\to u$ along a subsequence.  Test the discrete
resolvent equation against $S_h\psi$, with $\psi$ any smooth doubled spinor.
Self-adjointness and \eqref{eq:atlas-global-core} pass the equation to
$(\widehat D-z)u=f$ in the distributional sense.  Since smooth spinors are a
core for the self-adjoint operator $\widehat D$, this places $u$ in its
adjoint domain, which is $H^1$, and gives $u=R(z)f$.  Every convergent
subsequence has this limit, so $R_h(z)\to R(z)$ strongly.  Applying the same
argument at $\bar z$ also gives strong convergence of the adjoints.

For clarity, this last fact is needed for the upgrade to norm convergence.
If that convergence failed, choose unit vectors $f_h$ with
$\|(R_h(z)-R(z))f_h\|\ge\epsilon$.  Passing to subsequences, let
$f_h\rightharpoonup f$ and use collective compactness to make $R_h(z)f_h$
converge in norm.  For each fixed $v$, strong adjoint convergence implies
\[
 \langle v,R_h(z)f_h\rangle
 =\langle R_h(\bar z)v,f_h\rangle
 \longrightarrow\langle R(\bar z)v,f\rangle.
\]
Hence its norm limit is $R(z)f$.  Compactness of $R(z)$ also gives
$R(z)f_h\to R(z)f$, a contradiction.  This proves
\eqref{eq:supp-global-norm-resolvent}.  On the periodic extension,
\eqref{eq:local-reconstruction-estimates},
\eqref{eq:supp-general-Garding}, and \eqref{eq:supp-general-core} provide the
same hypotheses.  The spectral-projection conclusion is the usual consequence
of self-adjoint norm-resolvent convergence for isolated spectral clusters.
\end{proof}

\begin{corollary}[Unconditional periodic flat-spin norm-resolvent convergence]
\label{cor:supp-flat-torus-spin}
Let $M=\R^d/\Z^d$ carry a constant flat metric and the periodic spin structure.
Use the uniform periodic lattice $Q_h=h\Z^d/\Z^d$, constant Clifford
coefficients, trivial spin parallel transport, and the Wilson operator
\eqref{eq:supp-general-Wilson} with fixed $\varpi>0$.  With the isometric
piecewise-constant embedding $\mathcal J_h^0$,
\begin{equation}
 \left\|\mathcal J_h^0(D_{g,h}^{\rm W}-z)^{-1}(\mathcal J_h^0)^*
 -(\widehat D_M-z)^{-1}\right\|\longrightarrow0,
 \qquad z\in\C\setminus\R.
 \label{eq:supp-flat-torus-resolvent}
\end{equation}
Thus the compatible-stability hypotheses are automatically satisfied on this
periodic flat class.
\end{corollary}

\begin{proof}
There is no atlas-gluing problem.  The constant-coefficient Fourier estimate
\eqref{eq:frozen-uniform-ellipticity} is global and gives
\eqref{eq:atlas-global-stability}; the periodic affine and piecewise-constant
reconstructions satisfy \eqref{eq:atlas-reconstruction}; and
Lemma~\ref{lem:supp-general-core} gives \eqref{eq:atlas-global-core} globally on
smooth periodic spinors.  Moreover $\mathcal J_h^0(\mathcal J_h^0)^*$ is the
cellwise conditional expectation and converges strongly to the identity on
$L^2(M;S\otimes\C^2)$.  Hence all conditions of
Definition~\ref{def:stable-spin-atlas} hold and
Theorem~\ref{thm:supp-compact-spin-resolvent} gives
\eqref{eq:supp-flat-torus-resolvent}.
\end{proof}

\subsection{Independent metric approximation and the spin conclusion}

The endpoint-block construction and its complete uniform geodesic-convergence
proof are given in Section~\ref{app:endpoint-metric}.  They are independent of
the spin-assembly hypotheses and are not repeated here.

\begin{corollary}[Conditional stable compact-spin approximation]
\label{cor:supp-compact-spin-image}
A closed marked Riemannian spin manifold admitting a compatible stable spin
discretization has a finite-dimensional approximation of its doubled spin
Dirac operator in generalized norm resolvent, together with convergent
independent metric heads.  Consequently, essential surjectivity of this
particular atlas-based spin construction follows if such compatible stable
discretizations are supplied for all target manifolds.
\end{corollary}

\begin{proof}
Apply Theorem~\ref{thm:supp-compact-spin-resolvent} to the spin operators and
Theorem~\ref{thm:supp-global-metric} to the independent metric heads.
The last assertion records explicitly the additional existence requirement.
\end{proof}

\section{Explicit minimal realization of positive Stieltjes memory}
\label{app:positive-dilation}

Section~\ref{app:memory-details} uses the classical source-minimal positive
realization theorem.  For completeness, this extension retains the explicit
construction and uniqueness proof.  It starts from the same bounded positive
operator-valued measure on $(0,\infty)$ as Definition~\ref{def:supp-POVM};
no additional static or moment determinacy hypothesis is imposed here.

\begin{construction}[Canonical minimal Stieltjes dilation]
\label{con:supp-Naimark-tail}
Let $\mathscr S_\mu$ be the finite $H$-valued Borel simple functions and set
\[
 \left\langle
 \sum_j\one_{E_j}\xi_j,
 \sum_k\one_{F_k}\eta_k
 \right\rangle_\mu
 =
 \sum_{j,k}
 \langle\xi_j,\mu(E_j\cap F_k)\eta_k\rangle.
\]
Quotient by the null space and complete to $K_\mu$.  Multiplication by
$\one_E$ defines a projection-valued measure $P_\mu(E)$, and
\[
 C_\mu=\int\lambda\,\dd P_\mu(\lambda),
 \qquad
 V_\mu\xi=[\one_{(0,\infty)}\xi],
 \qquad
 B_\mu=V_\mu^*.
\]
\end{construction}

\begin{theorem}[Operator-measure realization and minimality]
\label{thm:supp-Naimark-realization}
Every bounded positive operator-valued measure satisfies
\begin{equation}
 {
 \Sigma_\mu(z)
 =B_\mu(C_\mu-z)^{-1}B_\mu^*.}
 \label{eq:extension-Naimark-realization}
\end{equation}
Moreover
\begin{equation}
 {
 K_\mu
 =\overline{\Span}\{
 P_\mu(E)B_\mu^*\xi:
 E\in\mathfrak B((0,\infty)),\ \xi\in H\}.}
 \label{eq:supp-Naimark-minimal}
\end{equation}
Conversely, if $C\ge0$ is self-adjoint on $K$ with $E_C(\{0\})=0$,
$B:K\to H$ is bounded, and
\[
 K=\overline{\Span}\{E_C(E)B^*H\},
\]
then $\mu(E)=BE_C(E)B^*$ is a bounded positive operator-valued measure with
self-energy $B(C-z)^{-1}B^*$.  Two minimal realizations of the same
$\mu$ are related by a unique unitary $U$ satisfying
\begin{equation}
 {UC_1=C_2U,\qquad UB_1^*=B_2^*.}
 \label{eq:supp-Naimark-unitary}
\end{equation}
\end{theorem}

\begin{proof}
The spectral calculus gives
\[
 B_\mu(C_\mu-z)^{-1}B_\mu^*
 =V_\mu^*
 \left(\int\frac{1}{\lambda-z}\,\dd P_\mu(\lambda)\right)V_\mu
 =\int\frac{1}{\lambda-z}\,\dd\mu(\lambda).
\]
Every simple function is a finite sum of vectors
$P_\mu(E)V_\mu\xi$, proving \eqref{eq:supp-Naimark-minimal}.

For the converse, the spectral theorem gives positivity and countable
additivity of $BE_C(E)B^*$.  On the algebraic span of spectral-source vectors
define
\[
 U_0\left(\sum_jE_C(E_j)B^*\xi_j\right)
 =
 \sum_jP_\mu(E_j)V_\mu\xi_j.
\]
The squared norms on the two sides are both
\[
 \sum_{j,k}
 \langle\xi_j,BE_C(E_j\cap E_k)B^*\xi_k\rangle,
\]
so $U_0$ is well defined and isometric.  Source minimality makes its extension
unitary.  It intertwines every spectral projection, hence $C$, and sends
$B^*$ to $V_\mu$.  Applying the same construction to two realizations proves
existence and uniqueness of \eqref{eq:supp-Naimark-unitary}.
\end{proof}

\section{Finite rational indefinite memory}
\label{app:rational-krein}
% =============================================================================

The positive Stieltjes theory in Section~\ref{app:memory-details} is not
the only finite realization class considered here.  A finite Krein/Pontryagin tail is instead controlled by the number of
negative squares of its generalized Nevanlinna kernel.

\begin{definition}[Normalized rational Hermitian dynamic function]
\label{def:supp-rational-Hermitian}
Let $H$ be finite dimensional and let $Q$ be an $\End(H)$-valued rational function
satisfying
\begin{equation}
Q(\bar z)^*=Q(z),
\qquad
\lim_{z\to\infty}Q(z)=0.
\label{eq:supp-rational-symmetry}
\end{equation}
Write its Laurent expansion at infinity as
\begin{equation}
Q(z)=-\sum_{n\ge0}M_nz^{-n-1},
\qquad M_n=M_n^*.
\label{eq:supp-rational-Laurent}
\end{equation}
Let $\mathcal K_Q$ be the span of Hankel columns
$c_j(h)=(M_{j+r}h)_{r\ge0}$ and define
\begin{equation}
[c_i(u),c_j(v)]_Q=\langle u,M_{i+j}v\rangle.
\label{eq:supp-Hankel-form}
\end{equation}
After quotienting the radical of the Hermitian form, let $A_Qc_j(h)=c_{j+1}(h)$ and
$\Gamma_Qh=c_0(h)$.
\end{definition}

\begin{theorem}[Complete normalized rational generalized-Nevanlinna realization]
\label{thm:supp-complete-rational-Krein}
The construction above is a minimal finite Pontryagin realization and
\begin{equation}
{Q(z)=\Gamma_Q^+(A_Q-z)^{-1}\Gamma_Q.}
\label{eq:supp-rational-realization}
\end{equation}
Its ordinary dimension is the stabilized coefficient-Hankel rank
\begin{equation}
\dim\mathcal K_Q=\McM(Q),
\label{eq:supp-rational-degree}
\end{equation}
and its negative index is the inertia of
\eqref{eq:supp-Hankel-form}.  The Nevanlinna kernel
\begin{equation}
N_Q(z,w)=\frac{Q(z)-Q(w)^*}{z-\bar w}
\label{eq:supp-Nevanlinna-kernel}
\end{equation}
has exactly that many negative squares.  Every other minimal finite
Pontryagin realization of $Q$ is related to this one by a unique source-fixing
Pontryagin unitary.
\end{theorem}

\begin{proof}
By construction, the output functional $L_Q$ on the Hankel column model obeys
$L_QA_Q^n\Gamma_Q=M_n$.  For large $z$,
\[
L_Q(A_Q-z)^{-1}\Gamma_Q
=-\sum_{n\ge0}\frac{M_n}{z^{n+1}}=Q(z),
\]
and rational continuation proves \eqref{eq:supp-rational-realization}.  The
power-cyclic span of $\Gamma_QH$ is all of $\mathcal K_Q$ by definition.  In finite
dimension, the resolvent-source span equals the power-cyclic span: expansion at
infinity gives one inclusion and a finite Vandermonde inversion at sufficiently many
regular spectral parameters gives the other.  Thus the realization is source
minimal.

The dimension is the stabilized Hankel rank by construction.  Choosing a Hilbert
majorant and a congruence which writes the nondegenerate Hankel form as a fundamental
symmetry $J_Q$ gives $A_Q^*J_Q=J_QA_Q$ and
$Q(z)=\Gamma_Q^*J_Q(A_Q-z)^{-1}\Gamma_Q$.  The standard resolvent identity gives
\[
N_Q(z,w)
=\Gamma_Q^*J_Q(A_Q-z)^{-1}(A_Q-\bar w)^{-1}\Gamma_Q.
\]
Every finite kernel Gram is therefore a compression of the Pontryagin form, so it
has at most its negative index.  Source minimality makes the resolvent-source vectors
span the whole state space and attains the full index.  Finally, if two
minimal realizations have the same $Q$, their complete Hankel moments agree.
Mapping $A_Q^j\Gamma_Qh$ to the corresponding cyclic vector in the second
realization preserves the indefinite form, is onto, fixes the source, and
intertwines the state operators.  Uniqueness follows from cyclicity.
\end{proof}

\begin{theorem}[Pole--Hankel invariants]
\label{thm:supp-pole-Hankel}
For a normalized rational Hermitian $Q$:
\begin{enumerate}[label=\textup{(P\arabic*)},leftmargin=3em]
\item the principal part at each pole determines the local McMillan degree and all
      Jordan block sizes through shifted principal-part Hankel ranks;
\item at each real pole, the principal-part Hankel Gram determines the complete
      local Pontryagin inertia;
\item every nonreal conjugate pole pair forms a neutral paired root space, and its
      upper-half-plane local McMillan degree contributes equally to positive and
      negative index;
\item the sum of the local contributions equals the minimum global negative index.
\end{enumerate}
\end{theorem}

\begin{proof}
At a pole $\lambda$, multiply by a sufficiently high power of $(z-\lambda)$ and
read the finite principal-part coefficient sequence.  The ranks of its shifted
Hankel matrices are the controllability/observability ranks of the nilpotent local
state operator and their successive differences recover the Jordan chain lengths.
For a real pole the local Hankel form is Hermitian and Sylvester inertia is invariant
under every minimal state-coordinate change, hence gives the local Pontryagin
signature.  Hermitian symmetry pairs a nonreal root space with its conjugate; their
indefinite form is neutral off-diagonal, so the two signs occur with equal
multiplicity.  The primary decomposition is Pontryagin orthogonal, and the global
index is the sum of local indices.
\end{proof}

A general finite rational Hermitian function has a unique split
\begin{equation}
Q(z)=P(z)+Q_0(z),
\qquad
\lim_{z\to\infty}Q_0(z)=0,
\label{eq:supp-rational-split}
\end{equation}
with a Hermitian polynomial $P$.  The following descriptor statement uses
the minimal regular, source-fixing Pontryagin descriptor convention and its
finite/infinity primary decomposition; it is separate from the explicit
strictly proper Hankel construction above.

\begin{corollary}[All finite rational Hermitian dynamics]
\label{cor:supp-all-rational}
The strictly proper part $Q_0$ is realized by
\Cref{thm:supp-complete-rational-Krein}.  The nonconstant polynomial jet $P$ has a
canonical finite infinity-chain/regular descriptor realization.  Their primary
direct sum gives an equivalence, up to source-fixing descriptor Pontryagin unitary,
\begin{equation}
{
\begin{gathered}
\text{finite rational Hermitian dynamic functions}\\
\longleftrightarrow\\
\text{minimal regular Pontryagin descriptor realizations}.
\end{gathered}}
\label{eq:supp-all-rational-equivalence}
\end{equation}
This includes finite real and nonreal poles, arbitrary finite Jordan chains, and all
polynomial jets at infinity.
\end{corollary}

\begin{proof}
The pencil of any regular finite descriptor realization has a primary decomposition
at infinity and at finite spectral points.  On the finite primary block the transfer
function is strictly proper and is minimized by
\Cref{thm:supp-complete-rational-Krein}.  On the generalized zero eigenspace of an
anchored descriptor pencil, the resolvent expansion is polynomial and is represented
by the canonical nilpotent infinity chain.  The direct sum attains both minimum
state dimensions and indices.  Any second minimal realization splits into the
same invariant primary blocks, and the unique source-fixing unitaries on the finite
and infinity blocks combine to the unique descriptor equivalence.
\end{proof}

\begin{remark}[Positive Stieltjes memory is the index-zero branch]
\label{scope:supp-positive-index-zero}
A positive Stieltjes memory packet has zero negative-square index.  A nonzero
physical Pontryagin index requires an independently occurring fundamental-symmetry
or signed packet.  Neither OS reflection, temporal orientation, charge conjugation,
nor a determinant/Pfaffian line determines that sign by itself.  The unrestricted
nonrational infinite-dimensional indefinite essential image is not claimed here.
\end{remark}

\subsection{Effective-space compression of indefinite memory}

\begin{proposition}[Negative-square index under effective-space compression]
\label{prop:extension-indefinite-compression}
Let $Q$ be a normalized rational Hermitian dynamic function and
$V:H'\to H$ an isometry.  Put $Q_V(z)=V^*Q(z)V$.
Then the negative-square index of $N_{Q_V}$ does not exceed that of $N_Q$,
and the minimal finite realization dimension of $Q_V$ does not exceed that
of $Q$.  Nested isometries give the same compressed function as their composite.
\end{proposition}

\begin{proof}
For any finite set of regular spectral parameters, the kernel Gram of
$N_{Q_V}(z,w)=V^*N_Q(z,w)V$ is a congruence of the corresponding kernel Gram
of $N_Q$.  Its number of negative eigenvalues therefore cannot increase.
A source-minimal realization of $Q$ gives a realization of $Q_V$ on the same
state space by replacing its source map $\Gamma$ by $\Gamma V$; minimizing
cannot increase the state dimension.  Finally,
$V_2^*V_1^*Q(z)V_1V_2=(V_1V_2)^*Q(z)(V_1V_2)$ gives composition.
\end{proof}

\section{Joint realizations and completed inverse fibres}
\label{app:completed-fibres}

This extension concerns the declared retained-coordinate realization model.
The elementary joint Gram classification and compact inverse-system pruning
are distinguished from the broader fibre descriptions, which use the
full-reduction and compatibility hypotheses stated below.  In particular,
these hypotheses are not derived from the stationary-current theorem, and the
section does not assert that current, memory, and monodromy exhaust every
possible predictive realization.

\subsection{Joint Gram realizations and the finite stratified fibre}

Let $E_1,\ldots,E_m$ be finite Hilbert spaces representing minimal marginal
realization spaces, normalized so that their marginal Gram operators are the
identity.  A joint realization is encoded by a block correlation operator
\begin{equation}
\mathcal K=[K_{ab}]_{a,b=1}^m,
\qquad
K_{aa}=I_{E_a},
\qquad
\mathcal K\succeq0.
\label{eq:supp-correlation-spectrahedron}
\end{equation}
Let $\Corr_X$ denote this correlation spectrahedron and let $\mathcal G_X$ be
the product of the unitary gauge groups of the marginal realization spaces.

\begin{theorem}[Joint Gram classification]
\label{thm:supp-common-source}
Every minimal joint realization of the retained marginal Gram blocks is
classified, up to a unitary on the common realization space, by a point of
$\Corr_X/\mathcal G_X$.  The support rank of $\mathcal K$ is the minimum
possible joint realization dimension.  A representative $\mathcal K$ has
residual gauge group exactly $\Stab_{\mathcal G_X}(\mathcal K)$.
\end{theorem}

\begin{proof}
A joint realization has a Gram factorization $\mathcal K=S^*S$, with the
identity diagonal expressing the chosen normalization of the marginal blocks.
Conversely every positive block operator with identity diagonal admits such a
factorization.  Restricting the realization space to $\overline{\Ran S}$ gives
the minimal realization, whose dimension is $\rank\mathcal K$ in the finite
case.  Two minimal Gram factorizations of the same positive operator are
related by a unique unitary on the common range.  Independent changes of
marginal frames act by block congruence, and the residual action is precisely
the stated stabilizer.
\end{proof}

\begin{proposition}[Pairwise positivity need not give a joint realization]
\label{cth:supp-pairwise-not-global}
Pairwise admissible Gram blocks need not assemble into one positive joint Gram
operator.  For example,
\begin{equation}
\begin{pmatrix}
1&-3/4&-3/4\\
-3/4&1&-3/4\\
-3/4&-3/4&1
\end{pmatrix}
\label{eq:supp-three-block-counterexample}
\end{equation}
has every $2\times2$ principal block positive but has eigenvalue $-1/2$.
Thus pairwise compatibility does not imply global Gram compatibility.
\end{proposition}

Let $\mathcal J_X$ be the current fibre, $\mathcal M_X$ the complete
dynamic-memory fibre over the fixed spectral data, $\mathcal P_X$ the retained
auxiliary-marking fibre, $\mathcal C_X^{\rm cal}$ the residual calibration
class, and $\mathcal R_X$ the marked monodromy representation space.  We call a
finite target \emph{fully reduced} when all null and spectator directions have
been quotiented out, the domain, weight, orientation, determinant-line, normalization, and auxiliary-marking
conventions are fixed, every retained Gram block is minimal, and a specified
finite family of bounded separating functionals (equivalently, finite-rank
defect operators) has trivial common kernel on any additional retained
direction.  This last condition is the standard separation hypothesis that
replaces the programme-specific ``zero-exhaustion'' terminology.

\begin{theorem}[Finite stratified fibre in the fully reduced coordinate model]
\label{thm:supp-finite-exhaustion}
For a completed fully reduced finite target $\mathsf S^\sharp$,
\begin{equation}
{
\RReal_X^{\sharp,\min}(\mathsf S)
\cong
\bigsqcup_{[\mathcal K]\in\Corr_X/\mathcal G_X}
\frac{
\mathcal J_X\times\mathcal M_X\times\mathcal P_X
\times\mathcal C_X^{\rm cal}\times\mathcal R_X
}{\Stab_{\mathcal G_X}(\mathcal K)}.}
\label{eq:supp-finite-fibre}
\end{equation}
The factors vary, respectively, directed dynamics, dynamical memory,
auxiliary markings, calibration, and global bundle data while preserving
the fixed spectral packet.  In general the quotient is stratified and not a
Cartesian product.
\end{theorem}

\begin{proof}
The declared coordinate-completeness part of full reduction restricts the
realization class to the listed retained directions after the null quotients.
The separating family distinguishes those directions; it does not establish
that every predictive model belongs to this class.  The graph inverse theorem classifies the directed-current freedom; the
positive and Pontryagin/descriptor theorems classify returned dynamic memory;
the retained auxiliary markings give $\mathcal P_X$; the calibration data give
$\mathcal C_X^{\rm cal}$; and marked bundle descent gives $\mathcal R_X$.
Theorem~\ref{thm:supp-common-source} classifies all compatible joint Gram
realizations and identifies their residual stabilizers.  Equivalence of two
fully reduced realizations is therefore simultaneous equivalence in these
factors under the common stabilizer, yielding \eqref{eq:supp-finite-fibre}.
Proposition~\ref{cth:supp-pairwise-not-global} explains why the general fibre
cannot be replaced by a product of independently quotiented factors.
\end{proof}

\subsection{Quasilocal and semifinite completion}

For completion one must distinguish the finite inverse fibre relevant to strictly
positive realizations from a compact finite-stage carrier used in the projective
limit.  For the current factor define the closed polytope
\begin{equation}
 \overline{\mathcal J}(G,c)
 =\{j\in\Ker B:\ |j_e|\le c_e\},
 \label{eq:supp-closed-current-polytope}
\end{equation}
whose relative interior is the strict-rate fibre $\mathcal J(G,c)$.  It is compact.
For positive memory, a finite-stage compactification is obtained by placing the
operator-valued tail measure on the one-point compactification of the retained
energy range and imposing a uniform bound on its total trace mass; in finite source
dimension this set is weak-$*$ compact.  Correlation spectrahedra are compact,
unitary gauge groups are compact, and for a finite incidence graph with compact
packet automorphism group the marked-monodromy space is a compact quotient of a
finite product of that group.  We therefore formulate the completion theorem on
compact finite-stage strata, rather than silently assuming compactness of the open
physical fibres.

\begin{definition}[Compact finite-stage inverse system]
\label{def:supp-compact-fibre-system}
A compact finite-stage inverse system consists of increasing finite-trace screens
$P_n\uparrow I$ and nonempty compact Hausdorff fully reduced strata obtained as
the diagonal compact-group quotient
\begin{equation}
 \overline{\mathfrak F}_n
 :=\frac{
 \Corr_n\times\overline{\mathcal J}_n\times\overline{\mathcal M}_n
 \times\overline{\mathcal P}_n\times\overline{\mathcal C}_n^{\rm cal}
 \times\overline{\mathcal R}_n
 }{\mathcal G_n}.
 \label{eq:supp-compact-finite-stratum}
\end{equation}
Equivalently, its orbit-type decomposition is
\[
 \overline{\mathfrak F}_n
 \cong
 \bigsqcup_{[\mathcal K_n]\in\Corr_n/\mathcal G_n}
 \frac{
 \overline{\mathcal J}_n\times\overline{\mathcal M}_n
 \times\overline{\mathcal P}_n\times\overline{\mathcal C}_n^{\rm cal}
 \times\overline{\mathcal R}_n
 }{\Stab(\mathcal K_n)}.
\]
The displayed disjoint union is an orbit-type stratification, not a disjoint-union
topology.  The factors use the closed current polytope
\eqref{eq:supp-closed-current-polytope}, weak-$*$ compact tight memory sets with
uniform trace-mass bounds, compact marking/calibration sets, and compact monodromy
carriers, with the natural continuous $\mathcal G_n$ actions.  Since
$\mathcal G_n$ is compact, \eqref{eq:supp-compact-finite-stratum} is compact
Hausdorff.  There are continuous restriction maps
$r_{m,n}:\overline{\mathfrak F}_m\to\overline{\mathfrak F}_n$ for $m\ge n$.
Compatible self-adjoint resolvent data and a compatible normal semifinite weight
are retained at each stage.  The restriction maps satisfy
$r_{n,n}=\id$ and $r_{\ell,n}=r_{m,n}r_{\ell,m}$ for $\ell\ge m\ge n$.
The operator-level completion statement additionally assumes that the
transported limiting resolvent family satisfies the resolvent identity,
adjoint symmetry, injectivity, and dense range, and that the retained
finite-trace weights admit the compatible normal semifinite extension used
below.  These requirements are distinct from compactness of the coordinate
inverse system.

For $m\ge n$ put
\begin{equation}
 \overline{\mathfrak F}_n^{[m]}
 =r_{m,n}(\overline{\mathfrak F}_m),
 \qquad
 \overline{\mathfrak F}_n^{\infty}
 =\bigcap_{m\ge n}\overline{\mathfrak F}_n^{[m]}.
 \label{eq:supp-stable-fibre-images}
\end{equation}
Finally, choose compatible finite families $\mathscr F_n$ of bounded separating
functionals or finite-rank defect operators such that the cofinal family
$\bigcup_n\mathscr F_n$ has trivial common kernel on every additional retained
quasilocal direction after quotienting by the displayed coordinates.
\end{definition}

\begin{lemma}[Canonical pruning of compact finite-stage fibres]
\label{lem:supp-fibre-pruning}
For every system of Definition~\ref{def:supp-compact-fibre-system}, each stable
image $\overline{\mathfrak F}_n^\infty$ is nonempty and compact, and
\begin{equation}
 r_{n+1,n}(\overline{\mathfrak F}_{n+1}^\infty)
 =\overline{\mathfrak F}_n^\infty.
 \label{eq:supp-stable-surjectivity}
\end{equation}
Moreover $\overline{\mathfrak F}_n^\infty$ is exactly the level-$n$ projection of
$\varprojlim_m\overline{\mathfrak F}_m$ and is the largest subsystem in which every
retained finite-stage point extends indefinitely.
\end{lemma}

\begin{proof}
For fixed $n$, the sets
$\overline{\mathfrak F}_n^{[m]}$ are nonempty compact and decrease with $m$, so
their intersection is nonempty compact.  The forward inclusion in
\eqref{eq:supp-stable-surjectivity} is immediate.  For the reverse inclusion, fix
$a\in\overline{\mathfrak F}_n^\infty$ and consider
\[
 B_m=\{b\in\overline{\mathfrak F}_{n+1}^{[m]}:
       r_{n+1,n}(b)=a\},\qquad m\ge n+1.
\]
Each $B_m$ is nonempty compact and the family is decreasing, so its intersection
contains a point of $\overline{\mathfrak F}_{n+1}^\infty$ mapping to $a$.
The projection and maximality statements follow directly from the definition of
the stable images.
\end{proof}

\begin{theorem}[Completed retained-coordinate inverse-fibre theorem]
\label{thm:supp-semifinite-exhaustion}
For a compact finite-stage inverse system satisfying the compatible resolvent,
weight, transport, tightness, and cofinal-separation hypotheses of
Definition~\ref{def:supp-compact-fibre-system}, the retained completed inverse
fibre is
\begin{equation}
 \boxed{
 \RReal_\infty^{\sharp,\min}(\mathsf S)
 \cong
 \varprojlim_n\overline{\mathfrak F}_n^\infty.}
 \label{eq:supp-semifinite-fibre}
\end{equation}
The inverse limit is nonempty and compact.  It contains no additional retained
quasilocal coordinate beyond the displayed finite-stage factors.  Its strict-rate
physical part is obtained by intersecting the compatible threads with the relative
interiors $\mathcal J_n\subset\overline{\mathcal J}_n$ and with the corresponding
nondegenerate interiors of the other compactified factors.  On a factor-preserving
stratum the limit may be written componentwise before the compatible compact gauge
quotient; without factor-preserving transport the stratified form is essential.
\end{theorem}

\begin{proof}
Lemma~\ref{lem:supp-fibre-pruning} replaces any possibly nonsurjective finite-stage
system by its canonical stable subsystem, whose bonding maps are surjective.  The
product $\prod_n\overline{\mathfrak F}_n^\infty$ is compact, and the compatibility
equations define closed subsets with the finite-intersection property; hence the
inverse limit is nonempty compact.

Tightness and the weak-$*$ compact memory completion prevent positive dynamic mass
from disappearing without being represented by either a finite-energy tail or the
explicit point-at-infinity contact component.  Compatible resolvents determine a
single closed self-adjoint operator by the resolvent identity, adjoint symmetry,
injectivity and dense range.  Compatible finite-trace weights increase to a unique
lower-semicontinuous normal semifinite weight by monotone convergence, and
compatible bundle transport fixes the global marked data.  Finally, if an
additional retained quasilocal coordinate existed, the cofinal separating family
would contain a functional or defect operator nonzero on it, contradicting the
trivial-common-kernel hypothesis.  The finite classification is
Theorem~\ref{thm:supp-finite-exhaustion}; taking its indefinitely extendible compact
threads proves \eqref{eq:supp-semifinite-fibre}.
\end{proof}

\clearpage
\phantomsection
\addcontentsline{toc}{section}{References}
\begin{thebibliography}{99}

\bibitem{Connes1994}
A.~Connes,
\newblock \emph{Noncommutative Geometry},
\newblock Academic Press, San Diego, 1994.

\bibitem{RennieVarilly2008}
A.~Rennie and J.~C.~V\'arilly,
\newblock Reconstruction of manifolds in noncommutative geometry,
\newblock arXiv:math/0610418, revised 2008.

\bibitem{connesreconstruction}
A.~Connes,
\newblock On the spectral characterization of manifolds,
\newblock \emph{Journal of Noncommutative Geometry} \textbf{7} (2013), no.~1,
1--82.

\bibitem{Rieffel1999}
M.~A.~Rieffel,
\newblock Metrics on state spaces,
\newblock \emph{Documenta Mathematica} \textbf{4} (1999), 559--600.

\bibitem{Rieffel2004}
M.~A.~Rieffel,
\newblock Gromov--Hausdorff distance for quantum metric spaces,
\newblock \emph{Memoirs of the American Mathematical Society} \textbf{168}
(2004), no.~796, 1--65.
\newblock doi:10.1090/memo/0796.

\bibitem{Latremoliere2022}
F.~Latr\'emoli\`ere,
\newblock The Gromov--Hausdorff propinquity for metric spectral triples,
\newblock \emph{Advances in Mathematics} \textbf{404} (2022), 108393.
\newblock doi:10.1016/j.aim.2022.108393.

\bibitem{FarsiLatremolierePacker2024}
C.~Farsi, F.~Latr\'emoli\`ere, and J.~Packer,
\newblock Convergence of inductive sequences of spectral triples for the spectral
propinquity,
\newblock \emph{Advances in Mathematics} \textbf{437} (2024), 109442.
\newblock doi:10.1016/j.aim.2023.109442.

\bibitem{PelissierRenewalSpacetime2026}
A.~Pelissier,
\newblock Renewal Geometry and the emergence of Lorentzian spacetime,
\newblock HAL preprint hal-05744772, 2026.
\newblock \url{https://hal.science/hal-05744772}

\bibitem{Krajewski1998}
T.~Krajewski,
\newblock Classification of finite spectral triples,
\newblock \emph{Journal of Geometry and Physics} \textbf{28} (1998), 1--30.
\newblock doi:10.1016/S0393-0440(97)00068-5.

\bibitem{ChristensenIvan2006}
E.~Christensen and C.~Ivan,
\newblock Spectral triples for AF $C^*$-algebras and metrics on the Cantor set,
\newblock \emph{Journal of Operator Theory} \textbf{56} (2006), no.~1, 17--46.

\bibitem{PearsonBellissard2009}
J.~Pearson and J.~Bellissard,
\newblock Noncommutative Riemannian geometry and diffusion on ultrametric Cantor sets,
\newblock \emph{Journal of Noncommutative Geometry} \textbf{3} (2009), no.~3,
447--480.

\bibitem{ChristensenIvanSchrohe2012}
E.~Christensen, C.~Ivan, and E.~Schrohe,
\newblock Spectral triples and the geometry of fractals,
\newblock \emph{Journal of Noncommutative Geometry} \textbf{6} (2012), no.~2,
249--274.
\newblock doi:10.4171/JNCG/91.

\bibitem{KellendonkSavinien2012}
J.~Kellendonk and J.~Savinien,
\newblock Spectral triples and characterization of aperiodic order,
\newblock \emph{Proceedings of the London Mathematical Society} \textbf{104} (2012),
no.~1, 123--157.

\bibitem{Paterson2014}
A.~L.~T.~Paterson,
\newblock Contractive spectral triples for crossed products,
\newblock \emph{Mathematica Scandinavica} \textbf{114} (2014), no.~2, 275--298.
\newblock doi:10.7146/math.scand.a-17112.

\bibitem{JulienPutnam2016}
A.~Julien and I.~F.~Putnam,
\newblock Spectral triples for subshifts,
\newblock \emph{Journal of Functional Analysis} \textbf{270} (2016), no.~3,
1031--1063.
\newblock doi:10.1016/j.jfa.2015.12.002.

\bibitem{AielloGuidoIsola2022}
V.~Aiello, D.~Guido, and T.~Isola,
\newblock Spectral triples on irreversible $C^*$-dynamical systems,
\newblock \emph{International Journal of Mathematics} \textbf{33} (2022), no.~1,
2250005.
\newblock doi:10.1142/S0129167X22500057.

\bibitem{Majid2025}
S.~Majid,
\newblock Dirac operator associated to a quantum metric,
\newblock \emph{Journal of Noncommutative Geometry} \textbf{19} (2025), no.~2,
573--600.

\bibitem{CiprianiSauvageot2003}
F.~Cipriani and J.-L.~Sauvageot,
\newblock Derivations as square roots of Dirichlet forms,
\newblock \emph{Journal of Functional Analysis} \textbf{201} (2003), no.~1,
78--120.

\bibitem{HinzTeplyaev2013}
M.~Hinz and A.~Teplyaev,
\newblock Dirac and magnetic Schr\"odinger operators on fractals,
\newblock \emph{Journal of Functional Analysis} \textbf{265} (2013), no.~11,
2830--2854.
\newblock doi:10.1016/j.jfa.2013.07.021.

\bibitem{HinzKelleherTeplyaev2015}
M.~Hinz, D.~J.~Kelleher, and A.~Teplyaev,
\newblock Metric and spectral triples for Dirichlet and resistance forms,
\newblock \emph{Journal of Noncommutative Geometry} \textbf{9} (2015), no.~2,
359--390.
\newblock doi:10.4171/JNCG/195.

\bibitem{Cipriani2016}
F.~Cipriani,
\newblock Noncommutative potential theory: A survey,
\newblock \emph{Journal of Geometry and Physics} \textbf{105} (2016), 25--59.

\bibitem{Arhancet2024}
C.~Arhancet,
\newblock Spectral triples, Coulhon--Varopoulos dimension and heat kernel estimates,
\newblock \emph{Advances in Mathematics} \textbf{451} (2024), 109794.

\bibitem{ChiribellaDArianoPerinotti2009}
G.~Chiribella, G.~M.~D'Ariano, and P.~Perinotti,
\newblock Theoretical framework for quantum networks,
\newblock \emph{Physical Review A} \textbf{80} (2009), 022339.
\newblock doi:10.1103/PhysRevA.80.022339.

\bibitem{PollockProcessTensor2018}
F.~A.~Pollock, C.~Rodr\'iguez-Rosario, T.~Frauenheim, M.~Paternostro, and K.~Modi,
\newblock Non-Markovian quantum processes: Complete framework and efficient
characterization,
\newblock \emph{Physical Review A} \textbf{97} (2018), 012127.
\newblock doi:10.1103/PhysRevA.97.012127.

\bibitem{PollockEtAl2018}
F.~A.~Pollock, C.~Rodr\'iguez-Rosario, T.~Frauenheim, M.~Paternostro, and K.~Modi,
\newblock Operational Markov condition for quantum processes,
\newblock \emph{Physical Review Letters} \textbf{120} (2018), 040405.
\newblock doi:10.1103/PhysRevLett.120.040405.

\bibitem{Jaeger2000}
H.~Jaeger,
\newblock Observable operator models for discrete stochastic time series,
\newblock \emph{Neural Computation} \textbf{12} (2000), no.~6, 1371--1398.
\newblock doi:10.1162/089976600300015411.

\bibitem{LittmanSuttonSingh2001}
M.~L.~Littman, R.~S.~Sutton, and S.~Singh,
\newblock Predictive representations of state,
\newblock in \emph{Advances in Neural Information Processing Systems 14},
MIT Press, 2002, pp.~1555--1561.

\bibitem{SinghJamesRudary2004}
S.~Singh, M.~R.~James, and M.~R.~Rudary,
\newblock Predictive state representations: A new theory for modeling dynamical
systems,
\newblock in \emph{Proceedings of the 20th Conference on Uncertainty in Artificial
Intelligence}, 2004, pp.~512--519.

\bibitem{CrutchfieldYoung1989}
J.~P.~Crutchfield and K.~Young,
\newblock Inferring statistical complexity,
\newblock \emph{Physical Review Letters} \textbf{63} (1989), 105--108.
\newblock doi:10.1103/PhysRevLett.63.105.

\bibitem{shalizicrutchfield}
C.~R.~Shalizi and J.~P.~Crutchfield,
\newblock Computational mechanics: Pattern and prediction, structure and simplicity,
\newblock \emph{Journal of Statistical Physics} \textbf{104} (2001), 817--879.

\bibitem{TraversCrutchfield2025}
N.~F.~Travers and J.~P.~Crutchfield,
\newblock Equivalence of history and generator $\epsilon$-machines,
\newblock \emph{Symmetry} \textbf{17} (2025), no.~1, 78.
\newblock doi:10.3390/sym17010078.

\bibitem{MarzenCrutchfield2017}
S.~E.~Marzen and J.~P.~Crutchfield,
\newblock Structure and randomness of continuous-time, discrete-event processes,
\newblock \emph{Journal of Statistical Physics} \textbf{169} (2017), 303--315.

\bibitem{PardoGuerraThapaWashburn2026}
S.~Pardo-Guerra, A.~Thapa, and J.~Washburn,
\newblock The cactus criterion: When nonlinear Hodge theory reduces to linear on graphs,
\newblock Preprint, arXiv:2604.17775v1 (2026).

\bibitem{Rieffel1998}
M.~A.~Rieffel,
\newblock Metrics on states from actions of compact groups,
\newblock \emph{Documenta Mathematica} \textbf{3} (1998), 215--230.
\newblock doi:10.4171/DM/41.

\bibitem{SkalskiZacharias2009}
A.~Skalski and J.~Zacharias,
\newblock A note on spectral triples and quasidiagonality,
\newblock \emph{Expositiones Mathematicae} \textbf{27} (2009), no.~2, 137--141.
\newblock doi:10.1016/j.exmath.2008.10.007.

\bibitem{vanSuijlekom2021}
W.~D.~van Suijlekom,
\newblock Gromov--Hausdorff convergence of state spaces for spectral truncations,
\newblock \emph{Journal of Geometry and Physics} \textbf{162} (2021), 104075.
\newblock doi:10.1016/j.geomphys.2020.104075.

\bibitem{Leimbach2025}
M.~Leimbach,
\newblock Convergence of Peter--Weyl truncations of compact quantum groups,
\newblock \emph{Journal of Noncommutative Geometry} (2025), published online first.
\newblock doi:10.4171/JNCG/634.

\bibitem{Davies1993}
E.~B.~Davies,
\newblock Analysis on graphs and noncommutative geometry,
\newblock \emph{Journal of Functional Analysis} \textbf{111} (1993), no.~2,
398--430.

\bibitem{Requardt2002}
M.~Requardt,
\newblock Dirac operators and the calculation of the Connes metric on arbitrary
(infinite) graphs,
\newblock \emph{Journal of Physics A: Mathematical and General} \textbf{35}
(2002), no.~3, 759--779.
\newblock doi:10.1088/0305-4470/35/3/319.

\bibitem{LandryLapidusLatremoliere2021}
T.-M.~Landry, M.~L.~Lapidus, and F.~Latr\'emoli\`ere,
\newblock Metric approximations of spectral triples on the Sierpi\'nski gasket
and other fractal curves,
\newblock \emph{Advances in Mathematics} \textbf{385} (2021), 107771.
\newblock doi:10.1016/j.aim.2021.107771.

\bibitem{MirandaParra2023}
P.~Miranda and D.~Parra,
\newblock Continuum limit for a discrete Hodge--Dirac operator on square lattices,
\newblock \emph{Letters in Mathematical Physics} \textbf{113} (2023), article~45.
\newblock doi:10.1007/s11005-023-01669-9.

\bibitem{YangDill2026}
Y.-J.~Yang and K.~A.~Dill,
\newblock A principled basis for nonequilibrium network flows,
\newblock \emph{Nature Communications} \textbf{17} (2026), 4980.
\newblock doi:10.1038/s41467-026-71239-9.

\bibitem{Schnakenberg1976}
J.~Schnakenberg,
\newblock Network theory of microscopic and macroscopic behavior of master equation systems,
\newblock \emph{Reviews of Modern Physics} \textbf{48} (1976), no.~4, 571--585.
\newblock \nolinkurl{doi:10.1103/RevModPhys.48.571}.

\bibitem{Besnard2021}
F.~Besnard,
\newblock Estimating noncommutative distances on graphs,
\newblock arXiv:2105.09056 (2021).

\bibitem{BisioDArianoPerinottiChiribella2011}
A.~Bisio, G.~M.~D'Ariano, P.~Perinotti, and G.~Chiribella,
\newblock Minimal computational-space implementation of multiround quantum protocols,
\newblock \emph{Physical Review A} \textbf{83} (2011), 022325.
\newblock doi:10.1103/PhysRevA.83.022325.

\bibitem{TraversCrutchfield2014}
N.~F.~Travers and J.~P.~Crutchfield,
\newblock Equivalence of history and generator $\epsilon$-machines,
\newblock \emph{Physical Review E} \textbf{89} (2014), 062107.

\bibitem{Stinespring1955}
W.~F.~Stinespring,
\newblock Positive functions on $C^*$-algebras,
\newblock \emph{Proceedings of the American Mathematical Society} \textbf{6}
(1955), no.~2, 211--216.
\newblock doi:10.1090/S0002-9939-1955-0069403-4.

\bibitem{Holevo1982}
A.~S.~Holevo,
\newblock \emph{Probabilistic and Statistical Aspects of Quantum Theory},
\newblock North-Holland Series in Statistics and Probability, vol.~1,
North-Holland, Amsterdam, 1982.

\bibitem{HorodeckiShorRuskai2003}
M.~Horodecki, P.~W.~Shor, and M.~B.~Ruskai,
\newblock Entanglement breaking channels,
\newblock \emph{Reviews in Mathematical Physics} \textbf{15} (2003), no.~6,
629--641.

\bibitem{EkertEtAl2002}
A.~K.~Ekert, C.~M.~Alves, D.~K.~L.~Oi, M.~Horodecki, P.~Horodecki, and L.~C.~Kwek,
\newblock Direct estimations of linear and nonlinear functionals of a quantum state,
\newblock \emph{Physical Review Letters} \textbf{88} (2002), 217901.

\bibitem{GesztesyTsekanovskii2000}
F.~Gesztesy and E.~R.~Tsekanovskii,
\newblock On matrix-valued Herglotz functions,
\newblock \emph{Mathematische Nachrichten} \textbf{218} (2000), 61--138.
\newblock doi:10.1002/1522-2616(200010)218:1<61::AID-MANA61>3.0.CO;2-D.

\bibitem{KreinNudelman1977}
M.~G.~Kre\u{\i}n and A.~A.~Nudelman,
\newblock \emph{The Markov Moment Problem and Extremal Problems},
\newblock Translations of Mathematical Monographs, vol.~50,
American Mathematical Society, Providence, RI, 1977.

\bibitem{nielsenninomiya1981a}
H.~B.~Nielsen and M.~Ninomiya,
\newblock Absence of neutrinos on a lattice. I. Proof by homotopy theory,
\newblock \emph{Nuclear Physics B} \textbf{185} (1981), 20--40.

\bibitem{nielsenninomiya1981b}
H.~B.~Nielsen and M.~Ninomiya,
\newblock Absence of neutrinos on a lattice. II. Intuitive topological proof,
\newblock \emph{Nuclear Physics B} \textbf{193} (1981), 173--194.

\bibitem{Nakamura2024}
S.~Nakamura,
\newblock Remarks on discrete Dirac operators and their continuum limits,
\newblock \emph{Journal of Spectral Theory} \textbf{14} (2024), no.~1,
255--269.
\newblock doi:10.4171/JST/478.

\bibitem{lawson-michelsohn}
H.~B.~Lawson and M.-L.~Michelsohn,
\newblock \emph{Spin Geometry},
\newblock Princeton University Press, Princeton, 1989.

\end{thebibliography}

\end{document}
